% Isométries de l'espace — Mathématiques Tle C — Cours signé OnBuch+
% Programme officiel MINESEC. Compiler avec : tectonic M18-isometries-de-lespace.tex
\newcommand{\DOCMATIERE}{Mathématiques}
\newcommand{\DOCNIVEAU}{Tle C}
\newcommand{\DOCMODULE}{Géométrie dans l'espace}
\newcommand{\DOCLECON}{18}
\newcommand{\DOCTITRE}{Isométries de l'espace}
% OnBuch+ — préambule « cours signés » ENRICHI (compilé avec tectonic / XeLaTeX)
% Système visuel « Pop » d'OnBuch+ : fond crème (jamais blanc), bordures encre
% épaisses, ombres « dures » décalées, titres Archivo Black en capitales.
% Version MATHÉMATIQUES : graphes (pgfplots/tikz), boîtes pédagogiques
% (définition, propriété, méthode, attention, le savais-tu, exercice résolu,
% exercice, TP, bilan), page de garde et signature OnBuch+ ancrée en bas.
%
% Macros attendues dans main.tex AVANT \input{preamble} :
%   \DOCMATIERE  \DOCNIVEAU  \DOCMODULE  \DOCLECON (numéro)  \DOCTITRE
\documentclass[11pt]{article}

\usepackage{fontspec}
\usepackage[french]{babel}
\usepackage[a4paper,margin=2cm,top=3.4cm,bottom=2.8cm,headheight=1.4cm,footskip=1.3cm]{geometry}
\usepackage{amsmath,amssymb,amsfonts}
\usepackage{mathtools} % \xmapsto, \coloneqq... souvent utilisés par les modèles
\usepackage{cancel} % simplifications barrées
% \abs{z} (module, valeur absolue) et \norm{u} (norme), utilisés par les modèles
\providecommand{\abs}{}\let\abs\relax\DeclarePairedDelimiter{\abs}{\lvert}{\rvert}
\providecommand{\norm}{}\let\norm\relax\DeclarePairedDelimiter{\norm}{\lVert}{\rVert}
\usepackage[table]{xcolor} % [table] : fournit aussi \rowcolors (alternance de lignes)
\usepackage{pagecolor}
\usepackage{tikz}
\usetikzlibrary{arrows.meta,positioning,calc,shapes.geometric,shapes.misc,shapes.arrows,decorations.pathmorphing,decorations.markings,patterns,shadows,mindmap,trees,fit,backgrounds,matrix,intersections}
\usepackage{pgfplots}
\usepackage{pgf-pie} % diagrammes circulaires \pie (statistiques)
\pgfplotsset{compat=1.18}
\usepgfplotslibrary{fillbetween,groupplots} % aires entre courbes (calcul intégral)
\usepackage{enumitem}
\usepackage{fancyhdr}
% [most] charge toutes les bibliothèques tcolorbox (listings, minted, raster,
% documentation...) et gonfle inutilement la mémoire TeX sur un document long
% (~300 boîtes) : on ne charge que ce qui est réellement utilisé.
\usepackage[skins,breakable]{tcolorbox}
\usepackage{graphicx}
\usepackage{colortbl}
\usepackage{booktabs}
\usepackage{tabularx}
\usepackage{diagbox} % case d'en-tête diagonale des tableaux à double entrée
% Colonnes extensibles centrée / gauche / droite (C, L, R), souvent utilisées par les modèles
\newcolumntype{C}{>{\centering\arraybackslash}X}
\newcolumntype{L}{>{\raggedright\arraybackslash}X}
\newcolumntype{R}{>{\raggedleft\arraybackslash}X}
\usepackage{array}
\usepackage{multicol}
\usepackage{siunitx}
% parse-numbers=false : accepte aussi \SI{2,5 \times 10^{-3}}{...}
\sisetup{locale=FR,output-decimal-marker={,},group-minimum-digits=4,parse-numbers=false}
\DeclareSIUnit\ohm{\ensuremath{\symup{\Omega}}} % U+2126 absent de la police
\usepackage{qrcode}
% \degree : commande fréquemment utilisée par les modèles pour un angle ou une
% température en dehors de \SI{}{} (non fournie par les paquets chargés).
\providecommand{\degree}{^{\circ}}
% \qed (fin de démonstration), utilisé par les modèles sans amsthm
\providecommand{\qed}{\hfill$\square$}
% \barre : double barre des valeurs interdites dans un tableau de variations (tkz-tab)
\providecommand{\barre}{\ensuremath{\|}}
% environnement proof (amsthm non chargé) utilisé par les modèles
\ifdefined\proof\else\newenvironment{proof}[1][Démonstration]{\par\noindent\textit{#1.}\ }{\qed\par}\fi
\usepackage{lastpage}
\usepackage{hyperref}
\hypersetup{hidelinks}

% ============================================================
% Palette « Pop » OnBuch+ — mêmes hex que la classe P (plus_ui.dart)
% ============================================================
\definecolor{poporange}{HTML}{F59321}
\definecolor{popdark}{HTML}{E07A0C}
\definecolor{popcream}{HTML}{FFF6E8}
\definecolor{popink}{HTML}{1C1714}
\definecolor{poppurple}{HTML}{7A5AE0}
\definecolor{popgreen}{HTML}{2E9E6A}
\definecolor{poppink}{HTML}{E0457B}
\definecolor{popblue}{HTML}{2F7DE1}
\definecolor{popgold}{HTML}{C98A12}
\definecolor{popmuted}{HTML}{8A7F76}
\definecolor{popline}{HTML}{EADFD2}
\definecolor{popgray}{HTML}{8A7F76}
\colorlet{popgrey}{popgray}
% Alias tolérants (le modèle nomme parfois les couleurs différemment)
\colorlet{popwhite}{white}
\colorlet{popblack}{popink}
\colorlet{popred}{poppink}
\colorlet{popyellow}{popgold}
% Teintes claires pour fonds de tableaux / zones
\colorlet{popblueL}{popblue!10!white}
\colorlet{poporangeL}{poporange!14!white}
\colorlet{popgreenL}{popgreen!12!white}
\colorlet{poppurpleL}{poppurple!12!white}
\colorlet{poppinkL}{poppink!12!white}
\colorlet{popgoldL}{popgold!14!white}

\pagecolor{popcream}

% ============================================================
% Typographie
% ============================================================
\defaultfontfeatures{Ligatures=TeX}
\setmainfont{PlusJakartaSans-Medium.ttf}[Path=../../../fonts/,BoldFont=PlusJakartaSans-Bold.ttf]
% Maths en sans-serif (Fira Math) avec chiffres et lettres droites de Plus
% Jakarta, pour que les formules se fondent dans le texte.
\usepackage[math-style=ISO,bold-style=ISO]{unicode-math}
\setmathfont{FiraMath-Regular.otf}[Path=../../../fonts/]
\setmathfont{PlusJakartaSans-Medium.ttf}[Path=../../../fonts/,range={up/num,up/{Latin,latin},\mathup}]
\usepackage{icomma} % virgule décimale sans espace en mode math
% Caractères mathématiques Unicode absents des polices de texte (Plus Jakarta,
% Archivo Black) : rendus par leur équivalent mathématique, en texte comme en
% formule — sinon ils s'affichent en carré vide (ex. « Arithmétique dans ℤ »).
\usepackage{newunicodechar}
\newunicodechar{ℝ}{\ensuremath{\mathbb{R}}}
\newunicodechar{ℤ}{\ensuremath{\mathbb{Z}}}
\newunicodechar{ℂ}{\ensuremath{\mathbb{C}}}
\newunicodechar{ℕ}{\ensuremath{\mathbb{N}}}
\newunicodechar{ℚ}{\ensuremath{\mathbb{Q}}}
\newunicodechar{𝔻}{\ensuremath{\mathbb{D}}}
\newunicodechar{α}{\ensuremath{\alpha}}
\newunicodechar{β}{\ensuremath{\beta}}
\newunicodechar{γ}{\ensuremath{\gamma}}
\newunicodechar{δ}{\ensuremath{\delta}}
\newunicodechar{ε}{\ensuremath{\varepsilon}}
\newunicodechar{θ}{\ensuremath{\theta}}
\newunicodechar{λ}{\ensuremath{\lambda}}
\newunicodechar{μ}{\ensuremath{\mu}}
\newunicodechar{σ}{\ensuremath{\sigma}}
\newunicodechar{τ}{\ensuremath{\tau}}
\newunicodechar{φ}{\ensuremath{\varphi}}
\newunicodechar{ψ}{\ensuremath{\psi}}
\newunicodechar{ω}{\ensuremath{\omega}}
\newunicodechar{Σ}{\ensuremath{\Sigma}}
% majuscules produites par \MakeUppercase dans les titres (α-aminés -> Α)
\newunicodechar{Α}{\ensuremath{\alpha}}
\newunicodechar{Β}{\ensuremath{\beta}}
\newunicodechar{Γ}{\ensuremath{\Gamma}}
\newunicodechar{Δ}{\ensuremath{\Delta}}
\newunicodechar{Ε}{\ensuremath{\varepsilon}}
\newunicodechar{Θ}{\ensuremath{\Theta}}
\newunicodechar{Λ}{\ensuremath{\Lambda}}
\newunicodechar{Μ}{\ensuremath{\mu}}
\newunicodechar{Τ}{\ensuremath{\tau}}
\newunicodechar{Φ}{\ensuremath{\Phi}}
\newunicodechar{Ψ}{\ensuremath{\Psi}}
\newunicodechar{Ω}{\ensuremath{\Omega}}
\newunicodechar{↦}{\ensuremath{\mapsto}}
\newunicodechar{∈}{\ensuremath{\in}}
\newunicodechar{∉}{\ensuremath{\notin}}
\newunicodechar{∧}{\ensuremath{\wedge}}
\newunicodechar{⇔}{\ensuremath{\Leftrightarrow}}
\newunicodechar{⟺}{\ensuremath{\Longleftrightarrow}}
\newunicodechar{⇒}{\ensuremath{\Rightarrow}}
\newunicodechar{⟹}{\ensuremath{\Longrightarrow}}
\newunicodechar{∘}{\ensuremath{\circ}}
\newunicodechar{∪}{\ensuremath{\cup}}
\newunicodechar{∩}{\ensuremath{\cap}}
\newunicodechar{≡}{\ensuremath{\equiv}}
\newunicodechar{⊂}{\ensuremath{\subset}}
\newunicodechar{⊥}{\ensuremath{\perp}}
\newunicodechar{∃}{\ensuremath{\exists}}
\newunicodechar{∀}{\ensuremath{\forall}}
\newunicodechar{∛}{\ensuremath{\sqrt[3]{\,}}}
\newunicodechar{∜}{\ensuremath{\sqrt[4]{\,}}}
\newunicodechar{⃗}{\ensuremath{{}^{\to}}}
\newunicodechar{ⁿ}{\ifmmode^{n}\else\textsuperscript{\ensuremath{n}}\fi}
\newunicodechar{ˣ}{\ifmmode^{x}\else\textsuperscript{\ensuremath{x}}\fi}
\newunicodechar{ᵇ}{\ifmmode^{b}\else\textsuperscript{\ensuremath{b}}\fi}
\newunicodechar{ʳ}{\ifmmode^{r}\else\textsuperscript{\ensuremath{r}}\fi}
\newunicodechar{ᵘ}{\ifmmode^{u}\else\textsuperscript{\ensuremath{u}}\fi}
\newunicodechar{ʸ}{\ifmmode^{y}\else\textsuperscript{\ensuremath{y}}\fi}
\newunicodechar{ᵃ}{\ifmmode^{a}\else\textsuperscript{\ensuremath{a}}\fi}
\newunicodechar{ᵉ}{\ifmmode^{e}\else\textsuperscript{\ensuremath{e}}\fi}
\newunicodechar{ᵏ}{\ifmmode^{k}\else\textsuperscript{\ensuremath{k}}\fi}
\newunicodechar{ᵖ}{\ifmmode^{p}\else\textsuperscript{\ensuremath{p}}\fi}
\newunicodechar{ᴺ}{\ifmmode^{N}\else\textsuperscript{\ensuremath{N}}\fi}
\newunicodechar{ᶜ}{\ifmmode^{c}\else\textsuperscript{\ensuremath{c}}\fi}
\newunicodechar{ʰ}{\ifmmode^{h}\else\textsuperscript{\ensuremath{h}}\fi}
\newunicodechar{ᵠ}{\ifmmode^{q}\else\textsuperscript{\ensuremath{q}}\fi}
\newunicodechar{ₙ}{\ifmmode_{n}\else\textsubscript{\ensuremath{n}}\fi}
\newunicodechar{ₐ}{\ifmmode_{a}\else\textsubscript{\ensuremath{a}}\fi}
\newunicodechar{ᵢ}{\ifmmode_{i}\else\textsubscript{\ensuremath{i}}\fi}
\newunicodechar{ᵣ}{\ifmmode_{r}\else\textsubscript{\ensuremath{r}}\fi}
\newunicodechar{ₖ}{\ifmmode_{k}\else\textsubscript{\ensuremath{k}}\fi}
\newunicodechar{ᵦ}{\ifmmode_{\beta}\else\textsubscript{\ensuremath{\beta}}\fi}
\newunicodechar{ₓ}{\ifmmode_{x}\else\textsubscript{\ensuremath{x}}\fi}
\newfontfamily\popdisplay{ArchivoBlack-Regular.ttf}[Path=../../../fonts/]
\newfontfamily\poplogo{SpaceGrotesk-Bold.ttf}[Path=../../../fonts/]
\newfontfamily\popsemi{PlusJakartaSans-SemiBold.ttf}[Path=../../../fonts/]
\newfontfamily\popxbold{PlusJakartaSans-ExtraBold.ttf}[Path=../../../fonts/]
\newfontfamily\popxboldls{PlusJakartaSans-ExtraBold.ttf}[Path=../../../fonts/,LetterSpace=3.0]

\setlist[enumerate]{leftmargin=1.7em,itemsep=5pt,topsep=4pt}
\setlist[itemize]{leftmargin=1.7em,itemsep=4pt,topsep=4pt,label={\color{poporange}\textbullet}}
\setlength{\parindent}{0pt}
\setlength{\parskip}{5pt}
\renewcommand{\arraystretch}{1.25}

% pgfplots : style Pop par défaut
\pgfplotsset{
  every axis/.append style={axis line style={popink,line width=1pt},tick style={popink},
    grid style={popline,line width=0.5pt},label style={font=\small},tick label style={font=\footnotesize}},
  popcurve/.style={poporange,line width=1.8pt,no markers,smooth},
}
% Même style disponible dans un \draw TikZ simple (hors pgfplots)
\tikzset{popcurve/.style={poporange,line width=1.8pt}}

% ============================================================
% Pill / squiggle
% ============================================================
\newcommand{\poppill}[3]{%
  \tikz[baseline=-0.55ex]\node[draw=popink,line width=1.2pt,rounded corners=2.6pt,%
    fill=#2,inner xsep=6.5pt,inner ysep=3pt]{%
    \popxboldls\fontsize{7}{7}\selectfont\color{#3}\MakeUppercase{#1}};%
}
\newcommand{\popsquiggle}[1]{%
  \tikz[baseline=-0.4ex]{
    \draw[#1,line width=2.6pt,line cap=round]
      plot [smooth] coordinates {(0,0.09) (0.34,0.01) (0.68,0.21) (1.02,0.01) (1.36,0.10)};
  }%
}
% Mot-clé surligné (marqueur orange sous le texte)
\newcommand{\cle}[1]{{\popxbold\color{popdark}#1}}

% ============================================================
% En-tête / pied
% ============================================================
\pagestyle{fancy}
\fancyhf{}
\renewcommand{\headrulewidth}{0pt}
\renewcommand{\footrulewidth}{0pt}
\lhead{\raisebox{-0.15ex}{\poplogo\fontsize{13}{13}\selectfont\color{popink}OnBuch\color{poporange}+}}
\rhead{\poppill{\DOCMATIERE\ \textbullet\ \DOCNIVEAU\ \textbullet\ Leçon \DOCLECON}{white}{popink}}
\lfoot{\popxbold\fontsize{7.6}{7.6}\selectfont\color{popmuted}\MakeUppercase{Cours signé OnBuch+}\ \textbullet\ {\popsemi\DOCTITRE}}
\rfoot{\tikz[baseline=-0.6ex]\node[rounded corners=8.5pt,fill=popink,minimum height=17pt,minimum width=17pt,inner xsep=5pt,inner sep=0]{\popxbold\fontsize{8}{8}\selectfont\color{popcream}\thepage\,/\,\pageref*{LastPage}};}

% ============================================================
% Titres
% ============================================================
\newcommand{\chaptitre}[2]{%
  \begin{tcolorbox}[enhanced,colback=poporange,colframe=popink,boxrule=2.6pt,arc=11pt,
      drop shadow=popink,left=15pt,right=15pt,top=12pt,bottom=11pt,before skip=3pt,after skip=18pt]
    \poppill{#2}{popink}{popcream}\par\vspace{9pt}
    {\raggedright\hyphenpenalty=10000\exhyphenpenalty=10000\popdisplay\fontsize{22}{24}\selectfont\color{popcream}\MakeUppercase{#1}\par}
  \end{tcolorbox}
}
% Section numérotée : pastille orange avec le numéro + titre Archivo
\newcounter{popsec}
\newcommand{\coursec}[1]{%
  \stepcounter{popsec}\setcounter{popsub}{0}%
  \par\vspace{16pt}\noindent
  \begin{minipage}{\linewidth}
  \tikz[baseline=-0.7ex]\node[draw=popink,line width=1.6pt,fill=poporange,rounded corners=4pt,minimum size=22pt,inner sep=0]{\popdisplay\fontsize{12}{12}\selectfont\color{popcream}\thepopsec};%
  \hspace{8pt}{\popdisplay\fontsize{15}{17}\selectfont\color{popink}\MakeUppercase{#1}}\par
  \vspace{4pt}\noindent{\color{poporange}\rule{36pt}{3.6pt}}
  \end{minipage}\par\nopagebreak\vspace{8pt}
}
\newcounter{popsub}
\newcommand{\courssub}[1]{%
  \stepcounter{popsub}%
  \par\vspace{9pt}\noindent{\popxbold\fontsize{11.5}{13}\selectfont\color{popink}{\color{poporange}\thepopsec.\thepopsub}\hspace{6pt}#1}\par\nopagebreak\vspace{4pt}
}

% ============================================================
% Boîtes pédagogiques — même recette : fond blanc, bordure épaisse colorée,
% ombre dure encre, badge plein. Titre optionnel [#1].
% ============================================================
\tcbset{popbase/.style={breakable,enhanced,colback=white,boxrule=2.3pt,arc=9pt,
  shadow={3pt}{-3pt}{0pt}{popink},left=14pt,right=14pt,top=11pt,bottom=11pt,before skip=11pt,after skip=15pt,
  fontupper=\popsemi\fontsize{10.6}{14.5}\selectfont\color{popink},
  fontlower=\popsemi\fontsize{10.6}{14.5}\selectfont\color{popink}}}
% Coche dessinée (le glyphe ✓ n'existe pas dans les polices du document)
\newcommand{\popcheck}[1]{\tikz[baseline=-0.5ex]\draw[#1,line width=1.6pt,line cap=round,line join=round](-0.13,0.0)--(-0.04,-0.09)--(0.13,0.1);}
\newcommand{\popboxhead}[3]{\poppill{#1}{#2}{white}\ifx\relax#3\relax\else\hspace{6pt}{\popxbold\fontsize{10.6}{12}\selectfont\color{popink}#3}\fi\par\vspace{7pt}}

\newtcolorbox{aretenir}[1][]{popbase,colframe=popgreen,before upper={\popboxhead{À retenir}{popgreen}{#1}}}
\newtcolorbox{exemplebox}[1][]{popbase,colframe=poporange,before upper={\popboxhead{Exemple}{poporange}{#1}}}
\newtcolorbox{definition}[1][]{popbase,colframe=popblue,before upper={\popboxhead{Définition}{popblue}{#1}}}
\newtcolorbox{propriete}[1][]{popbase,colframe=poppurple,before upper={\popboxhead{Propriété}{poppurple}{#1}}}
\newtcolorbox{methode}[1][]{popbase,colframe=popgold,before upper={\popboxhead{Méthode}{popgold}{#1}}}
\newtcolorbox{attention}[1][]{popbase,colframe=poppink,before upper={\popboxhead{Attention}{poppink}{#1}}}
\newtcolorbox{experience}[1][]{popbase,colframe=popblue,colback=popblueL,before upper={\popboxhead{Expérience}{popblue}{#1}}}
% « Le savais-tu ? » : carte violette pleine (contexte, culture, Cameroun)
\newtcolorbox{savaistu}[1][]{popbase,colback=poppurple,colframe=popink,
  fontupper=\popsemi\fontsize{10.4}{14.2}\selectfont\color{white},
  before upper={\poppill{Le savais-tu\,?}{white}{poppurple}\ifx\relax#1\relax\else\hspace{6pt}{\popxbold\color{white}#1}\fi\par\vspace{7pt}}}
% Exercice résolu : énoncé en haut, \tcblower puis la solution sur fond crème
\newcounter{popexo}\newcounter{popexr}
\newtcolorbox{exoresolu}[1][]{popbase,colframe=poporange,colbacklower=popcream,
  segmentation style={popink,line width=1.2pt,dashed},
  before upper={\stepcounter{popexr}\popboxhead{Exercice résolu \thepopexr}{poporange}{#1}},
  before lower={\poppill{Solution}{popink}{popcream}\par\vspace{6pt}}}
% Exercice (non résolu en place) — la correction est dans la partie Corrigés
\newtcolorbox{exercice}[1][]{popbase,colframe=popink,
  before upper={\stepcounter{popexo}\popboxhead{Exercice \thepopexo}{popink}{#1}}}
% Corrigé
\newtcolorbox{corrige}[1][]{popbase,colframe=popgreen,colback=popgreenL,
  before upper={\popboxhead{Corrigé}{popgreen}{#1}}}
% Objectifs / prérequis / bilan
\newtcolorbox{objectifs}{popbase,colframe=popink,colback=poporangeL,
  before upper={\popboxhead{Objectifs de la leçon}{popink}{}}}
\newtcolorbox{prerequis}{popbase,colframe=popink,colback=popblueL,
  before upper={\popboxhead{Prérequis}{popblue}{}}}
\newtcolorbox{bilan}{popbase,colframe=popink,colback=white,boxrule=2.8pt,
  before upper={\popboxhead{Fiche bilan}{poporange}{L'essentiel à connaître pour l'examen}}}
% Figure encadrée avec légende
\newtcolorbox{popfigure}[1][]{enhanced,breakable=false,colback=white,colframe=popline,boxrule=1.4pt,arc=8pt,
  left=10pt,right=10pt,top=10pt,bottom=8pt,before skip=10pt,after skip=12pt,halign=center,
  lower separated=false,halign lower=center,
  fontlower=\popsemi\fontsize{9}{11.5}\selectfont\color{popmuted},
  #1}
\newcommand{\legende}[1]{\par\vspace{6pt}{\popsemi\fontsize{9}{11.5}\selectfont\color{popmuted}#1}}

% ============================================================
% Signature OnBuch+
% ============================================================
\newcommand{\poplogoinline}[1]{{\poplogo\fontsize{#1}{#1}\selectfont\color{popink}OnBuch\color{poporange}+}}
\newcommand{\signatureonbuch}{%
  \begin{tcolorbox}[enhanced,colback=popink,colframe=popink,boxrule=2.2pt,arc=10pt,
      drop shadow=poporange,left=14pt,right=14pt,top=10pt,bottom=10pt,before skip=0pt,after skip=0pt]
    \tikz[baseline=-0.7ex]\node[circle,fill=popgreen,draw=popcream,line width=1.3pt,minimum size=22pt,inner sep=0]{\popcheck{white}};%
    \hspace{9pt}\begin{minipage}[c]{0.62\linewidth}
      {\popxbold\fontsize{12}{13}\selectfont\color{popcream}Signé\ \textbullet\ {\poplogo OnBuch\color{poporange}+}}\par\vspace{2pt}
      {\raggedright\popsemi\fontsize{8}{10.5}\selectfont\color{popline}\MakeUppercase{Équipe pédagogique OnBuch+}\\\MakeUppercase{Conforme au programme officiel MINESEC}\par}
    \end{minipage}\hfill
    {\poplogo\fontsize{22}{22}\selectfont\color{popcream}OnBuch\color{poporange}+}
  \end{tcolorbox}%
}

% ============================================================
% Page de garde
% #1 titre · #2 matière · #3 niveau · #4 module · #5 n° leçon · #6 durée ·
% #7 description / objectifs (texte ou itemize)
% ============================================================
\newcommand{\pagegarde}[7]{%
  \thispagestyle{empty}
  \newgeometry{a4paper,margin=1.7cm,top=1.6cm,bottom=1.4cm}%
  \begin{tikzpicture}[remember picture,overlay]
    \fill[poporange!18] ([xshift=-3.2cm,yshift=-3.6cm]current page.north east) circle (4.6cm);
    \fill[poppurple!14] ([xshift=2.4cm,yshift=7.6cm]current page.south west) circle (3.8cm);
  \end{tikzpicture}%
  {\poplogo\fontsize{30}{30}\selectfont\color{popink}OnBuch\color{poporange}+}\hfill
  \raisebox{0.6ex}{\poppill{Cours signé \textbullet\ Premium}{popink}{popcream}}\par
  \vspace{1pt}\popsquiggle{poppurple}\par
  \vspace{8pt}
  \begin{tcolorbox}[enhanced,colback=poporange,colframe=popink,boxrule=2.8pt,arc=13pt,
      drop shadow=popink,left=18pt,right=18pt,top=12pt,bottom=13pt,after skip=10pt]
    \poppill{#2 \textbullet\ #3}{popink}{popcream}\hspace{5pt}\poppill{#4}{white}{popink}\par\vspace{8pt}
    {\popdisplay\fontsize{14}{14}\selectfont\color{popink}LEÇON #5}\par\vspace{4pt}
    {\raggedright\hyphenpenalty=10000\exhyphenpenalty=10000\popdisplay\fontsize{25}{27}\selectfont\color{popcream}\MakeUppercase{#1}\par}
    \vspace{8pt}
    {\popxbold\fontsize{9.5}{9.5}\selectfont\color{popink}\MakeUppercase{Durée conseillée : #6}}
  \end{tcolorbox}
  \begin{tcolorbox}[enhanced,colback=white,colframe=popink,boxrule=2.2pt,arc=10pt,
      drop shadow=popink,left=16pt,right=16pt,top=10pt,bottom=10pt,after skip=8pt]
    \poppill{Dans cette leçon}{poppurple}{white}\par\vspace{5pt}
    \popsemi\fontsize{9.3}{12.6}\selectfont\color{popink}#7
  \end{tcolorbox}
  \noindent\begin{minipage}[t]{0.487\linewidth}
    \begin{tcolorbox}[enhanced,colback=popgreen,colframe=popink,boxrule=2.2pt,arc=10pt,
        drop shadow=popink,left=11pt,right=11pt,top=10pt,bottom=10pt,height=3.1cm,valign=center]
      \tcbox[colback=white,colframe=popink,boxrule=1.3pt,arc=5pt,boxsep=0pt,nobeforeafter,
        left=3pt,right=3pt,top=3pt,bottom=3pt,valign=center]{\qrcode[height=1.9cm]{https://play.google.com/store/apps/details?id=com.luvvix.onbuch&hl=fr}}%
      \hspace{8pt}\begin{minipage}[c]{0.5\linewidth}
        {\popdisplay\fontsize{11}{12.5}\selectfont\color{white}\MakeUppercase{Télécharge OnBuch}}\par\vspace{4pt}
        {\raggedright\popsemi\fontsize{8.4}{11}\selectfont\color{white}Gratuit \textbullet\ Google Play\\Tuteur IA, annales, concours, résultats\par}
      \end{minipage}
    \end{tcolorbox}
  \end{minipage}\hfill
  \begin{minipage}[t]{0.487\linewidth}
    \begin{tcolorbox}[enhanced,colback=poppurple,colframe=popink,boxrule=2.2pt,arc=10pt,
        drop shadow=popink,left=11pt,right=11pt,top=10pt,bottom=10pt,height=3.1cm,valign=center]
      \tikz[baseline=-0.7ex]\node[circle,fill=white,draw=popink,line width=1.3pt,minimum size=1.6cm,inner sep=0]{\popdisplay\fontsize{24}{24}\selectfont\color{poppurple}+};%
      \hspace{8pt}\begin{minipage}[c]{0.6\linewidth}
        {\popdisplay\fontsize{11}{12.5}\selectfont\color{white}\MakeUppercase{Passe à OnBuch+}}\par\vspace{4pt}
        {\raggedright\popsemi\fontsize{8.4}{11}\selectfont\color{white}Tous les cours signés, Tuteur IA illimité, fiches PDF — onglet OnBuch+ de l'app\par}
      \end{minipage}
    \end{tcolorbox}
  \end{minipage}
  \vspace{8pt}
  \popcontenu
  \vfill
  \signatureonbuch
  \restoregeometry
  \newpage
}

% Rangée « Ce cours contient » (page de garde)
\newcommand{\poptile}[3]{% #1 couleur #2 gros texte #3 légende
  \begin{tcolorbox}[enhanced,colback=#1,colframe=popink,boxrule=2pt,arc=9pt,shadow={2.5pt}{-2.5pt}{0pt}{popink},
      left=6pt,right=6pt,top=9pt,bottom=9pt,halign=center,valign=center,height=1.75cm,width=\linewidth,nobeforeafter]
    {\popdisplay\fontsize{12.5}{13}\selectfont\color{white}\MakeUppercase{#2}}\par\vspace{3pt}
    {\centering\popsemi\fontsize{7.8}{9.5}\selectfont\color{white}#3\par}
  \end{tcolorbox}}
\newcommand{\popcontenu}{%
  \par\noindent{\popxboldls\fontsize{7.5}{7.5}\selectfont\color{popmuted}CE COURS SIGNÉ CONTIENT}\par\vspace{7pt}
  \noindent\begin{minipage}[t]{0.235\linewidth}\poptile{popblue}{Cours}{complet, illustré, pas à pas}\end{minipage}\hfill
  \begin{minipage}[t]{0.235\linewidth}\poptile{poporange}{Méthodes}{et exercices résolus}\end{minipage}\hfill
  \begin{minipage}[t]{0.235\linewidth}\poptile{poppink}{Exercices}{type examen + corrigés}\end{minipage}\hfill
  \begin{minipage}[t]{0.235\linewidth}\poptile{popgreen}{Fiche bilan}{l'essentiel à retenir}\end{minipage}\par}

% Sommaire
\newtcolorbox{sommaire}{enhanced,colback=white,colframe=popink,boxrule=2.2pt,arc=10pt,
  drop shadow=popink,left=18pt,right=18pt,top=13pt,bottom=13pt,before skip=6pt,after skip=20pt,
  before upper={\poppill{Sommaire}{poppurple}{white}\par\vspace{9pt}\popsemi\fontsize{10.6}{14.5}\selectfont\color{popink}}}

% Signature de fin de cours — poussée tout en bas de la dernière page
\newcommand{\findecours}{%
  \par\vfill\nopagebreak
  \begin{center}\popsquiggle{poporange}\end{center}\vspace{4pt}
  \signatureonbuch
}
\AtBeginDocument{\def\llbracket{\mathopen{[\mkern-3.5mu[}}\def\rrbracket{\mathclose{]\mkern-3.5mu]}}} % intervalles d'entiers (glyphe ⟦ absent de la police)
% Glyphes absents de Fira Math : redéfinis à partir de glyphes disponibles
\AtBeginDocument{%
  \def\setminus{\mathbin{\backslash}}\let\smallsetminus\setminus
  \def\vdots{\mathord{\vbox{\baselineskip3.2pt\lineskiplimit0pt\kern4pt\hbox{.}\hbox{.}\hbox{.}}}}%
  \def\ddots{\mathinner{\mkern1mu\raise6.5pt\hbox{.}\mkern2mu\raise3.6pt\hbox{.}\mkern2mu\raise0.7pt\hbox{.}\mkern1mu}}%
  \def\triangle{\mathord{\scalebox{0.9}{$\symup{\Delta}$}}}%
  \def\nabla{\mathord{\raisebox{\depth}{\scalebox{1}[-1]{$\symup{\Delta}$}}}}%
  \def\checkmark{\popcheck{popdark}}%
  \def\star{\mathbin{\ast}}\def\bigstar{\mathord{\ast}}%
  \def\diamond{\mathbin{\scalebox{0.6}{$\Diamond$}}}%
}

%% FIGFIX-BEGIN v3 — correctifs globaux des figures (docs/onbuch-generation/tools/figfix)
\usepackage{adjustbox}\usepackage{etoolbox}
% toute figure TikZ/pgfplots plus large que la ligne est réduite à la largeur disponible
\BeforeBeginEnvironment{tikzpicture}{\begin{adjustbox}{max width=\linewidth}}
\AfterEndEnvironment{tikzpicture}{\end{adjustbox}}
% légendes pgfplots lisibles par-dessus les courbes
\pgfplotsset{every axis/.append style={legend style={fill=white,fill opacity=0.92,text opacity=1,draw=black!15,inner sep=3pt}}}
% étiquettes d'arêtes (syntaxe edge["..."]) avec fond blanc
\tikzset{every edge quotes/.append style={fill=white,inner sep=1.2pt}}
% bulles de cartes mentales : texte à la ligne dans la bulle, bulle un peu plus grande
\tikzset{every concept/.append style={text width=2.0cm,align=center,minimum size=2.3cm,inner sep=2pt}}
%% FIGFIX-END
\begin{document}

\pagegarde{Isométries de l'espace}{Mathématiques}{Tle C}{Géométrie dans l'espace}{18}{12 h}{Cette leçon introduit les isométries fondamentales de l'espace : la réflexion par rapport à un plan et le demi-tour autour d'une droite. On y étudie leurs expressions analytiques et la nature des composées de deux réflexions ou de deux demi-tours, avec des applications à la résolution de problèmes de géométrie dans l'espace.
\vspace{4pt}
\begin{multicols}{2}\small
\begin{itemize}
\item À la fin de cette leçon, je sais définir une réflexion de l'espace et caractériser son plan de symétrie
\item À la fin de cette leçon, je sais déterminer l'expression analytique d'une réflexion donnée par son plan
\item À la fin de cette leçon, je sais définir un demi-tour de l'espace et caractériser son axe
\item À la fin de cette leçon, je sais déterminer l'expression analytique d'un demi-tour donné par son axe
\item À la fin de cette leçon, je sais reconnaître la nature d'une isométrie donnée par son expression analytique
\item À la fin de cette leçon, je sais déterminer la nature et les éléments caractéristiques de la composée de deux réflexions de plans parallèles ou perpendiculaires
\end{itemize}\end{multicols}}

\begin{sommaire}
\begin{multicols}{2}
\begin{enumerate}
\item Réflexion de l'espace : définition et propriétés
\item Expression analytique d'une réflexion
\item Composée de deux réflexions
\item Demi-tour : définition et expression analytique
\item Composée de deux demi-tours
\item Reconnaissance d'isométries et résolution de problèmes
\item Activité d'intégration
\item Méthodes et exercices résolus
\item Exercices
\item Corrigés des exercices
\item Fiche bilan
\end{enumerate}
\end{multicols}
\end{sommaire}

\begin{objectifs}
\begin{itemize}
\item À la fin de cette leçon, je sais définir une réflexion de l'espace et caractériser son plan de symétrie
\item À la fin de cette leçon, je sais déterminer l'expression analytique d'une réflexion donnée par son plan
\item À la fin de cette leçon, je sais définir un demi-tour de l'espace et caractériser son axe
\item À la fin de cette leçon, je sais déterminer l'expression analytique d'un demi-tour donné par son axe
\item À la fin de cette leçon, je sais reconnaître la nature d'une isométrie donnée par son expression analytique
\item À la fin de cette leçon, je sais déterminer la nature et les éléments caractéristiques de la composée de deux réflexions de plans parallèles ou perpendiculaires
\item À la fin de cette leçon, je sais déterminer la nature et les éléments caractéristiques de la composée de deux demi-tours d'axes parallèles ou orthogonaux
\item À la fin de cette leçon, je sais utiliser les isométries de l'espace pour résoudre un problème de géométrie (alignement, distances, sections planes)
\end{itemize}
\end{objectifs}

\begin{prerequis}
\begin{itemize}
\item Symétries dans le plan (symétrie axiale, symétrie centrale) vues en Première
\item Produit scalaire dans l'espace et équation cartésienne d'un plan
\item Projeté orthogonal d'un point sur un plan ou une droite, distance d'un point à un plan
\item Repérage dans l'espace : repère orthonormé, coordonnées, vecteur normal à un plan
\item Notion de composée de deux applications
\end{itemize}
\end{prerequis}

\newpage

\coursec{Réflexion de l'espace : définition et propriétés}

\textbf{Situation d'accroche.} À Douala, un technicien installe un grand miroir mural dans le salon d'une villa de Bonanjo. En regardant le reflet du lustre, il a l'impression que celui-ci se trouve \emph{derrière} le mur, à égale distance du miroir. À Yaoundé, dans un atelier de menuiserie du quartier Mvog-Ada, on retourne une armoire autour de son arête verticale pour l'encastrer dans un angle : le meuble pivote d'un demi-tour. Comment décrire mathématiquement ces transformations de l'espace ? Comment prévoir le résultat de deux symétries successives, par exemple l'image d'un objet réfléchi par les deux miroirs parallèles d'une cabine d'essayage ?

\textbf{Problématique.} Quelles sont les transformations de l'espace qui conservent les distances ? Comment les définir rigoureusement, les construire, et que conservent-elles exactement ? Dans cette première section, nous étudions la plus fondamentale d'entre elles : la \cle{réflexion}, c'est-à-dire la symétrie orthogonale par rapport à un plan.

\courssub{Projeté orthogonal sur un plan et distance d'un point à un plan}

Avant de parler de symétrie, il faut savoir « descendre perpendiculairement » d'un point vers un plan, exactement comme un fil à plomb descend verticalement du lustre vers le sol.

\begin{definition}[Projeté orthogonal d'un point sur un plan]
Soit $\mathcal{P}$ un plan de l'espace et $M$ un point. Le \cle{projeté orthogonal} de $M$ sur $\mathcal{P}$ est le point $H$ intersection de $\mathcal{P}$ et de la droite passant par $M$ et perpendiculaire à $\mathcal{P}$.
\begin{itemize}
\item Si $M \in \mathcal{P}$, alors $H = M$.
\item Si $M \notin \mathcal{P}$, la droite $(MH)$ est perpendiculaire à $\mathcal{P}$ : elle est perpendiculaire à \emph{toute} droite de $\mathcal{P}$ passant par $H$.
\end{itemize}
\end{definition}

\begin{propriete}[Distance d'un point à un plan]
Soit $\mathcal{P}$ un plan, $M$ un point et $H$ son projeté orthogonal sur $\mathcal{P}$. Pour tout point $K$ de $\mathcal{P}$ distinct de $H$, on a $MK > MH$. La \cle{distance} de $M$ au plan $\mathcal{P}$, notée $d(M, \mathcal{P})$, est donc :
\[ d(M, \mathcal{P}) = MH. \]
\end{propriete}

\begin{experience}[Justification par le théorème de Pythagore]
Pourquoi $H$ est-il le point de $\mathcal{P}$ le plus proche de $M$ ? Raisonnons pas à pas.
\begin{enumerate}
\item Soit $K$ un point quelconque de $\mathcal{P}$, distinct de $H$.
\item La droite $(MH)$ est perpendiculaire au plan $\mathcal{P}$, donc perpendiculaire à la droite $(HK)$ qui est contenue dans $\mathcal{P}$. Le triangle $MHK$ est donc \textbf{rectangle en $H$}.
\item Le théorème de Pythagore donne : $MK^2 = MH^2 + HK^2$.
\item Comme $HK^2 > 0$ (car $K \neq H$), on obtient $MK^2 > MH^2$, donc $MK > MH$.
\end{enumerate}
Le projeté orthogonal réalise bien le plus court chemin de $M$ vers $\mathcal{P}$. C'est ce même théorème de Pythagore qui justifiera plus loin la conservation des distances par réflexion.
\end{experience}

\begin{exemplebox}[Distance d'un point à un plan]
Dans un repère orthonormé $(O; \vec{i}, \vec{j}, \vec{k})$, soit $\mathcal{P}$ le plan d'équation $z = 3$ (plan horizontal à l'altitude 3) et $M(2; -1; 7)$. La perpendiculaire à $\mathcal{P}$ passant par $M$ est dirigée par $\vec{k}$ ; elle coupe $\mathcal{P}$ en $H(2; -1; 3)$. Alors :
\[ d(M, \mathcal{P}) = MH = |7 - 3| = 4. \]
\end{exemplebox}

\courssub{Définition de la réflexion de plan $\mathcal{P}$}

Regarde ton visage dans un miroir mural : ton image semble située derrière le mur, à la même distance du miroir que toi, et le segment joignant un point de ton visage à son image est perpendiculaire au miroir. C'est exactement cela, une réflexion.

\begin{definition}[Réflexion de plan $\mathcal{P}$]
Soit $\mathcal{P}$ un plan de l'espace. La \cle{réflexion} de plan $\mathcal{P}$ (ou \cle{symétrie orthogonale} par rapport à $\mathcal{P}$), notée $s_{\mathcal{P}}$, est la transformation de l'espace qui à tout point $M$ associe le point $M'$ défini par :
\begin{itemize}
\item si $M \in \mathcal{P}$, alors $M' = M$ ;
\item si $M \notin \mathcal{P}$, alors $\mathcal{P}$ est le \cle{plan médiateur} du segment $[MM']$, c'est-à-dire le plan perpendiculaire à $[MM']$ en son milieu.
\end{itemize}
On écrit $M' = s_{\mathcal{P}}(M)$ et on dit que $M'$ est l'\emph{image} de $M$ par $s_{\mathcal{P}}$.
\end{definition}

\begin{propriete}[Caractérisation vectorielle]
Soit $\mathcal{P}$ un plan, $M$ un point, $H$ son projeté orthogonal sur $\mathcal{P}$ et $M' = s_{\mathcal{P}}(M)$. Alors :
\[ \overrightarrow{HM'} = -\overrightarrow{HM} \qquad \text{ou, de façon équivalente,} \qquad \overrightarrow{MM'} = 2\,\overrightarrow{MH}. \]
Autrement dit, $H$ est le \textbf{milieu} du segment $[MM']$.
\end{propriete}

Cette caractérisation est fondamentale : c'est elle qui permettra, dans la section suivante, d'obtenir les expressions analytiques des réflexions.

\begin{popfigure}
\begin{center}
\begin{tikzpicture}[scale=1.1, line cap=round]
\% Plan P vu de profil (parallélogramme)
\fill[popblueL, opacity=0.55] (-3.2,0) -- (-1.6,-1.1) -- (3.4,-1.1) -- (1.8,0) -- cycle;
\draw[popblue, thick] (-3.2,0) -- (-1.6,-1.1) -- (3.4,-1.1) -- (1.8,0) -- cycle;
\node[popblue] at (2.7,-0.75) {$\mathcal{P}$};
\% Points
\coordinate (M) at (-0.6,2.4);
\coordinate (H) at (-0.6,-0.55);
\coordinate (Mp) at (-0.6,-3.5);
\% Segment MM' perpendiculaire
\draw[popdark, dashed, thick] (M) -- (Mp);
\draw[popink, very thick] (M) -- (H);
\draw[poporange, very thick] (H) -- (Mp);
\% Angle droit
\draw[popdark] (-0.6,-0.25) -- (-0.3,-0.25) -- (-0.3,-0.55);
\% Points et labels
\fill[popink] (M) circle (2.2pt) node[above left] {$M$};
\fill[popdark] (H) circle (2.2pt) node[right=3pt] {$H$};
\fill[poporange] (Mp) circle (2.2pt) node[below left] {$M' = s_{\mathcal{P}}(M)$};
\% Distances égales
\node[popink, left] at (-0.75,0.9) {$MH$};
\node[poporange, left] at (-0.75,-2.1) {$HM'$};
\node[popmuted, align=center] at (2.9,1.6) {$MH = HM'$\\ $(MM') \perp \mathcal{P}$};
\end{tikzpicture}
\end{center}
\legende{La réflexion de plan $\mathcal{P}$ : $H$ est le projeté orthogonal de $M$ sur $\mathcal{P}$ et $H$ est le milieu de $[MM']$. Le segment $[MM']$ est perpendiculaire à $\mathcal{P}$.}
\end{popfigure}

\begin{methode}[Construire l'image d'un point par réflexion]
Pour construire (ou calculer) l'image $M'$ d'un point $M$ par la réflexion de plan $\mathcal{P}$ :
\begin{enumerate}
\item \textbf{Tester} si $M \in \mathcal{P}$. Si oui, $M' = M$ et c'est terminé.
\item \textbf{Construire le projeté orthogonal} $H$ de $M$ sur $\mathcal{P}$ : mener par $M$ la perpendiculaire à $\mathcal{P}$ et noter $H$ son intersection avec $\mathcal{P}$.
\item \textbf{Reporter} : placer $M'$ sur la droite $(MH)$, de l'autre côté de $\mathcal{P}$, tel que $HM' = MH$, c'est-à-dire $\overrightarrow{HM'} = -\overrightarrow{HM}$.
\end{enumerate}
En calcul, l'étape 2 revient à résoudre un système (paramétrage de la perpendiculaire + équation du plan), et l'étape 3 utilise la relation vectorielle ci-dessus.
\end{methode}

\begin{exoresolu}[Image d'un point par réflexion]
L'espace est muni d'un repère orthonormé $(O; \vec{i}, \vec{j}, \vec{k})$. Soit $\mathcal{P}$ le plan d'équation $z = 0$ (le plan $(O; \vec{i}, \vec{j})$) et $A(1; -2; 5)$.
Déterminer le projeté orthogonal $H$ de $A$ sur $\mathcal{P}$, puis l'image $A'$ de $A$ par $s_{\mathcal{P}}$.
\tcblower
\textbf{Solution détaillée.}
\textbf{Étape 1.} $A \notin \mathcal{P}$ car sa cote $5$ est non nulle.
\textbf{Étape 2.} Le plan $\mathcal{P}$ est horizontal ; la perpendiculaire à $\mathcal{P}$ passant par $A$ est la droite verticale passant par $A$, dirigée par $\vec{k}$. Elle coupe $\mathcal{P}$ en $H(1; -2; 0)$.
\textbf{Étape 3.} On utilise $\overrightarrow{HA'} = -\overrightarrow{HA}$. Or $\overrightarrow{HA}(0; 0; 5)$, donc $\overrightarrow{HA'}(0; 0; -5)$, d'où :
\[ A' = H + \overrightarrow{HA'} = (1; -2; -5). \]
\textbf{Vérification.} $H$ est bien le milieu de $[AA']$ : $\left(\dfrac{1+1}{2}; \dfrac{-2-2}{2}; \dfrac{5+(-5)}{2}\right) = (1; -2; 0) = H$, et $(AA')$, verticale, est perpendiculaire à $\mathcal{P}$. La réflexion par rapport au plan $z = 0$ transforme donc $(x; y; z)$ en $(x; y; -z)$ : seule la cote change de signe.
\end{exoresolu}

\courssub{Premières propriétés : involution et points fixes}

\begin{propriete}[La réflexion est une involution]
Soit $\mathcal{P}$ un plan. La réflexion $s_{\mathcal{P}}$ est une \cle{involution}, c'est-à-dire :
\[ s_{\mathcal{P}} \circ s_{\mathcal{P}} = \mathrm{Id}_{\mathcal{E}}, \]
où $\mathrm{Id}_{\mathcal{E}}$ désigne l'application identique de l'espace. En particulier, $s_{\mathcal{P}}$ est une bijection de l'espace, et sa bijection réciproque est elle-même : $s_{\mathcal{P}}^{-1} = s_{\mathcal{P}}$.
\end{propriete}

\begin{experience}[Démonstration guidée de l'involutivité]
Montrons que pour tout point $M$, $s_{\mathcal{P}}(s_{\mathcal{P}}(M)) = M$.
\begin{enumerate}
\item \textbf{Cas 1 : $M \in \mathcal{P}$.} Alors $s_{\mathcal{P}}(M) = M$, donc $s_{\mathcal{P}}(s_{\mathcal{P}}(M)) = s_{\mathcal{P}}(M) = M$.
\item \textbf{Cas 2 : $M \notin \mathcal{P}$.} Posons $M' = s_{\mathcal{P}}(M)$. Par définition, $\mathcal{P}$ est le plan médiateur de $[MM']$.
\item Or la relation « $\mathcal{P}$ est le plan médiateur de $[MM']$ » est \emph{symétrique} en $M$ et $M'$ : le plan médiateur de $[M'M]$ est le même que celui de $[MM']$. Donc $\mathcal{P}$ est aussi le plan médiateur de $[M'M]$.
\item Par définition de $s_{\mathcal{P}}$, cela signifie exactement que $s_{\mathcal{P}}(M') = M$, c'est-à-dire $s_{\mathcal{P}}(s_{\mathcal{P}}(M)) = M$.
\end{enumerate}
Dans les deux cas, $s_{\mathcal{P}} \circ s_{\mathcal{P}}(M) = M$ : la réflexion est involutive. \emph{Image de l'image, on revient au point de départ — comme quand on regarde le reflet du reflet dans le même miroir.}
\end{experience}

\begin{propriete}[Points fixes]
L'ensemble des points \cle{invariants} (fixes) par la réflexion $s_{\mathcal{P}}$ est \textbf{exactement} le plan $\mathcal{P}$ :
\[ s_{\mathcal{P}}(M) = M \iff M \in \mathcal{P}. \]
\end{propriete}

En effet, si $M \notin \mathcal{P}$, alors $H \neq M$ et $M'$ est le symétrique de $M$ par rapport à $H$ sur la perpendiculaire, donc $M' \neq M$. Réciproquement, tout point de $\mathcal{P}$ est fixe par définition. Cette propriété est précieuse pour \emph{reconnaître} une réflexion : si une isométrie laisse fixes tous les points d'un plan, et seulement ceux-là, c'est la réflexion de ce plan (ou l'identité si tout l'espace est fixe).

\courssub{La réflexion est une isométrie : conservation des distances}

Voici la propriété qui justifie le titre de la leçon.

\begin{propriete}[Conservation des distances — admise]
La réflexion $s_{\mathcal{P}}$ \cle{conserve les distances} : pour tous points $A$ et $B$ d'images $A'$ et $B'$,
\[ A'B' = AB. \]
On dit que $s_{\mathcal{P}}$ est une \cle{isométrie} de l'espace. Cette propriété est admise ; elle se justifie par le théorème de Pythagore dans l'espace, en décomposant le segment $[AB]$ selon une composante parallèle à $\mathcal{P}$ (inchangée) et une composante perpendiculaire (dont seul le sens change, pas la longueur).
\end{propriete}

\begin{propriete}[Conséquences de la conservation des distances]
La réflexion $s_{\mathcal{P}}$ conserve :
\begin{itemize}
\item les \cle{milieux} : si $I$ est le milieu de $[AB]$, alors $I' = s_{\mathcal{P}}(I)$ est le milieu de $[A'B']$ ;
\item l'\cle{alignement} : l'image de trois points alignés est formée de trois points alignés ;
\item le \cle{produit scalaire} : pour tous vecteurs $\vec{u}$ et $\vec{v}$ d'images $\vec{u}'$ et $\vec{v}'$, on a $\vec{u}' \cdot \vec{v}' = \vec{u} \cdot \vec{v}$ ; en particulier la réflexion conserve les \cle{angles géométriques} et l'\cle{orthogonalité} ;
\item les \cle{aires} et les \cle{volumes}.
\end{itemize}
\end{propriete}

La conservation du produit scalaire découle de celle des distances via l'identité de polarisation :
\[ \vec{u} \cdot \vec{v} = \frac{1}{2}\left(\|\vec{u} + \vec{v}\|^2 - \|\vec{u}\|^2 - \|\vec{v}\|^2\right), \]
où le second membre ne fait intervenir que des normes, donc des distances, qui sont conservées. La conservation des milieux vient de la caractérisation vectorielle : $I$ milieu de $[AB]$ signifie $\overrightarrow{AI} = \frac{1}{2}\overrightarrow{AB}$, relation préservée car une réflexion agit linéairement sur les vecteurs.

\begin{attention}[Réflexion ou symétrie centrale : ne pas confondre !]
Dans l'espace, il existe \textbf{trois} symétries orthogonales distinctes, et les confondre est l'erreur la plus fréquente au Baccalauréat :
\begin{itemize}
\item la \textbf{réflexion} : symétrie par rapport à un \emph{plan} $\mathcal{P}$ — les points fixes forment un plan ;
\item le \textbf{demi-tour} : symétrie orthogonale par rapport à une \emph{droite} $(d)$ — les points fixes forment une droite (étudié en section 4) ;
\item la \textbf{symétrie centrale} : symétrie par rapport à un \emph{point} $\Omega$ — un seul point fixe.
\end{itemize}
La symétrie centrale de centre $\Omega$ (définie par $\overrightarrow{\Omega M'} = -\overrightarrow{\Omega M}$) n'est \textbf{pas} une réflexion : elle transforme $(x;y;z)$ en $(2a-x;\, 2b-y;\, 2c-z)$ et renverse les trois coordonnées, alors qu'une réflexion n'en renverse qu'une seule (dans un repère bien choisi). Dans le plan, la symétrie centrale coïncide avec un demi-tour ; dans l'espace, ce n'est plus le cas. Retiens : \emph{réflexion = miroir = plan}.
\end{attention}

\courssub{Effet sur les figures usuelles et renversement de l'orientation}

\begin{propriete}[Images des figures usuelles par une réflexion]
Soit $s_{\mathcal{P}}$ une réflexion.
\begin{itemize}
\item L'image d'une \textbf{droite} $(d)$ est une droite $(d')$ ; si $(d)$ est parallèle à $\mathcal{P}$, $(d')$ lui est parallèle ; si $(d)$ coupe $\mathcal{P}$ en un point $A$, $(d')$ passe aussi par $A$ et $(d)$ et $(d')$ sont symétriques par rapport à $\mathcal{P}$.
\item L'image d'un \textbf{plan} $\mathcal{Q}$ est un plan $\mathcal{Q}'$ ; si $\mathcal{Q} \parallel \mathcal{P}$, alors $\mathcal{Q}' \parallel \mathcal{P}$, à égale distance de l'autre côté.
\item L'image d'une \textbf{sphère} $\mathcal{S}(\Omega, R)$ est la sphère $\mathcal{S}(\Omega', R)$ de \textbf{même rayon}, où $\Omega' = s_{\mathcal{P}}(\Omega)$.
\item L'image d'un \textbf{solide} est un solide de \textbf{même volume}.
\end{itemize}
\end{propriete}

Tout est donc conservé... sauf une chose, subtile et essentielle : l'\cle{orientation}. Une réflexion \textbf{renverse l'orientation} de l'espace : l'image d'un trièdre direct $(\vec{i}, \vec{j}, \vec{k})$ est un trièdre \emph{indirect}. C'est la différence fondamentale avec les translations, rotations et demi-tours, qui, eux, la conservent. On dit que la réflexion est un \cle{antidéplacement} (ou isométrie négative), tandis que translations et rotations sont des \cle{déplacements}.

\begin{popfigure}
\begin{center}
\begin{tikzpicture}[scale=0.85, line cap=round, line join=round]
\% Plan vertical (le miroir)
\fill[popblueL, opacity=0.5] (3.4,-0.6) -- (3.4,3.4) -- (5.6,4.4) -- (5.6,0.4) -- cycle;
\draw[popblue, thick] (3.4,-0.6) -- (3.4,3.4) -- (5.6,4.4) -- (5.6,0.4) -- cycle;
\node[popblue] at (4.9,3.6) {$\mathcal{P}$};
\% Cube original (à gauche du plan)
\begin{scope}[shift={(0,0)}]
\draw[popink, thick] (0,0) -- (1.4,0) -- (1.4,1.4) -- (0,1.4) -- cycle;
\draw[popink, thick] (0.5,0.5) -- (1.9,0.5) -- (1.9,1.9) -- (0.5,1.9) -- cycle;
\draw[popink, thick] (0,0) -- (0.5,0.5) (1.4,0) -- (1.9,0.5) (1.4,1.4) -- (1.9,1.9) (0,1.4) -- (0.5,1.9);
\fill[popink] (0,0) circle (2pt) node[below left] {\small $A$};
\fill[popink] (1.4,0) circle (2pt) node[below] {\small $B$};
\fill[popink] (1.4,1.4) circle (2pt) node[right] {\small $C$};
\node[popink] at (0.7,-0.5) {\small Cube direct};
\end{scope}
\% Cube image (à droite du plan, symétrique)
\begin{scope}[shift={(6.8,0)}, xscale=-1]
\draw[poporange, thick] (0,0) -- (1.4,0) -- (1.4,1.4) -- (0,1.4) -- cycle;
\draw[poporange, thick] (0.5,0.5) -- (1.9,0.5) -- (1.9,1.9) -- (0.5,1.9) -- cycle;
\draw[poporange, thick] (0,0) -- (0.5,0.5) (1.4,0) -- (1.9,0.5) (1.4,1.4) -- (1.9,1.9) (0,1.4) -- (0.5,1.9);
\fill[poporange] (0,0) circle (2pt);
\fill[poporange] (1.4,0) circle (2pt);
\fill[poporange] (1.4,1.4) circle (2pt);
\end{scope}
\node[poporange] at (7.5,-0.5) {\small Cube image};
\node[poporange] at (8.6,0.1) {\small $A'$};
\node[poporange] at (7.3,0.1) {\small $B'$};
\node[poporange] at (7.3,1.6) {\small $C'$};
\% Correspondances pointillées
\draw[popmuted, dashed] (0,0) -- (6.8,0);
\draw[popmuted, dashed] (1.4,0) -- (5.4,0);
\draw[popmuted, dashed] (1.4,1.4) -- (5.4,1.4);
\% Flèche orientation
\draw[->, popgreen, thick] (0.7,2.6) arc[start angle=200, end angle=-20, radius=0.55];
\node[popgreen] at (0.7,3.35) {\small sens direct};
\draw[->, poppink, thick] (7.5,2.6) arc[start angle=-20, end angle=200, radius=0.55];
\node[poppink] at (7.5,3.35) {\small sens renversé};
\end{tikzpicture}
\end{center}
\legende{Un cube et son image par la réflexion de plan vertical $\mathcal{P}$ : les longueurs, les angles et le volume sont conservés, mais le sens de parcours des faces est renversé — l'orientation change.}
\end{popfigure}

\begin{savaistu}[Pourquoi le miroir « inverse »-t-il la gauche et la droite ?]
Tiens un texte imprimé devant un miroir : les lettres apparaissent illisibles, « inversées ». On dit souvent que le miroir échange la gauche et la droite. En réalité, il fait mieux (ou pire !) : il échange le \emph{devant} et le \emph{derrière}, c'est-à-dire la direction perpendiculaire à sa surface. Mathématiquement, la réflexion renverse l'orientation de l'espace tout entier : une main droite devient une main gauche dans le miroir (c'est pourquoi ta main droite semble être la main gauche de ton reflet). Cette propriété a des conséquences concrètes : en chimie, deux molécules images l'une de l'autre dans un miroir (dites \emph{énantiomères}) peuvent avoir des effets très différents sur le corps humain ; en mécanique, une pièce et son image miroir ne sont pas interchangeables. Les opticiens de Douala qui montent les miroirs des cabines d'essayage, ou les techniciens qui alignent les réflecteurs des antennes, utilisent sans le savoir la théorie des réflexions de l'espace !
\end{savaistu}

\begin{exoresolu}[Exploiter les propriétés de conservation]
Soit $\mathcal{P}$ un plan et $s_{\mathcal{P}}$ la réflexion associée. Soit $\mathcal{S}$ la sphère de centre $\Omega(2; 1; 4)$ et de rayon $R = 3$, où $\mathcal{P}$ est le plan d'équation $z = 0$.
\begin{enumerate}
\item Déterminer l'image $\mathcal{S}'$ de $\mathcal{S}$ par $s_{\mathcal{P}}$.
\item Les sphères $\mathcal{S}$ et $\mathcal{S}'$ sont-elles sécantes ? Si oui, préciser la nature de leur intersection.
\end{enumerate}
\tcblower
\textbf{Solution détaillée.}
\textbf{1.} La réflexion conserve les distances, donc l'image d'une sphère est une sphère de \textbf{même rayon} $R' = 3$, centrée en $\Omega' = s_{\mathcal{P}}(\Omega)$. La réflexion par rapport au plan $z = 0$ transforme $(x;y;z)$ en $(x;y;-z)$, donc $\Omega'(2; 1; -4)$. Ainsi :
\[ \mathcal{S}' : \quad (x-2)^2 + (y-1)^2 + (z+4)^2 = 9. \]
\textbf{2.} Calculons la distance des centres : $\Omega\Omega' = |4 - (-4)| = 8$. Or $R + R' = 6$. Comme $\Omega\Omega' = 8 > 6 = R + R'$, les deux sphères sont \textbf{extérieures l'une à l'autre} : elles ne sont pas sécantes, leur intersection est vide. \emph{Remarque :} si $\Omega$ avait été à une distance de $\mathcal{P}$ inférieure à $3$ (par exemple $\Omega(2;1;2)$, donnant $\Omega\Omega' = 4 < 6$), les sphères se seraient coupées selon un \textbf{cercle} contenu dans un plan parallèle à $\mathcal{P}$.
\end{exoresolu}

\begin{aretenir}[L'essentiel de la section]
\begin{itemize}
\item Le \cle{projeté orthogonal} de $M$ sur un plan $\mathcal{P}$ est le point $H$ tel que $(MH) \perp \mathcal{P}$ ; il réalise la distance minimale : $d(M, \mathcal{P}) = MH$ (justification par Pythagore).
\item La \cle{réflexion} $s_{\mathcal{P}}$ associe à $M$ le point $M'$ tel que $\mathcal{P}$ soit le plan médiateur de $[MM']$ ; vectoriellement : $\overrightarrow{HM'} = -\overrightarrow{HM}$.
\item $s_{\mathcal{P}}$ est une \cle{involution} : $s_{\mathcal{P}} \circ s_{\mathcal{P}} = \mathrm{Id}$ ; ses seuls points fixes sont les points de $\mathcal{P}$.
\item $s_{\mathcal{P}}$ est une \cle{isométrie} : elle conserve distances, milieux, alignement, produit scalaire, angles, aires et volumes.
\item Mais elle \cle{renverse l'orientation} : c'est un antidéplacement — d'où l'effet « miroir ».
\item Ne jamais confondre réflexion (plan), demi-tour (droite) et symétrie centrale (point).
\end{itemize}
\end{aretenir}

\begin{exercice}[Pour t'entraîner]
L'espace est muni d'un repère orthonormé. Soit $\mathcal{P}$ le plan d'équation $y = 2$ et $A(3; 5; -1)$, $B(1; 2; 4)$.
\begin{enumerate}
\item Déterminer les images $A'$ et $B'$ de $A$ et $B$ par $s_{\mathcal{P}}$.
\item Vérifier que $A'B' = AB$.
\item Soit $I$ le milieu de $[AB]$. Vérifier que $s_{\mathcal{P}}(I)$ est le milieu de $[A'B']$.
\end{enumerate}
\end{exercice}

\begin{corrige}[Exercice 1]
\textbf{1.} Le plan $\mathcal{P} : y = 2$ est parallèle au plan $(O;\vec{i},\vec{k})$. La perpendiculaire à $\mathcal{P}$ est dirigée par $\vec{j}$ : seule l'ordonnée change. Le projeté de $A(3;5;-1)$ est $H_A(3;2;-1)$ et $\overrightarrow{H_A A'} = -\overrightarrow{H_A A} = (0;-3;0)$, d'où $A'(3; -1; -1)$. Pour $B(1;2;4) \in \mathcal{P}$ (car $y = 2$) : $B' = B(1;2;4)$.
\textbf{2.} $AB = \sqrt{(1-3)^2 + (2-5)^2 + (4+1)^2} = \sqrt{4 + 9 + 25} = \sqrt{38}$. $A'B' = \sqrt{(1-3)^2 + (2+1)^2 + (4+1)^2} = \sqrt{4+9+25} = \sqrt{38}$. L'égalité est vérifiée.
\textbf{3.} $I = \left(2; \frac{7}{2}; \frac{3}{2}\right)$, donc $I' = \left(2; \frac{1}{2}; \frac{3}{2}\right)$. Le milieu de $[A'B']$ est $\left(\frac{3+1}{2}; \frac{-1+2}{2}; \frac{-1+4}{2}\right) = \left(2; \frac{1}{2}; \frac{3}{2}\right) = I'$. La conservation des milieux est confirmée.
\end{corrige}

Dans la section suivante, nous traduirons tout cela en formules : l'\emph{expression analytique} d'une réflexion dans un repère orthonormé, outil indispensable pour les exercices du Baccalauréat.

\coursec{Expression analytique d'une réflexion}

Dans la section précédente, nous avons défini la réflexion de l'espace de manière purement géométrique : étant donné un plan $P$, l'image d'un point $M$ est le point $M'$ tel que $P$ soit le plan médiateur du segment $[MM']$, c'est-à-dire tel que la droite $(MM')$ soit perpendiculaire à $P$ et que le milieu de $[MM']$ appartienne à $P$. Cette définition est belle, mais pour calculer concrètement — et au Baccalauréat, on calcule ! — il nous faut traduire cette géométrie en \cle{coordonnées}. C'est toute l'ambition de cette section : obtenir une \cle{formule} qui, à partir des coordonnées $(x\,;y\,;z)$ d'un point $M$ et de l'équation du plan $P$, donne directement les coordonnées $(x'\,;y'\,;z')$ de son image $M'$.

\courssub{Cadre de travail : le repère orthonormé}

\begin{aretenir}[Cadre obligatoire]
Toute cette section se déroule dans l'espace muni d'un repère \cle{orthonormé} $(O\,;\vec{i},\vec{j},\vec{k})$ : les trois axes sont deux à deux perpendiculaires et $\|\vec{i}\|=\|\vec{j}\|=\|\vec{k}\|=1$.
\end{aretenir}

Pourquoi cette exigence ? Parce que la réflexion est une symétrie \emph{orthogonale} : la notion de perpendicularité et la notion de distance y jouent un rôle central. Or, dans un repère quelconque, le produit scalaire de deux vecteurs ne s'écrit pas simplement $xx'+yy'+zz'$, et la distance $MH$ ne se lit pas sur les coordonnées. C'est précisément le caractère orthonormé du repère qui autorise les formules simples que nous allons établir.

\begin{attention}[Vérifier le repère avant toute chose]
La formule de la réflexion que nous allons démontrer \textbf{exige un repère orthonormé}. Avant d'appliquer la moindre formule, lis attentivement l'énoncé : il doit préciser « l'espace est muni d'un repère orthonormé ». Si le repère est seulement « orthogonal » ou si rien n'est dit, la formule est en principe inapplicable. Au Baccalauréat, cette précision figure presque toujours dans l'énoncé : repère-la et cite-la dans ta rédaction.
\end{attention}

\courssub{Cas fondamental : les plans de coordonnées}

Commençons par les situations les plus simples, celles où le plan de réflexion est l'un des trois plans de coordonnées. L'intuition est immédiate : se réfléchir par rapport au plan horizontal $(xOy)$, c'est « retourner » l'espace comme une crêpe au-dessus de ce plan — les coordonnées horizontales ne bougent pas, seule la hauteur change de signe.

\begin{propriete}[Réflexions par rapport aux plans de coordonnées]
Dans un repère orthonormé $(O\,;\vec{i},\vec{j},\vec{k})$, soit $M(x\,;y\,;z)$ un point de l'espace et $M'(x'\,;y'\,;z')$ son image.
\begin{itemize}
\item Par la réflexion de plan $(xOy)$ : $\begin{cases} x'=x \\ y'=y \\ z'=-z \end{cases}$
\item Par la réflexion de plan $(xOz)$ : $\begin{cases} x'=x \\ y'=-y \\ z'=z \end{cases}$
\item Par la réflexion de plan $(yOz)$ : $\begin{cases} x'=-x \\ y'=y \\ z'=z \end{cases}$
\end{itemize}
\end{propriete}

\begin{experience}[Justification dans le cas du plan $(xOy)$]
Soit $M(x\,;y\,;z)$ et $M'(x'\,;y'\,;z')$ son image par la réflexion $s$ de plan $(xOy)$.
\begin{enumerate}
\item Le projeté orthogonal de $M$ sur $(xOy)$ est le point $H(x\,;y\,;0)$ : c'est le point du plan situé « à la verticale » de $M$, puisque l'axe $(Oz)$ est perpendiculaire au plan $(xOy)$.
\item Par définition de la réflexion, $H$ est le milieu de $[MM']$. Ses coordonnées sont donc $\left(\dfrac{x+x'}{2}\,;\dfrac{y+y'}{2}\,;\dfrac{z+z'}{2}\right)$.
\item En identifiant : $\dfrac{x+x'}{2}=x$ donne $x'=x$ ; $\dfrac{y+y'}{2}=y$ donne $y'=y$ ; $\dfrac{z+z'}{2}=0$ donne $z'=-z$. \quad$\blacksquare$
\end{enumerate}
\end{experience}

\begin{exemplebox}[Lecture directe]
L'image du point $B(4\,;-1\,;7)$ par la réflexion de plan $(xOy)$ est $B'(4\,;-1\,;-7)$. L'image de ce même point par la réflexion de plan $(yOz)$ est $B''(-4\,;-1\,;7)$.
\end{exemplebox}

\begin{attention}[Quelle coordonnée change de signe ?]
C'est la coordonnée \textbf{absente du nom du plan} qui change de signe : le plan $(xOy)$ ne contient pas l'axe $(Oz)$, donc c'est $z$ qui devient $-z$. Une erreur classique consiste à changer le signe d'une mauvaise coordonnée : retiens l'astuce « le plan garde ses deux lettres, la troisième coordonnée bascule ».
\end{attention}

\courssub{Cas général : réflexion par rapport à un plan quelconque}

Passons maintenant au cœur de la section. Soit $P$ un plan d'équation cartésienne $ax+by+cz+d=0$, de vecteur normal $\vec{n}(a\,;b\,;c)$. Soit $M(x_M\,;y_M\,;z_M)$ un point, $H$ son projeté orthogonal sur $P$ et $M'$ son image par la réflexion $s_P$.

\textbf{Intuition.} Pour aller de $M$ à son image $M'$, on se déplace \emph{uniquement le long de la normale} au plan : on avance perpendiculairement à $P$ jusqu'à toucher le plan (on arrive en $H$), puis on continue de la même distance de l'autre côté. Le déplacement total $\overrightarrow{MM'}$ est donc colinéaire à $\vec{n}$, et sa « longueur algébrique » vaut deux fois la distance orientée de $M$ au plan, avec un signe moins car on va vers le plan.

\begin{propriete}[Expression analytique de la réflexion de plan $P$]
Dans un repère orthonormé, soit $P$ le plan d'équation $ax+by+cz+d=0$ et $\vec{n}(a\,;b\,;c)$ un vecteur normal à $P$. Pour tout point $M(x_M\,;y_M\,;z_M)$, on pose
\[ k=\dfrac{ax_M+by_M+cz_M+d}{a^2+b^2+c^2}. \]
Alors le projeté orthogonal $H$ de $M$ sur $P$ est donné par
\[ \overrightarrow{MH}=-k\,\vec{n} \qquad\text{c'est-à-dire}\qquad H=M-k\,\vec{n}, \]
et l'image $M'$ de $M$ par la réflexion de plan $P$ est donnée par
\[ \overrightarrow{MM'}=-2k\,\vec{n} \qquad\text{c'est-à-dire}\qquad \boxed{\;M'=M-2\,\dfrac{ax_M+by_M+cz_M+d}{a^2+b^2+c^2}\,\vec{n}\;} \]
soit, en coordonnées :
\[ \begin{cases} x'=x_M-2ka \\ y'=y_M-2kb \\ z'=z_M-2kc \end{cases} \]
Cette démonstration est \textbf{exigible} : elle est guidée ci-dessous.
\end{propriete}

\begin{experience}[Démonstration guidée de la formule]
On cherche le point $M'$ tel que $\overrightarrow{MM'}=-2k\,\vec{n}$ soit l'image de $M$. Procédons par étapes.

\textbf{Étape 1 : le déplacement est colinéaire à $\vec{n}$.}
Par définition de la réflexion, la droite $(MM')$ est perpendiculaire à $P$, donc dirigée par $\vec{n}$. Il existe donc un réel $t$ tel que $\overrightarrow{MM'}=t\,\vec{n}$. Tout le problème est de trouver $t$.

\textbf{Étape 2 : le milieu de $[MM']$ appartient à $P$.}
Notons $I$ le milieu de $[MM']$. Alors $\overrightarrow{MI}=\dfrac{1}{2}\overrightarrow{MM'}=\dfrac{t}{2}\,\vec{n}$, donc les coordonnées de $I$ sont :
\[ I\left(x_M+\dfrac{ta}{2}\;;\;y_M+\dfrac{tb}{2}\;;\;z_M+\dfrac{tc}{2}\right). \]
Par définition de la réflexion, $I\in P$, donc ses coordonnées vérifient l'équation du plan :
\[ a\left(x_M+\dfrac{ta}{2}\right)+b\left(y_M+\dfrac{tb}{2}\right)+c\left(z_M+\dfrac{tc}{2}\right)+d=0. \]

\textbf{Étape 3 : résolution.}
En développant et en regroupant :
\[ (ax_M+by_M+cz_M+d)+\dfrac{t}{2}(a^2+b^2+c^2)=0. \]
Comme $\vec{n}\neq\vec{0}$, on a $a^2+b^2+c^2>0$, et l'on peut diviser :
\[ \dfrac{t}{2}=-\dfrac{ax_M+by_M+cz_M+d}{a^2+b^2+c^2}=-k \qquad\text{donc}\qquad t=-2k. \]

\textbf{Étape 4 : conclusion.}
$\overrightarrow{MM'}=-2k\,\vec{n}$, ce qui donne exactement les formules annoncées. Le projeté $H$ correspond au milieu du trajet : $\overrightarrow{MH}=-k\,\vec{n}$. \quad$\blacksquare$
\end{experience}

\begin{aretenir}[Lecture géométrique du coefficient $k$]
Le nombre $k=\dfrac{ax_M+by_M+cz_M+d}{\|\vec{n}\|^2}$ mesure la « distance algébrique » de $M$ au plan, comptée dans l'unité $\|\vec{n}\|$ le long de la normale. En particulier :
\begin{itemize}
\item $k=0$ signifie que $M\in P$ : le point est alors \textbf{invariant}, $M'=M$ ;
\item la distance de $M$ au plan vaut $MH=|k|\cdot\|\vec{n}\|=\dfrac{|ax_M+by_M+cz_M+d|}{\sqrt{a^2+b^2+c^2}}$.
\end{itemize}
\end{aretenir}

\begin{popfigure}
\begin{tikzpicture}[scale=1.05, >=Stealth]
\fill[popblueL, opacity=0.55] (-3.4,-1.3) -- (3.4,-1.3) -- (4.4,0.9) -- (-2.4,0.9) -- cycle;
\draw[popblue, thick] (-3.4,-1.3) -- (3.4,-1.3) -- (4.4,0.9) -- (-2.4,0.9) -- cycle;
\node[popblue] at (-2.6,0.55) {$P : ax+by+cz+d=0$};
\coordinate (M) at (0.4,3.1);
\coordinate (H) at (0.4,0.15);
\coordinate (Mp) at (0.4,-2.8);
\fill[popink] (M) circle (2pt) node[above right] {$M(x_M;y_M;z_M)$};
\fill[popgreen] (H) circle (2pt) node[right=2pt] {$H=M-k\vec{n}$};
\fill[poporange] (Mp) circle (2pt) node[below right] {$M'=M-2k\vec{n}$};
\draw[popink, thick, ->] (M) -- (H) node[midway, right] {$-k\vec{n}$};
\draw[poporange, thick, ->] (H) -- (Mp) node[midway, right] {$-k\vec{n}$};
\draw[popdark, dashed] (M) -- (Mp);
\draw[poppurple, very thick, ->] (2.2,0.15) -- (2.2,1.65) node[above] {$\vec{n}(a;b;c)$};
\draw[poppurple] (2.2,0.15) -- (2.55,0.15) -- (2.55,0.5);
\draw[popdark] (0.4,0.15) ++(0,0.28) -- ++(0.28,0.28) -- ++(0.28,0);
\draw[popmuted, ->] (-3.1,-2.6) -- (-2.1,-2.6) node[right] {\footnotesize $\vec{j}$};
\draw[popmuted, ->] (-3.1,-2.6) -- (-3.1,-1.6) node[above] {\footnotesize $\vec{k}$};
\draw[popmuted, ->] (-3.1,-2.6) -- (-3.75,-3.15) node[below] {\footnotesize $\vec{i}$};
\fill[popmuted] (-3.1,-2.6) circle (1.2pt) node[below left] {\footnotesize $O$};
\end{tikzpicture}
\legende{Réflexion de plan $P$ : le vecteur $\overrightarrow{MM'}$ est colinéaire au vecteur normal $\vec{n}$ et vaut $-2k\,\vec{n}$, où $k=\dfrac{ax_M+by_M+cz_M+d}{a^2+b^2+c^2}$. Le projeté $H$ est le milieu de $[MM']$.}
\end{popfigure}

\courssub{Méthode de calcul et exemple complet}

\begin{methode}[Obtenir l'image d'un point par une réflexion en trois étapes]
Pour déterminer l'image $M'$ d'un point $M(x_M\,;y_M\,;z_M)$ par la réflexion de plan $P : ax+by+cz+d=0$ :
\begin{enumerate}
\item \textbf{Identifier} le vecteur normal $\vec{n}(a\,;b\,;c)$ en lisant les coefficients de l'équation du plan, puis calculer $\|\vec{n}\|^2=a^2+b^2+c^2$.
\item \textbf{Calculer le coefficient} $k=\dfrac{ax_M+by_M+cz_M+d}{a^2+b^2+c^2}$ : on remplace simplement $x$, $y$, $z$ par les coordonnées de $M$ dans le membre de gauche de l'équation.
\item \textbf{Soustraire $2k\,\vec{n}$} aux coordonnées de $M$ : $x'=x_M-2ka$, $y'=y_M-2kb$, $z'=z_M-2kc$.
\end{enumerate}
\end{methode}

\begin{exoresolu}[Image de $A(2\,;0\,;3)$ par la réflexion de plan $P : x+y-z+1=0$]
L'espace est muni d'un repère orthonormé $(O\,;\vec{i},\vec{j},\vec{k})$. On considère le plan $P$ d'équation $x+y-z+1=0$ et le point $A(2\,;0\,;3)$. Déterminer les coordonnées de $A'$, image de $A$ par la réflexion de plan $P$.
\tcblower
\textbf{Étape 1 : vecteur normal.} L'équation $x+y-z+1=0$ donne $\vec{n}(1\,;1\,;-1)$, donc
\[ \|\vec{n}\|^2=1^2+1^2+(-1)^2=3. \]
\textbf{Étape 2 : coefficient $k$.} On évalue le membre de gauche de l'équation en $A$ :
\[ ax_A+by_A+cz_A+d=2+0-3+1=0. \]
Donc $k=\dfrac{0}{3}=0$. Le point $A$ appartient au plan $P$ (vérification : $2+0-3+1=0$), il est donc invariant : $A'=A(2\,;0\,;3)$.

Pour illustrer pleinement la méthode sur un point hors du plan, considérons le point $B(2\,;0\,;4)$ :
\[ k=\dfrac{2+0-4+1}{3}=\dfrac{-1}{3}. \]
\textbf{Étape 3 : coordonnées de $B'$.} On soustrait $2k\,\vec{n}=-\dfrac{2}{3}\vec{n}$ :
\begin{align*}
x' &= 2-2\times\left(-\dfrac{1}{3}\right)\times 1 = 2+\dfrac{2}{3}=\dfrac{8}{3},\\
y' &= 0-2\times\left(-\dfrac{1}{3}\right)\times 1 = \dfrac{2}{3},\\
z' &= 4-2\times\left(-\dfrac{1}{3}\right)\times(-1) = 4-\dfrac{2}{3}=\dfrac{10}{3}.
\end{align*}
Conclusion : $B'\left(\dfrac{8}{3}\,;\dfrac{2}{3}\,;\dfrac{10}{3}\right)$.
\end{exoresolu}

\begin{exoresolu}[Un calcul complet sans point invariant]
Dans un repère orthonormé, soit $P$ le plan d'équation $x+y-z+1=0$ et $M(1\,;2\,;1)$. Déterminer l'image $M'$ de $M$ par la réflexion de plan $P$, puis vérifier le résultat.
\tcblower
\textbf{Calcul.} Vecteur normal : $\vec{n}(1\,;1\,;-1)$, $\|\vec{n}\|^2=3$.
\[ k=\dfrac{1+2-1+1}{3}=\dfrac{3}{3}=1. \]
On soustrait $2k\,\vec{n}=2\vec{n}=(2\,;2\,;-2)$ :
\[ x'=1-2=-1,\qquad y'=2-2=0,\qquad z'=1-(-2)=3. \]
Donc $M'(-1\,;0\,;3)$.

\textbf{Vérification 1 : le milieu de $[MM']$ est sur $P$.} Le milieu est $I\left(\dfrac{1+(-1)}{2}\,;\dfrac{2+0}{2}\,;\dfrac{1+3}{2}\right)=(0\,;1\,;2)$. Testons : $0+1-2+1=0$. \checkmark{} Le point $I$ appartient bien à $P$.

\textbf{Vérification 2 : $\overrightarrow{MM'}$ est colinéaire à $\vec{n}$.} $\overrightarrow{MM'}(-2\,;-2\,;2)=-2\,(1\,;1\,;-1)=-2\vec{n}$. \checkmark{} Le vecteur est bien colinéaire au vecteur normal.

Les deux conditions de la définition sont satisfaites : le résultat est confirmé.
\end{exoresolu}

\courssub{Vérification du résultat : le réflexe qui sauve des points}

\begin{methode}[Astuce Bac : le double contrôle]
Après avoir calculé l'image $M'$ d'un point $M$ par une réflexion de plan $P$, effectue systématiquement les deux vérifications suivantes, qui découlent directement de la définition :
\begin{enumerate}
\item \textbf{Le milieu $I$ de $[MM']$ appartient à $P$} : remplace ses coordonnées dans l'équation du plan, tu dois obtenir $0$ ;
\item \textbf{Le vecteur $\overrightarrow{MM'}$ est colinéaire à $\vec{n}$} : les coordonnées de $\overrightarrow{MM'}$ doivent être proportionnelles à $(a\,;b\,;c)$.
\end{enumerate}
Ce double contrôle prend moins d'une minute et détecte la quasi-totalité des erreurs de calcul. Rédige-le explicitement : un correcteur de Baccalauréat apprécie toujours un candidat qui vérifie ses résultats.
\end{methode}

\begin{attention}[Pièges fréquents]
\begin{itemize}
\item \textbf{Oublier le facteur 2} : la formule du projeté utilise $-k\vec{n}$, celle de la réflexion utilise $-2k\vec{n}$. Confondre les deux donne le projeté $H$ au lieu de l'image $M'$.
\item \textbf{Oublier le terme constant $d$} dans le numérateur de $k$ : c'est $ax_M+by_M+cz_M\,\textbf{+}\,d$, pas seulement $ax_M+by_M+cz_M$.
\item \textbf{Mauvais signe sur $c$} : pour le plan $x+y-z+1=0$, le vecteur normal est $(1\,;1\,;-1)$ et non $(1\,;1\,;1)$. Lis les coefficients avec leurs signes.
\item \textbf{Diviser par $\|\vec{n}\|$ au lieu de $\|\vec{n}\|^2$} : le dénominateur de $k$ est $a^2+b^2+c^2$ (le carré de la norme), pas $\sqrt{a^2+b^2+c^2}$.
\end{itemize}
\end{attention}

\begin{savaistu}[Des miroirs aux scanners]
La formule $M'=M-2k\vec{n}$ est l'un des outils de base de l'infographie et de l'imagerie médicale : lorsqu'un logiciel de scanner ou d'IRM reconstruit une image 3D du corps humain, il exploite en permanence des symétries par rapport à des plans de coupe. Plus près de nous, les ingénieurs qui conçoivent les antennes relais de téléphonie mobile (MTN, Orange, Camtel) utilisent les réflexions de l'espace pour modéliser les ondes réfléchies par les façades des immeubles : chaque réflexion d'une onde sur un mur suit exactement la loi géométrique que tu viens d'étudier.
\end{savaistu}

\begin{exercice}[À toi de jouer]
L'espace est muni d'un repère orthonormé $(O\,;\vec{i},\vec{j},\vec{k})$. Soit $P$ le plan d'équation $2x-y+2z-3=0$.
\begin{enumerate}
\item Déterminer l'image $M'$ du point $M(1\,;1\,;1)$ par la réflexion de plan $P$.
\item Vérifier que le milieu de $[MM']$ appartient à $P$ et que $\overrightarrow{MM'}$ est colinéaire au vecteur normal $\vec{n}$.
\item Le point $N(2\,;1\,;0)$ appartient-il à $P$ ? Que peut-on en déduire pour son image ?
\end{enumerate}
\end{exercice}

\begin{corrige}[Exercice]
\begin{enumerate}
\item Vecteur normal : $\vec{n}(2\,;-1\,;2)$, donc $\|\vec{n}\|^2=4+1+4=9$. Coefficient :
\[ k=\dfrac{2(1)-1+2(1)-3}{9}=\dfrac{2-1+2-3}{9}=0. \]
Donc $M\in P$ (vérification directe : $2-1+2-3=0$) et $M'=M(1\,;1\,;1)$.
\item Ici $M'=M$, le milieu de $[MM']$ est $M$ lui-même, qui appartient à $P$, et $\overrightarrow{MM'}=\vec{0}$ est colinéaire à tout vecteur. La vérification est immédiate mais le cas est dégénéré : c'est le cas d'un point invariant.
\item Pour $N(2\,;1\,;0)$ : $2(2)-1+2(0)-3=4-1-3=0$, donc $N\in P$ également. Son image est donc $N$ lui-même : tout point du plan de réflexion est invariant.
\end{enumerate}
\end{corrige}

En résumé, tu disposes désormais d'un outil de calcul complet : une équation de plan et trois étapes mécaniques suffisent pour déterminer l'image de n'importe quel point par une réflexion. Dans la section suivante, nous composerons deux réflexions et découvrirons que ces composées cachent deux autres isométries fondamentales : la translation et le demi-tour.

\coursec{Composée de deux réflexions}

Dans la section précédente, nous avons appris à écrire l'expression analytique d'une réflexion $s_P$ par rapport à un plan $P$ donné par son équation cartésienne. Nous savons donc maintenant calculer l'image d'un point par une réflexion. Mais la géométrie de l'espace devient véritablement puissante lorsqu'on \emph{enchaîne} les transformations : que se passe-t-il si l'on applique une réflexion, \textbf{puis} une autre ? C'est toute la question de cette section, et tu vas découvrir un résultat magnifique : les réflexions, composées deux à deux, fabriquent les deux autres isométries fondamentales que sont la \cle{translation} et le \cle{demi-tour}.

\begin{savaistu}[Une intuition de la vie quotidienne]
Imagine un objet placé entre deux miroirs plans et parallèles, comme dans certains salons de coiffure de Yaoundé où deux glaces se font face : tu vois l'objet se répéter à l'infini, chaque image étant \emph{décalée} de la même quantité. Ce décalage régulier, c'est exactement une translation ! De même, deux miroirs disposés en angle droit (comme dans un coin de salle de bain) te montrent une image « retournée » d'un demi-tour. La physique des miroirs illustre parfaitement les théorèmes que nous allons démontrer.
\end{savaistu}

\courssub{Rappel : position relative de deux plans de l'espace}

Avant de composer deux réflexions $s_P$ et $s_Q$, il faut d'abord savoir comment les plans $P$ et $Q$ sont situés l'un par rapport à l'autre. Toute la suite en dépend.

\begin{propriete}[Position relative de deux plans]
Soient $P$ et $Q$ deux plans de l'espace. Trois cas, et trois seulement, sont possibles :
\begin{itemize}
\item $P$ et $Q$ sont \textbf{confondus} : $P = Q$ ;
\item $P$ et $Q$ sont \textbf{strictement parallèles} : $P \cap Q = \varnothing$ ;
\item $P$ et $Q$ sont \textbf{sécants} : leur intersection est une droite $(d)$. Le cas particulier important est celui où $P$ et $Q$ sont \textbf{perpendiculaires}.
\end{itemize}
Analytiquement, si $P : ax+by+cz+d=0$ et $Q : a'x+b'y+c'z+d'=0$, alors $P$ et $Q$ sont parallèles (confondus ou strictement) si et seulement si leurs vecteurs normaux $\vec{n}(a\,;b\,;c)$ et $\vec{n}'(a'\,;b'\,;c')$ sont colinéaires.
\end{propriete}

\begin{attention}[Cas trivial à évacuer immédiatement]
Si $P = Q$, alors $s_Q \circ s_P = s_P \circ s_P = \mathrm{Id}$ (l'identité de l'espace) : appliquer deux fois la même réflexion ramène chaque point à sa position initiale. Dans les exercices, on suppose donc toujours $P \neq Q$, et l'on se ramène aux deux cas intéressants : plans \textbf{strictement parallèles} ou plans \textbf{sécants}.
\end{attention}

\courssub{Composée de deux réflexions de plans strictement parallèles}

Commençons par une expérience de pensée. Soient $P$ et $Q$ deux plans strictement parallèles. Prenons un point $M$ quelconque, notons $M' = s_P(M)$ son symétrique par rapport à $P$, puis $M'' = s_Q(M')$ le symétrique de $M'$ par rapport à $Q$. La question est : comment passe-t-on \emph{directement} de $M$ à $M''$ ?

Observons d'abord que la droite $(MM')$ est perpendiculaire à $P$, et la droite $(M'M'')$ est perpendiculaire à $Q$. Comme $P \parallel Q$, ces deux directions sont les mêmes : les points $M$, $M'$ et $M''$ sont \textbf{alignés} sur une même droite perpendiculaire aux deux plans. Le déplacement de $M$ vers $M''$ se fait donc le long d'une direction fixe, celle de la normale commune. Reste à calculer l'amplitude de ce déplacement : c'est l'objet du théorème suivant.

\begin{propriete}[Théorème 1 : composée de deux réflexions de plans parallèles]
Soient $P$ et $Q$ deux plans \textbf{strictement parallèles}. Soient $H_1$ un point de $P$ et $H_2$ son projeté orthogonal sur $Q$ (de sorte que $\overrightarrow{H_1H_2}$ est perpendiculaire aux deux plans). Alors :
\[ s_Q \circ s_P = t_{\,2\overrightarrow{H_1H_2}} \]
c'est-à-dire la \cle{translation} de vecteur $2\overrightarrow{H_1H_2}$. Ce vecteur ne dépend pas du choix de $H_1$ : c'est le double de n'importe quel vecteur allant perpendiculairement de $P$ vers $Q$.
\par\textbf{Démonstration exigible au Bac.}
\end{propriete}

\begin{experience}[Démonstration guidée du Théorème 1]
Soit $M$ un point quelconque de l'espace. Notons $M' = s_P(M)$ et $M'' = s_Q(M')$. Notons $I$ le milieu de $[MM']$ (c'est le projeté orthogonal de $M$ sur $P$, donc $I \in P$) et $J$ le milieu de $[M'M'']$ (c'est le projeté orthogonal de $M'$ sur $Q$, donc $J \in Q$).
\begin{enumerate}
\item \textbf{Alignement.} $(MM') \perp P$ et $(M'M'') \perp Q$. Comme $P \parallel Q$, les droites $(MM')$ et $(M'M'')$, qui passent toutes deux par $M'$ et sont perpendiculaires à la même direction de plans, sont confondues. Donc $M$, $I$, $M'$, $J$, $M''$ sont alignés.
\item \textbf{Relation de Chasles.} Décomposons le vecteur $\overrightarrow{MM''}$ en passant par $M'$ :
\[ \overrightarrow{MM''} = \overrightarrow{MM'} + \overrightarrow{M'M''}. \]
\item \textbf{Utilisation des milieux.} Puisque $I$ est le milieu de $[MM']$, $\overrightarrow{MM'} = 2\overrightarrow{IM'}$. Puisque $J$ est le milieu de $[M'M'']$, $\overrightarrow{M'M''} = 2\overrightarrow{M'J}$. Donc :
\[ \overrightarrow{MM''} = 2\overrightarrow{IM'} + 2\overrightarrow{M'J} = 2\left(\overrightarrow{IM'} + \overrightarrow{M'J}\right) = 2\overrightarrow{IJ}. \]
\item \textbf{Conclusion.} Le vecteur $\overrightarrow{IJ}$ joint un point de $P$ à son projeté orthogonal sur $Q$ : il est égal à $\overrightarrow{H_1H_2}$ (tous les vecteurs perpendiculaires joignant $P$ à $Q$ sont égaux, car $P \parallel Q$). Ainsi, pour \emph{tout} point $M$ :
\[ \overrightarrow{MM''} = 2\overrightarrow{IJ} = 2\overrightarrow{H_1H_2}, \]
ce qui prouve que $s_Q \circ s_P$ est la translation de vecteur $2\overrightarrow{H_1H_2}$. \hfill$\blacksquare$
\end{enumerate}
\end{experience}

\begin{popfigure}
\begin{center}
\begin{tikzpicture}[scale=1.05, >=Stealth]
% Plan P
\draw[fill=popblueL, draw=popblue, thick] (-0.5,-0.3) -- (1.8,-1.3) -- (4.6,-1.3) -- (2.3,-0.3) -- cycle;
\node[popblue] at (4.1,-0.5) {$P$};
% Plan Q
\draw[fill=poporangeL, draw=poporange, thick] (-0.5,1.6) -- (1.8,0.6) -- (4.6,0.6) -- (2.3,1.6) -- cycle;
\node[poporange] at (4.1,1.4) {$Q$};
% Points M, M', M''
\coordinate (M) at (1.0,-1.9);
\coordinate (Mp) at (1.0,-0.35);
\coordinate (Mpp) at (1.0,1.2);
\coordinate (I) at (1.0,-0.85);
\coordinate (J) at (1.0,0.65);
\draw[popline, dashed] (M) -- (Mpp);
\filldraw[popdark] (M) circle (1.6pt) node[below left] {$M$};
\filldraw[popgreen] (Mp) circle (1.6pt) node[right] {$M'$};
\filldraw[poppink] (Mpp) circle (1.6pt) node[above left] {$M''$};
\filldraw[popblue] (I) circle (1.2pt) node[left] {\footnotesize $I$};
\filldraw[poporange] (J) circle (1.2pt) node[left] {\footnotesize $J$};
% flèches de symétrie
\draw[->, popgreen, thick] (0.75,-1.75) to[bend right=25] (0.75,-0.5);
\draw[->, popgreen, thick] (0.75,-0.2) to[bend right=25] (0.75,1.05);
% vecteur double
\draw[->, poppurple, very thick] (1.35,-1.9) -- (1.35,1.2) node[midway, right] {\footnotesize $2\overrightarrow{H_1H_2}$};
% vecteur H1H2
\coordinate (H1) at (3.1,-0.85);
\coordinate (H2) at (3.1,0.65);
\draw[->, popdark, thick] (H1) -- (H2) node[midway, right] {\footnotesize $\overrightarrow{H_1H_2}$};
\filldraw[popdark] (H1) circle (1.2pt) node[below] {\footnotesize $H_1$};
\filldraw[popdark] (H2) circle (1.2pt) node[above] {\footnotesize $H_2$};
\end{tikzpicture}
\end{center}
\legende{$P \parallel Q$ : la composée $s_Q \circ s_P$ envoie $M$ sur $M''$ par la translation de vecteur $2\overrightarrow{H_1H_2}$.}
\end{popfigure}

\begin{exemplebox}[Exemple chiffré fondamental]
L'espace est muni d'un repère orthonormé $(O\,;\vec{i},\vec{j},\vec{k})$. Soient $P : z = 0$ (le plan horizontal) et $Q : z = 3$. Déterminons la nature de $s_Q \circ s_P$.
\begin{itemize}
\item Les deux plans ont pour vecteur normal commun $\vec{k}(0\,;0\,;1)$ : ils sont strictement parallèles.
\item Prenons $H_1(0\,;0\,;0) \in P$ ; son projeté orthogonal sur $Q$ est $H_2(0\,;0\,;3)$. Donc $\overrightarrow{H_1H_2} = 3\vec{k}$ et $2\overrightarrow{H_1H_2} = 6\vec{k}$, de coordonnées $(0\,;0\,;6)$.
\item Vérifions par le calcul analytique. La réflexion par rapport à $P : z=0$ donne $M'(x\,;y\,;-z)$, puis la réflexion par rapport à $Q : z=3$ donne $M''(x\,;y\,;6-(-z)) = (x\,;y\,;z+6)$. On obtient bien $\overrightarrow{MM''}(0\,;0\,;6)$ : $s_Q \circ s_P$ est la translation de vecteur $6\vec{k}$. $\checkmark$
\end{itemize}
\end{exemplebox}

\begin{attention}[L'ordre de composition compte !]
Contrairement à une idée reçue, la composée d'isométries n'est \textbf{pas commutative} en général. Si $P \parallel Q$, alors :
\[ s_Q \circ s_P = t_{\,2\overrightarrow{H_1H_2}} \qquad \text{mais} \qquad s_P \circ s_Q = t_{\,2\overrightarrow{H_2H_1}} = t_{\,-2\overrightarrow{H_1H_2}}. \]
Les deux composées sont des translations de vecteurs \textbf{opposés}. Dans l'exemple ci-dessus, $s_P \circ s_Q$ est la translation de vecteur $(0\,;0\,;-6)$. Retiens la règle pratique : dans $s_Q \circ s_P$, on applique \emph{d'abord} $s_P$, et le vecteur de la translation va \emph{du premier plan vers le second}. Au Bac, inverser l'ordre est l'erreur la plus fréquente sur ce thème.
\end{attention}

\courssub{Composée de deux réflexions de plans perpendiculaires}

Supposons maintenant que $P$ et $Q$ sont \textbf{perpendiculaires}, sécants selon une droite $(d)$. Que vaut $s_Q \circ s_P$ ?

L'intuition est la suivante : un point de la droite $(d)$ appartient aux deux plans, il est donc invariant par chacune des deux réflexions, donc par la composée. La droite $(d)$ est ainsi fixée \emph{point par point}. Pour un point $M$ hors de $(d)$, les deux symétries successives le « font pivoter » d'un demi-tour autour de $(d)$. Formalisons cela.

\begin{propriete}[Théorème 2 : composée de deux réflexions de plans perpendiculaires]
Soient $P$ et $Q$ deux plans \textbf{perpendiculaires}, d'intersection la droite $(d)$. Alors :
\[ s_Q \circ s_P = s_{(d)} \]
c'est-à-dire le \cle{demi-tour} (symétrie orthogonale) d'axe $(d)$.
\par\textbf{Démonstration guidée ci-dessous (démarche exigible).}
\end{propriete}

\begin{experience}[Démonstration guidée du Théorème 2]
Notons $f = s_Q \circ s_P$. Procédons en trois étapes.
\begin{enumerate}
\item \textbf{Les points de $(d)$ sont invariants.} Soit $A \in (d)$. Alors $A \in P$, donc $s_P(A) = A$ ; et $A \in Q$, donc $s_Q(A) = A$. Par suite $f(A) = s_Q(s_P(A)) = A$ : la droite $(d)$ est invariante \emph{point par point} par $f$.
\item \textbf{Construction géométrique.} Soit $M \notin (d)$, $M' = s_P(M)$ et $M'' = s_Q(M') = f(M)$. Soit $H$ le projeté orthogonal de $M$ sur $(d)$. Les réflexions conservent les distances et les angles droits ; de plus $(d)$ est invariante point par point par $s_P$ et $s_Q$. Donc $H$, projeté de $M$ sur $(d)$, est aussi le projeté orthogonal de $M'$ sur $(d)$ (car $s_P$ conserve l'orthogonalité et fixe $(d)$), puis celui de $M''$ sur $(d)$. Les trois points $M$, $M'$, $M''$ ont donc le \emph{même} projeté orthogonal $H$ sur $(d)$, et $HM = HM' = HM''$.
\item \textbf{$H$ est le milieu de $[MM'']$.} Les points $M$, $M'$, $M''$ et $H$ sont coplanaires, dans le plan $(\Pi)$ passant par $M$ et perpendiculaire à $(d)$. Dans ce plan, les traces de $P$ et $Q$ sont deux droites \emph{perpendiculaires} passant par $H$ (car $P \perp Q$), et la situation plane est celle de la composée de deux symétries axiales d'axes perpendiculaires : c'est la symétrie centrale de centre $H$. Donc $H$ est le milieu de $[MM'']$, c'est-à-dire $\overrightarrow{HM''} = -\overrightarrow{HM}$.
\end{enumerate}
Conclusion : pour tout point $M$, $f(M)$ est le point $M''$ tel que le projeté orthogonal $H$ de $M$ sur $(d)$ soit le milieu de $[MM'']$ : c'est exactement la définition du \textbf{demi-tour d'axe $(d)$}. \hfill$\blacksquare$
\end{experience}

\begin{popfigure}
\begin{center}
\begin{tikzpicture}[scale=1.2, >=Stealth]
% Traces of P and Q in horizontal plane (perpendicular to d)
% P: x=0 -> trace is y-axis (vertical). Q: y=0 -> trace is x-axis (horizontal).
\draw[popblue, thick] (0,-2.5) -- (0,2.5) node[above] {\footnotesize Trace de $P$};
\draw[poporange, thick] (-2.5,0) -- (2.5,0) node[right] {\footnotesize Trace de $Q$};
% Right angle mark
\draw[popdark] (0.2,0) -- (0.2,0.2) -- (0,0.2);
% Point H (intersection of d with horizontal plane)
\filldraw[popdark] (0,0) circle (1.5pt) node[below left] {$H$};
% Points M, M', M''
\coordinate (M) at (-1.5,-1.0);
\coordinate (Mp) at (1.5,-1.0);
\coordinate (Mpp) at (1.5,1.0);
\filldraw[popdark] (M) circle (1.8pt) node[below left] {$M$};
\filldraw[popgreen] (Mp) circle (1.8pt) node[below right] {$M'$};
\filldraw[poppink] (Mpp) circle (1.8pt) node[above right] {$M''$};
% Segments for reflections
\draw[popline, dashed] (M) -- (Mp) node[midway, below] {\footnotesize $s_P$};
\draw[popline, dashed] (Mp) -- (Mpp) node[midway, right] {\footnotesize $s_Q$};
% Half-turn arrow
\draw[poppurple, very thick, <->] (M) -- (Mpp) node[midway, above left, sloped] {\footnotesize Demi-tour};
% Projection lines (optional clarity)
\draw[popblue, dashed, thin] (M) -- (-1.5,0);
\draw[popblue, dashed, thin] (Mp) -- (1.5,0);
\draw[poporange, dashed, thin] (Mp) -- (0,-1.0);
\draw[poporange, dashed, thin] (Mpp) -- (0,1.0);
\end{tikzpicture}
\end{center}
\legende{Vue de dessus dans un plan perpendiculaire à $(d)$ : $H$ est le milieu de $[MM'']$, la composée des deux réflexions axiales (traces de $P$ et $Q$) est la symétrie centrale de centre $H$, c'est-à-dire le demi-tour d'axe $(d)$.}
\end{popfigure}

\begin{exemplebox}[Exemple chiffré]
Soient $P : x = 0$ (plan $(Oyz)$) et $Q : y = 0$ (plan $(Oxz)$). Ces plans sont perpendiculaires (leurs normales $\vec{i}$ et $\vec{j}$ sont orthogonales) et sécants selon la droite $(d) : x = 0 \text{ et } y = 0$, c'est-à-dire l'axe $(Oz)$. Vérifions analytiquement : pour $M(x\,;y\,;z)$,
\[ M \xrightarrow{\;s_P\;} M'(-x\,;y\,;z) \xrightarrow{\;s_Q\;} M''(-x\,;-y\,;z). \]
Le projeté orthogonal de $M$ sur $(Oz)$ est $H(0\,;0\,;z)$, et l'on a bien $\overrightarrow{HM''}(-x\,;-y\,;0) = -\overrightarrow{HM}(x\,;y\,;0)$ : $H$ est le milieu de $[MM'']$. La composée $s_Q \circ s_P$ est donc le demi-tour d'axe $(Oz)$, d'expression analytique $(x\,;y\,;z) \mapsto (-x\,;-y\,;z)$. $\checkmark$
\end{exemplebox}

\begin{attention}[Plans perpendiculaires : ici, l'ordre ne change pas la nature]
Si $P \perp Q$, alors $s_Q \circ s_P$ et $s_P \circ s_Q$ sont \emph{toutes deux} le demi-tour d'axe $(d) = P \cap Q$ : on a même $s_Q \circ s_P = s_P \circ s_Q$ dans ce cas particulier. Mais attention : cette commutation est une exception ! Ne la généralise jamais aux autres situations (voir le cas des plans parallèles ci-dessus).
\end{attention}

\courssub{Réciproques : décomposer une translation ou un demi-tour en deux réflexions}

Les théorèmes 1 et 2 admettent des réciproques très utiles pour résoudre les problèmes du Bac : au lieu de partir des réflexions pour fabriquer une isométrie, on part de l'isométrie et on la \emph{décompose}.

\begin{propriete}[Réciproques]
\begin{itemize}
\item \textbf{Toute translation} de vecteur $\vec{u}$ est la composée de deux réflexions de plans strictement parallèles, perpendiculaires à $\vec{u}$, distants de $\dfrac{\|\vec{u}\|}{2}$. Précisément, si $P$ est un plan quelconque perpendiculaire à $\vec{u}$ et $Q = t_{\vec{u}/2}(P)$, alors $t_{\vec{u}} = s_Q \circ s_P$. Il y a donc une \emph{infinité} de telles décompositions.
\item \textbf{Tout demi-tour} d'axe $(d)$ est la composée de deux réflexions de plans perpendiculaires se coupant selon $(d)$. Si $P$ est un plan quelconque contenant $(d)$ et $Q$ le plan contenant $(d)$ et perpendiculaire à $P$, alors $s_{(d)} = s_Q \circ s_P$.
\end{itemize}
\end{propriete}

\begin{exemplebox}[Décomposition d'une translation]
Soit $t$ la translation de vecteur $\vec{u}(0\,;0\,;10)$. Pour écrire $t = s_Q \circ s_P$, il suffit de choisir deux plans horizontaux distants de $\frac{10}{2} = 5$, par exemple $P : z = 1$ et $Q : z = 6$. Alors $2\overrightarrow{H_1H_2} = 2 \times 5\,\vec{k} = 10\vec{k} = \vec{u}$, donc $t = s_Q \circ s_P$. On aurait pu choisir aussi bien $P : z = -2$ et $Q : z = 3$ : la décomposition n'est pas unique.
\end{exemplebox}

\begin{savaistu}[Complément Bac : plans sécants non perpendiculaires (résultat admis)]
Que se passe-t-il si $P$ et $Q$ sont sécants selon $(d)$ avec un angle $\theta$ quelconque ($0 < \theta < \frac{\pi}{2}$) ? Le résultat, \textbf{admis} au programme de Terminale, est que $s_Q \circ s_P$ est la \cle{rotation} d'axe $(d)$ et d'angle $2\theta$. Le cas perpendiculaire correspond à $\theta = \frac{\pi}{2}$, donc $2\theta = \pi$ : on retrouve exactement le demi-tour, qui est la rotation d'angle $\pi$. Ce résultat général unifie toute la section : les réflexions sont les « briques élémentaires » à partir desquelles on construit translations, demi-tours et rotations.
\end{savaistu}

\courssub{Méthode générale et exercice résolu}

\begin{methode}[Déterminer la nature et les éléments caractéristiques de $s_Q \circ s_P$]
\begin{enumerate}
\item \textbf{Étudier la position relative de $P$ et $Q$} : comparer leurs vecteurs normaux $\vec{n}_P$ et $\vec{n}_Q$.
\item \textbf{Si $\vec{n}_P$ et $\vec{n}_Q$ sont colinéaires} (plans strictement parallèles) : $s_Q \circ s_P$ est une \textbf{translation}. Calculer $\overrightarrow{H_1H_2}$ en projetant orthogonalement un point de $P$ sur $Q$, puis conclure : vecteur $2\overrightarrow{H_1H_2}$ (sens : de $P$ vers $Q$).
\item \textbf{Si $\vec{n}_P \cdot \vec{n}_Q = 0$ sans colinéarité} (plans perpendiculaires) : $s_Q \circ s_P$ est le \textbf{demi-tour} d'axe $(d) = P \cap Q$. Déterminer $(d)$ en résolvant le système formé par les deux équations de plans.
\item \textbf{Sinon} (plans sécants non perpendiculaires) : citer le résultat admis — rotation d'axe $P \cap Q$.
\item \textbf{Vérifier} éventuellement par le calcul analytique sur un point $M(x\,;y\,;z)$ générique.
\end{enumerate}
\end{methode}

\begin{exoresolu}[Exercice type Bac]
L'espace est muni d'un repère orthonormé $(O\,;\vec{i},\vec{j},\vec{k})$. On considère les plans $P : x - 2 = 0$ et $Q : x + 1 = 0$, puis les plans $R : x = 0$ et $S : z = 0$.
\begin{enumerate}
\item Déterminer la nature et les éléments caractéristiques de $f = s_Q \circ s_P$.
\item Déterminer la nature et les éléments caractéristiques de $g = s_S \circ s_R$, et donner son expression analytique.
\end{enumerate}
\tcblower
\textbf{1.} Les plans $P : x = 2$ et $Q : x = -1$ admettent tous deux $\vec{i}(1\,;0\,;0)$ pour vecteur normal : ils sont \textbf{strictement parallèles}. D'après le Théorème 1, $f = s_Q \circ s_P$ est une translation. Prenons $H_1(2\,;0\,;0) \in P$ ; son projeté orthogonal sur $Q$ est $H_2(-1\,;0\,;0)$, d'où $\overrightarrow{H_1H_2}(-3\,;0\,;0)$ et :
\[ f = t_{\,\vec{u}} \quad \text{avec} \quad \vec{u} = 2\overrightarrow{H_1H_2} = (-6\,;0\,;0). \]
Vérification analytique : $M(x\,;y\,;z) \xrightarrow{s_P} M'(4-x\,;y\,;z) \xrightarrow{s_Q} M''(-2-(4-x)\,;y\,;z) = (x-6\,;y\,;z)$. On retrouve bien $\overrightarrow{MM''}(-6\,;0\,;0)$. $\checkmark$

\textbf{2.} $R : x = 0$ a pour normal $\vec{i}$ et $S : z = 0$ a pour normal $\vec{k}$. Comme $\vec{i} \cdot \vec{k} = 0$, les plans $R$ et $S$ sont \textbf{perpendiculaires}. Leur intersection est la droite $(d) : \begin{cases} x = 0 \\ z = 0 \end{cases}$, c'est-à-dire l'axe $(Oy)$. D'après le Théorème 2, $g = s_S \circ s_R$ est le \textbf{demi-tour d'axe $(Oy)$}. Expression analytique :
\[ M(x\,;y\,;z) \xrightarrow{s_R} M'(-x\,;y\,;z) \xrightarrow{s_S} M''(-x\,;y\,;-z), \]
donc $g(x\,;y\,;z) = (-x\,;y\,;-z)$. Le projeté orthogonal de $M$ sur $(Oy)$ est $H(0\,;y\,;0)$, milieu de $[MM'']$ : la vérification est immédiate. $\checkmark$
\end{exoresolu}

\begin{exercice}[Pour t'entraîner]
Soient $P : y = 1$ et $Q : y = 4$, puis $P' : x + z = 0$ et $Q' : x - z = 0$.
\begin{enumerate}
\item Déterminer la nature et le vecteur de $s_Q \circ s_P$, puis celui de $s_P \circ s_Q$. Que constates-tu ?
\item Montrer que $P'$ et $Q'$ sont perpendiculaires, déterminer leur droite d'intersection $(d)$, puis donner la nature et l'expression analytique de $s_{Q'} \circ s_{P'}$.
\end{enumerate}
\end{exercice}

\begin{corrige}[Exercice]
\textbf{1.} $P \parallel Q$ (normale commune $\vec{j}$). Avec $H_1(0\,;1\,;0) \in P$ et $H_2(0\,;4\,;0) \in Q$ : $s_Q \circ s_P = t_{\vec{u}}$ avec $\vec{u} = 2\overrightarrow{H_1H_2} = (0\,;6\,;0)$. En inversant l'ordre : $s_P \circ s_Q = t_{-\vec{u}}$, de vecteur $(0\,;-6\,;0)$. Les deux composées sont des translations de vecteurs \textbf{opposés}.

\textbf{2.} $\vec{n}_{P'}(1\,;0\,;1)$ et $\vec{n}_{Q'}(1\,;0\,;-1)$ ; leur produit scalaire vaut $1 - 1 = 0$ : les plans sont perpendiculaires. Intersection : $\begin{cases} x + z = 0 \\ x - z = 0 \end{cases}$ donne $x = z = 0$ : c'est l'axe $(Oy)$. Donc $s_{Q'} \circ s_{P'}$ est le demi-tour d'axe $(Oy)$ : $(x\,;y\,;z) \mapsto (-x\,;y\,;-z)$.
\end{corrige}

\begin{aretenir}[L'essentiel de la section]
\begin{itemize}
\item $P \parallel Q$ strictement $\Rightarrow$ $s_Q \circ s_P = t_{\,2\overrightarrow{H_1H_2}}$, vecteur dirigé \emph{de $P$ vers $Q$}, perpendiculaire aux plans. L'ordre compte : inverser $P$ et $Q$ oppose le vecteur.
\item $P \perp Q$, $P \cap Q = (d)$ $\Rightarrow$ $s_Q \circ s_P$ est le \textbf{demi-tour d'axe $(d)$}.
\item Réciproquement : toute translation se décompose en deux réflexions de plans parallèles ; tout demi-tour se décompose en deux réflexions de plans perpendiculaires.
\item Cas admis : plans sécants d'angle $\theta$ $\Rightarrow$ rotation d'axe $P \cap Q$ et d'angle $2\theta$.
\item La première chose à faire face à une composée de réflexions : \textbf{étudier la position relative des plans}.
\end{itemize}
\end{aretenir}

\coursec{Demi-tour : définition et expression analytique}

Dans la section précédente, nous avons composé deux réflexions de plans perpendiculaires $(P)$ et $(Q)$, et nous avons découvert que la composée $s_Q \circ s_P$ est une nouvelle isométrie : le \cle{demi-tour} d'axe $(d) = (P) \cap (Q)$. Cette transformation mérite d'être étudiée pour elle-même. Elle possède une définition géométrique simple, une expression analytique maniable, et elle apparaît partout dans la vie quotidienne : une porte que l'on ouvre à $180$ degrés, une hélice qui tourne d'un demi-tour, un panneau publicitaire rotatif sur le boulevard de la Liberté à Douala. Cette section lui est entièrement consacrée.

\courssub{Définition du demi-tour}

Commençons par l'intuition. Prends une droite $(d)$ dans l'espace, par exemple l'arête verticale d'une porte (l'axe des charnières). Un point $M$ de la porte, lorsqu'on la fait pivoter d'un demi-tour autour de cet axe, décrit un demi-cercle dans un plan perpendiculaire à $(d)$, et arrive en un point $M'$ situé de l'autre côté de l'axe, à la même distance. Le point de l'axe le plus proche de $M$ — son projeté orthogonal $H$ sur $(d)$ — est alors exactement le milieu du segment $[MM']$. C'est cette observation qui fonde la définition.

\begin{definition}[Demi-tour d'axe $(d)$]
Soit $(d)$ une droite de l'espace. On appelle \cle{demi-tour} d'axe $(d)$, noté $s_d$, la transformation de l'espace qui, à tout point $M$, associe le point $M'$ défini par :
\begin{itemize}
\item si $M \in (d)$, alors $M' = M$ ;
\item si $M \notin (d)$, alors $M'$ est le point tel que le projeté orthogonal $H$ de $M$ sur $(d)$ soit le milieu du segment $[MM']$.
\end{itemize}
Le demi-tour $s_d$ est aussi appelé \cle{symétrie orthogonale par rapport à la droite} $(d)$, ou encore \emph{symétrie d'axe} $(d)$.
\end{definition}

\begin{attention}[Demi-tour ou symétrie centrale ?]
Ne confonds pas le demi-tour d'axe $(d)$ avec la symétrie centrale. Dans le demi-tour d'axe $(d)$, c'est le \textbf{projeté orthogonal} $H$ de $M$ sur $(d)$ qui est le milieu de $[MM']$, et $H$ dépend de $M$. Dans une symétrie centrale de centre $\Omega$, c'est le \textbf{même point fixe} $\Omega$ qui est le milieu de tous les segments $[MM']$. Le demi-tour possède toute une droite de points fixes ; la symétrie centrale n'en possède qu'un seul.
\end{attention}

\begin{propriete}[Lien avec la rotation]
Le demi-tour d'axe $(d)$ coïncide avec la rotation d'axe $(d)$ et d'angle $\pi$ : faire tourner un point d'un demi-tour ($180$ degrés) autour de $(d)$ le conduit exactement au point $M'$ de la définition. En effet, dans le plan perpendiculaire à $(d)$ passant par $M$, la rotation de centre $H$ et d'angle $\pi$ est la symétrie centrale de centre $H$, qui envoie $M$ sur le point $M'$ tel que $H$ soit le milieu de $[MM']$.
\end{propriete}

\begin{popfigure}
\begin{tikzpicture}[scale=1.15, line cap=round, line join=round]
% Axe vertical (d)
\draw[popblue, very thick] (0,-1.2) -- (0,3.4);
\node[popblue, right] at (0,3.4) {$(d)$};
% Plan perpendiculaire (suggéré)
\draw[popmuted, dashed] (-2.6,1.1) -- (2.6,1.1);
\draw[popmuted, dashed] (-2.6,1.1) -- (-1.4,0.3);
\draw[popmuted, dashed] (2.6,1.1) -- (1.4,0.3);
\draw[popmuted, dashed] (-1.4,0.3) -- (1.4,0.3);
\node[popmuted, font=\small] at (2.9,0.35) {plan $\perp (d)$};
% Points M, H, M'
\coordinate (H) at (0,1.1);
\coordinate (M) at (2.1,1.55);
\coordinate (Mp) at (-2.1,0.65);
% Demi-cercle de rotation
\draw[poporange, thick, ->] (2.1,1.42) arc[start angle=8, end angle=172, x radius=2.1, y radius=0.5];
\node[poporange, font=\small] at (0.2,2.05) {$180^\circ$};
% Segment MM'
\draw[popdark, thick] (M) -- (Mp);
% Angle droit en H
\draw[popdark] (0.16,1.1) -- (0.16,1.26) -- (0,1.26);
% Points
\fill[popink] (M) circle (1.8pt); \node[popink, above right] at (M) {$M$};
\fill[poppink] (Mp) circle (1.8pt); \node[poppink, below left] at (Mp) {$M' = s_d(M)$};
\fill[popblue] (H) circle (1.8pt); \node[popblue, left] at (H) {$H$};
% Milieu
\node[popgreen, font=\small, below right=2pt and -6pt of H] {$H$ milieu de $[MM']$};
\end{tikzpicture}
\legende{Demi-tour d'axe $(d)$ : le projeté orthogonal $H$ de $M$ sur $(d)$ est le milieu de $[MM']$. Le point $M$ décrit un demi-cercle dans le plan perpendiculaire à $(d)$ passant par $M$.}
\end{popfigure}

\begin{savaistu}[Une porte est un demi-tour en action]
Lorsque tu pousses une porte jusqu'à ce qu'elle soit complètement ouverte contre le mur, chaque point de la porte a effectué une rotation autour de la droite verticale des gonds (charnières). Si l'angle d'ouverture est exactement $180$ degrés, la position finale de la porte est l'image de sa position initiale par le \textbf{demi-tour} dont l'axe est la droite des charnières. Les poignées de porte, les ventilateurs de plafond, les éoliennes du nord du Cameroun : toutes ces rotations familières deviennent des demi-tours dès que l'angle atteint $\pi$.
\end{savaistu}

\courssub{Premières propriétés}

\begin{propriete}[Propriétés fondamentales du demi-tour]
Soit $(d)$ une droite de l'espace et $s_d$ le demi-tour d'axe $(d)$.
\begin{enumerate}
\item $s_d$ est une \cle{involution} : $s_d \circ s_d = \mathrm{Id}$, c'est-à-dire que $s_d$ est sa propre réciproque.
\item $s_d$ est une \cle{isométrie} : elle conserve les distances, c'est-à-dire que pour tous points $A$ et $B$, on a $s_d(A)\,s_d(B) = AB$.
\item L'ensemble des points fixes de $s_d$ est \textbf{exactement} la droite $(d)$ : tout point de $(d)$ est invariant, et aucun autre point ne l'est.
\item $s_d$ conserve l'alignement, le parallélisme, l'orthogonalité et les milieux (conséquences du fait que c'est une isométrie affine).
\end{enumerate}
\end{propriete}

Justifions rapidement ces affirmations. Pour le point 1 : si $M' = s_d(M)$, alors $H$ est le milieu de $[MM']$, donc aussi de $[M'M]$, ce qui signifie que $s_d(M') = M$. Pour le point 3 : si $M \notin (d)$, alors $H \neq M$, donc le point $M'$ tel que $H$ soit le milieu de $[MM']$ vérifie $M' \neq M$ : le point $M$ n'est pas fixe. Pour le point 2, on peut remarquer que $s_d$ est la composée de deux réflexions de plans perpendiculaires se coupant selon $(d)$ (section précédente), et qu'une composée d'isométries est une isométrie.

\courssub{Expression analytique : les cas fondamentaux}

Plaçons-nous dans un repère orthonormé $(O\,;\vec{i},\vec{j},\vec{k})$ de l'espace. Les demi-tours d'axes $(Ox)$, $(Oy)$, $(Oz)$ ont des expressions analytiques remarquablement simples, qu'il faut connaître par cœur.

Prenons l'axe $(Oz)$. Soit $M(x\,;y\,;z)$. Le projeté orthogonal de $M$ sur $(Oz)$ est le point $H(0\,;0\,;z)$ : on conserve la cote et on annule les deux autres coordonnées. Dire que $H$ est le milieu de $[MM']$ avec $M'(x'\,;y'\,;z')$ donne :
\[
\frac{x+x'}{2}=0,\qquad \frac{y+y'}{2}=0,\qquad \frac{z+z'}{2}=z,
\]
d'où $x'=-x$, $y'=-y$, $z'=z$. On retient :

\begin{aretenir}[Demi-tours des axes de coordonnées]
Dans un repère orthonormé $(O\,;\vec{i},\vec{j},\vec{k})$ :
\begin{itemize}
\item le demi-tour d'axe $(Oz)$ : $(x\,;y\,;z) \longmapsto (-x\,;-y\,;z)$ ;
\item le demi-tour d'axe $(Ox)$ : $(x\,;y\,;z) \longmapsto (x\,;-y\,;-z)$ ;
\item le demi-tour d'axe $(Oy)$ : $(x\,;y\,;z) \longmapsto (-x\,;y\,;-z)$.
\end{itemize}
Règle mnémotechnique : la coordonnée correspondant à l'axe est \textbf{conservée}, les deux autres sont \textbf{changées en leur opposé}.
\end{aretenir}

\begin{attention}[Le piège classique du Bac]
L'erreur la plus fréquente consiste à écrire que le demi-tour d'axe $(Oz)$ envoie $(x\,;y\,;z)$ sur $(-x\,;-y\,;-z)$. C'est \textbf{faux} ! La transformation $(x\,;y\,;z) \mapsto (-x\,;-y\,;-z)$ est la \cle{symétrie centrale} de centre $O$, qui n'a qu'un seul point fixe. Le demi-tour d'axe $(Oz)$ conserve la cote $z$ : tous les points de $(Oz)$ doivent rester fixes, ce qui n'est possible que si $z' = z$. Vérifie toujours ta réponse sur un point de l'axe : le point $(0\,;0\,;5)$ doit être invariant.
\end{attention}

\begin{popfigure}
\begin{tikzpicture}[scale=1.5, line cap=round, line join=round]
% Repère
\draw[->, popmuted] (0,0) -- (3.2,0) node[below left, font=\small]{$x$};
\draw[->, popmuted] (0,0) -- (-1.5,-1.1) node[below, font=\small]{$y$};
\draw[->, popmuted] (0,0) -- (0,3.1) node[right, font=\small]{$z$};
\fill (0,0) circle (1pt) node[below right, font=\small]{$O$};
% Axe (Oz) en gras
\draw[popblue, very thick] (0,-0.3) -- (0,3.1);
% Triangle ABC à z = 2
\coordinate (A) at (1.6,2.0);
\coordinate (B) at (2.5,1.55);
\coordinate (C) at (2.1,2.6);
% Images par demi-tour d'axe (Oz)
\coordinate (Ap) at (-1.6,2.0);
\coordinate (Bp) at (-2.5,1.55);
\coordinate (Cp) at (-2.1,2.6);
% Triangles
\draw[popink, thick, fill=popink!12] (A) -- (B) -- (C) -- cycle;
\draw[poppink, thick, fill=poppink!12] (Ap) -- (Bp) -- (Cp) -- cycle;
% Liaisons pointillées vers l'axe
\draw[popmuted, dashed] (A) -- (0,2.0) -- (Ap);
\draw[popmuted, dashed] (C) -- (0,2.6) -- (Cp);
\fill[popblue] (0,2.0) circle (1.2pt);
\fill[popblue] (0,2.6) circle (1.2pt);
% Labels
\node[popink, font=\small, right] at (A) {$A(1{,}6\,;0\,;2)$};
\node[popink, font=\small, right] at (B) {$B$};
\node[popink, font=\small, above] at (C) {$C$};
\node[poppink, font=\small, left] at (Ap) {$A'(-1{,}6\,;0\,;2)$};
\node[poppink, font=\small, left] at (Bp) {$B'$};
\node[poppink, font=\small, above] at (Cp) {$C'$};
\end{tikzpicture}
\legende{Un triangle $ABC$ et son image $A'B'C'$ par le demi-tour d'axe $(Oz)$ : les cotes sont conservées, les abscisses et ordonnées changent de signe. Les segments joignant un point à son image passent par l'axe $(Oz)$ et y sont perpendiculaires.}
\end{popfigure}

\courssub{Expression analytique : le cas général}

Passons maintenant au cas général : l'axe $(d)$ est une droite quelconque, passant par un point $A$ et dirigée par un vecteur $\vec{u}$ que l'on choisira \cle{unitaire} (c'est-à-dire $\|\vec{u}\| = 1$). Soit $M$ un point de l'espace, $H$ son projeté orthogonal sur $(d)$, et $M' = s_d(M)$.

La condition « $H$ est le milieu de $[MM']$ » se traduit vectoriellement par :
\[
\overrightarrow{HM} + \overrightarrow{HM'} = \vec{0}
\quad\Longleftrightarrow\quad
\overrightarrow{OM'} = 2\,\overrightarrow{OH} - \overrightarrow{OM},
\]
c'est-à-dire, en notant encore $M$, $H$, $M'$ les vecteurs positions :
\[
\boxed{\;M' = 2H - M\;}
\]
Tout le travail consiste donc à calculer $H$. Or $H$ est le projeté orthogonal de $M$ sur la droite $(A\,;\vec{u})$ : le vecteur $\overrightarrow{AH}$ est la projection orthogonale de $\overrightarrow{AM}$ sur $\vec{u}$. Comme $\vec{u}$ est unitaire, cette projection s'écrit $\left(\overrightarrow{AM}\cdot\vec{u}\right)\vec{u}$. D'où la formule complète.

\begin{propriete}[Expression vectorielle du demi-tour]
Soit $(d)$ la droite passant par $A$ et dirigée par le vecteur \textbf{unitaire} $\vec{u}$. Pour tout point $M$, son image $M'$ par le demi-tour d'axe $(d)$ est :
\[
M' = 2H - M
\qquad\text{avec}\qquad
H = A + \left(\overrightarrow{AM}\cdot\vec{u}\right)\vec{u}.
\]
Si $\vec{u}$ n'est pas unitaire, on remplace le produit scalaire par $\dfrac{\overrightarrow{AM}\cdot\vec{u}}{\|\vec{u}\|^{2}}\,\vec{u}$.
\end{propriete}

\begin{experience}[Démonstration guidée de la formule $M' = 2H - M$]
Nous allons établir rigoureusement la formule, étape par étape.
\begin{enumerate}
\item \textbf{Position du problème.} Soit $(d)$ une droite passant par $A$, dirigée par $\vec{u}$ unitaire, $M$ un point quelconque, $H$ son projeté orthogonal sur $(d)$ et $M' = s_d(M)$.
\item \textbf{Traduction de la définition.} Par définition du demi-tour, $H$ est le milieu de $[MM']$. La caractérisation vectorielle du milieu donne $\overrightarrow{OH} = \dfrac{1}{2}\left(\overrightarrow{OM} + \overrightarrow{OM'}\right)$, d'où $\overrightarrow{OM'} = 2\,\overrightarrow{OH} - \overrightarrow{OM}$, soit $M' = 2H - M$.
\item \textbf{Calcul de $H$.} Puisque $H \in (d)$, il existe un réel $t$ tel que $\overrightarrow{AH} = t\,\vec{u}$. La condition d'orthogonalité $\overrightarrow{HM} \perp \vec{u}$ s'écrit $\overrightarrow{HM}\cdot\vec{u} = 0$. Or $\overrightarrow{HM} = \overrightarrow{AM} - \overrightarrow{AH} = \overrightarrow{AM} - t\,\vec{u}$, donc :
\[
\left(\overrightarrow{AM} - t\,\vec{u}\right)\cdot\vec{u} = 0
\;\Longleftrightarrow\;
\overrightarrow{AM}\cdot\vec{u} = t\,\|\vec{u}\|^{2} = t
\quad(\text{car } \|\vec{u}\| = 1).
\]
Ainsi $\overrightarrow{AH} = \left(\overrightarrow{AM}\cdot\vec{u}\right)\vec{u}$, c'est-à-dire $H = A + \left(\overrightarrow{AM}\cdot\vec{u}\right)\vec{u}$.
\item \textbf{Conclusion.} En reportant dans l'étape 2 : $M' = 2A + 2\left(\overrightarrow{AM}\cdot\vec{u}\right)\vec{u} - M$. La formule est démontrée. $\blacksquare$
\end{enumerate}
\end{experience}

\begin{methode}[Calculer l'image d'un point par un demi-tour]
Pour déterminer l'image $M'$ d'un point $M$ par le demi-tour d'axe $(d)$ passant par $A$ et dirigé par $\vec{u}$ :
\begin{enumerate}
\item \textbf{Rendre $\vec{u}$ unitaire} si nécessaire, en le divisant par sa norme (ou garder la formule avec $\|\vec{u}\|^2$ au dénominateur).
\item \textbf{Calculer le produit scalaire} $\overrightarrow{AM}\cdot\vec{u}$.
\item \textbf{Obtenir le projeté} : $H = A + \left(\overrightarrow{AM}\cdot\vec{u}\right)\vec{u}$.
\item \textbf{Conclure} par $M' = 2H - M$.
\item \textbf{Vérifier} : le milieu de $[MM']$ doit appartenir à $(d)$ et $\overrightarrow{MM'}$ doit être orthogonal à $\vec{u}$.
\end{enumerate}
\end{methode}

\begin{exoresolu}[Image d'un point par un demi-tour]
L'espace est muni d'un repère orthonormé $(O\,;\vec{i},\vec{j},\vec{k})$. On considère le demi-tour $s_d$ d'axe $(d)$ passant par $A(1\,;0\,;2)$ et dirigé par $\vec{u}(0\,;0\,;1)$. Déterminer l'image $B'$ du point $B(3\,;1\,;5)$.
\tcblower
\textbf{Solution détaillée.}

\textbf{Étape 1.} Le vecteur $\vec{u}(0\,;0\,;1)$ est déjà unitaire : $\|\vec{u}\| = \sqrt{0^{2}+0^{2}+1^{2}} = 1$.

\textbf{Étape 2.} Calculons $\overrightarrow{AB} = (3-1\,;\,1-0\,;\,5-2) = (2\,;1\,;3)$, puis le produit scalaire :
\[
\overrightarrow{AB}\cdot\vec{u} = 2\times 0 + 1\times 0 + 3\times 1 = 3.
\]

\textbf{Étape 3.} Le projeté orthogonal $H$ de $B$ sur $(d)$ est :
\[
H = A + 3\,\vec{u} = (1\,;0\,;2) + (0\,;0\,;3) = (1\,;0\,;5).
\]

\textbf{Étape 4.} On applique $B' = 2H - B$ :
\begin{align*}
x_{B'} &= 2\times 1 - 3 = -1,\\
y_{B'} &= 2\times 0 - 1 = -1,\\
z_{B'} &= 2\times 5 - 5 = 5.
\end{align*}
Donc $\boxed{B'(-1\,;-1\,;5)}$.

\textbf{Étape 5 (vérification).} Le milieu de $[BB']$ est $\left(\dfrac{3-1}{2}\,;\dfrac{1-1}{2}\,;\dfrac{5+5}{2}\right) = (1\,;0\,;5) = H$, qui appartient bien à $(d)$. De plus $\overrightarrow{BB'} = (-4\,;-2\,;0)$ et $\overrightarrow{BB'}\cdot\vec{u} = 0$ : le segment $[BB']$ est bien perpendiculaire à l'axe. La réponse est cohérente.

\emph{Remarque :} l'axe $(d)$ est la droite verticale passant par $A$, parallèle à $(Oz)$. On retrouve le schéma du cas fondamental : la cote est conservée ($z' = 5$), seules les coordonnées horizontales sont « repliées » autour de l'axe.
\end{exoresolu}

\begin{exercice}[À toi de jouer]
Dans un repère orthonormé, on considère le demi-tour $s_d$ d'axe $(d)$ passant par $A(0\,;1\,;0)$ et dirigé par $\vec{u}(1\,;0\,;0)$.
\begin{enumerate}
\item Déterminer l'image du point $M(4\,;3\,;-2)$.
\item Le point $N(2\,;1\,;7)$ est-il fixe par $s_d$ ? Justifier.
\item Déterminer l'expression analytique de $s_d$, c'est-à-dire les coordonnées $(x'\,;y'\,;z')$ de l'image de $M(x\,;y\,;z)$ en fonction de $x$, $y$, $z$.
\end{enumerate}
\end{exercice}

\begin{corrige}[Exercice : demi-tour d'axe horizontal]
\textbf{1.} $\vec{u}$ est unitaire. $\overrightarrow{AM} = (4\,;2\,;-2)$ et $\overrightarrow{AM}\cdot\vec{u} = 4$. Donc $H = A + 4\vec{u} = (4\,;1\,;0)$, puis $M' = 2H - M = (8-4\,;\,2-3\,;\,0-(-2)) = (4\,;-1\,;2)$.

\textbf{2.} $N$ est fixe si et seulement si $N \in (d)$. Or $(d)$ est l'ensemble des points $(t\,;1\,;0)$, $t \in \mathbb{R}$. La cote de $N$ est $7 \neq 0$, donc $N \notin (d)$ : $N$ n'est \textbf{pas} fixe.

\textbf{3.} Pour $M(x\,;y\,;z)$ : $\overrightarrow{AM} = (x\,;y-1\,;z)$, $\overrightarrow{AM}\cdot\vec{u} = x$, donc $H = (x\,;1\,;0)$ et $M' = 2H - M$, d'où :
\[
\begin{cases}
x' = x,\\
y' = 2 - y,\\
z' = -z.
\end{cases}
\]
On reconnaît la structure du cas fondamental : la coordonnée le long de l'axe ($x$) est conservée ; les deux autres sont symétrisées par rapport à l'axe ($y$ par rapport à $1$, cote de l'axe ; $z$ par rapport à $0$).
\end{corrige}

\begin{aretenir}[L'essentiel de la section]
\begin{itemize}
\item Le \cle{demi-tour} d'axe $(d)$ est la symétrie orthogonale par rapport à la droite $(d)$ : $M' = s_d(M)$ signifie que le projeté orthogonal $H$ de $M$ sur $(d)$ est le milieu de $[MM']$.
\item C'est une \textbf{involution} et une \textbf{isométrie}, dont l'ensemble des points fixes est exactement $(d)$ ; elle coïncide avec la rotation d'axe $(d)$ et d'angle $\pi$.
\item Formule fondamentale : $M' = 2H - M$ avec $H = A + \left(\overrightarrow{AM}\cdot\vec{u}\right)\vec{u}$ pour $\vec{u}$ unitaire.
\item Cas fondamentaux : axe $(Oz)$ : $(x\,;y\,;z) \mapsto (-x\,;-y\,;z)$ — la coordonnée de l'axe est conservée, les autres changent de signe.
\end{itemize}
\end{aretenir}

Maintenant que le demi-tour est bien en main, une question naturelle se pose, symétrique de celle de la section précédente : que se passe-t-il lorsqu'on compose \emph{deux} demi-tours ? La réponse dépend, comme toujours, de la position relative des deux axes : c'est l'objet de la section suivante.

\coursec{Composée de deux demi-tours}

Dans la section précédente, tu as découvert le \cle{demi-tour}, symétrie orthogonale par rapport à une droite, et tu sais maintenant déterminer son expression analytique. Une question naturelle se pose alors, comme pour les réflexions : que se passe-t-il lorsqu'on enchaîne deux demi-tours ? La réponse dépend entièrement de la \cle{position relative} des deux axes. C'est ce que nous allons établir rigoureusement dans cette section, qui est l'un des points les plus classiques du Baccalauréat C et E.

\courssub{Position relative de deux droites de l'espace}

Avant d'énoncer les théorèmes, rappelons un point essentiel de géométrie de Première : dans l'espace, deux droites ne sont pas forcément sécantes ou parallèles.

\begin{definition}[Positions relatives de deux droites]
Soit $(d_1)$ et $(d_2)$ deux droites de l'espace. Trois cas, et trois seulement, peuvent se présenter :
\begin{enumerate}
\item $(d_1)$ et $(d_2)$ sont \cle{coplanaires} et \cle{parallèles} (strictement parallèles ou confondues) ;
\item $(d_1)$ et $(d_2)$ sont \cle{sécantes} en un point $O$ (elles sont alors coplanaires ; si de plus elles sont perpendiculaires, on dit qu'elles sont \cle{perpendiculaires}) ;
\item $(d_1)$ et $(d_2)$ sont \cle{non coplanaires} : elles ne sont ni parallèles ni sécantes.
\end{enumerate}
\end{definition}

\begin{attention}[Orthogonales ou perpendiculaires ?]
Deux droites sont dites \cle{orthogonales} lorsque leurs directions sont orthogonales : cela ne suppose \textbf{pas} qu'elles se coupent. Elles sont \cle{perpendiculaires} lorsqu'elles sont orthogonales \textbf{et} sécantes. Dans l'espace, deux droites orthogonales peuvent très bien être non coplanaires. Cette distinction sera capitale dans la suite : les théorèmes de cette section exigent des axes \emph{parallèles} ou \emph{perpendiculaires sécants}, et rien d'autre.
\end{attention}

\courssub{Composée de deux demi-tours d'axes parallèles : c'est une translation}

\textbf{Intuition.} Imagine deux poteaux électriques verticaux et parallèles le long de la route Yaoundé--Douala. Fais subir à un point $M$ un demi-tour autour du premier poteau, puis un demi-tour autour du second. En regardant la scène \emph{depuis le ciel} (c'est-à-dire en projection sur le sol, plan perpendiculaire aux poteaux), chaque demi-tour se lit comme une symétrie centrale. Or tu sais depuis la classe de Seconde que la composée de deux symétries centrales est une translation. Le résultat en trois dimensions va suivre exactement la même loi.

\begin{propriete}[Théorème 1 : composée de deux demi-tours d'axes parallèles]
Soit $(d_1)$ et $(d_2)$ deux droites strictement parallèles de l'espace. Soit $H_1 \in (d_1)$ et $H_2 \in (d_2)$ tels que la droite $(H_1H_2)$ soit perpendiculaire commune à $(d_1)$ et $(d_2)$. Alors :
\[ s_{d_2} \circ s_{d_1} = t_{\,2\overrightarrow{H_1H_2}} \]
c'est-à-dire la \cle{translation} de vecteur $2\overrightarrow{H_1H_2}$.
\end{propriete}

\begin{experience}[Démonstration du théorème 1 (exigible)]
On note $f = s_{d_2} \circ s_{d_1}$. Soit $M$ un point quelconque de l'espace, $M' = s_{d_1}(M)$ et $M'' = s_{d_2}(M') = f(M)$.

\textbf{Étape 1 : choisir un plan perpendiculaire aux deux axes.} Soit $(P)$ le plan passant par $M$ et perpendiculaire à $(d_1)$. Comme $(d_1) \parallel (d_2)$, le plan $(P)$ est aussi perpendiculaire à $(d_2)$. Notons $I_1 = (P) \cap (d_1)$ et $I_2 = (P) \cap (d_2)$.

\textbf{Étape 2 : les trois points $M$, $M'$, $M''$ sont dans $(P)$.} Le demi-tour $s_{d_1}$ envoie $M$ sur le point $M'$ tel que le milieu de $[MM']$ soit le projeté orthogonal de $M$ sur $(d_1)$. Or $(P)$ est perpendiculaire à $(d_1)$ et contient $M$ : le projeté orthogonal de $M$ sur $(d_1)$ est donc $I_1$, et $M' \in (P)$ avec $I_1$ milieu de $[MM']$. De même, $I_2$ est le milieu de $[M'M'']$ et $M'' \in (P)$.

\textbf{Étape 3 : lecture plane.} Dans le plan $(P)$, la restriction de $s_{d_1}$ est la symétrie centrale de centre $I_1$, et celle de $s_{d_2}$ est la symétrie centrale de centre $I_2$. Donc, dans $(P)$ :
\[ \overrightarrow{MM''} = \overrightarrow{MM'} + \overrightarrow{M'M''} = 2\overrightarrow{I_1M'} + 2\overrightarrow{M'I_2} = 2\overrightarrow{I_1I_2}. \]

\textbf{Étape 4 : le vecteur ne dépend pas de $M$.} Le vecteur $\overrightarrow{I_1I_2}$ est indépendant du point $M$ choisi : quand $M$ varie, le plan $(P)$ se translate parallèlement à lui-même le long des axes, et le segment $[I_1I_2]$ se déplace sans changer de vecteur. En particulier, en prenant le plan perpendiculaire commun passant par $H_1$, on obtient $\overrightarrow{I_1I_2} = \overrightarrow{H_1H_2}$. Ainsi, pour tout point $M$ :
\[ \overrightarrow{Mf(M)} = 2\overrightarrow{H_1H_2}, \]
ce qui prouve que $f$ est la translation de vecteur $2\overrightarrow{H_1H_2}$. $\blacksquare$
\end{experience}

\begin{popfigure}
\begin{center}
\begin{tikzpicture}[scale=1.05, >=Stealth]
% axes parallèles verticaux
\draw[very thick, popblue] (0,-0.6) -- (0,4.2) node[above, popblue] {$(d_1)$};
\draw[very thick, poporange] (3.4,-0.6) -- (3.4,4.2) node[above, poporange] {$(d_2)$};
% perpendiculaire commune
\draw[popmuted, dashed] (0,3.4) -- (3.4,3.4);
\fill[popblue] (0,3.4) circle (1.6pt) node[above left, popdark] {$H_1$};
\fill[poporange] (3.4,3.4) circle (1.6pt) node[above right, popdark] {$H_2$};
% points M, M', M''
\coordinate (M) at (-1.3,1.1);
\coordinate (Mp) at (1.3,1.1);
\coordinate (Mpp) at (2.1,1.1);
\fill[popink] (M) circle (2pt) node[below left] {$M$};
\fill[popgreen] (Mp) circle (2pt) node[above] {$M'$};
\fill[poppurple] (Mpp) circle (2pt) node[below right] {$M''$};
% projections sur les axes
\draw[popline, dashed] (M) -- (0,1.1);
\draw[popline, dashed] (Mp) -- (0,1.1);
\fill[popblue] (0,1.1) circle (1.4pt) node[left, popdark] {$I_1$};
\draw[popline, dashed] (Mp) -- (3.4,1.1);
\draw[popline, dashed] (Mpp) -- (3.4,1.1);
\fill[poporange] (3.4,1.1) circle (1.4pt) node[right, popdark] {$I_2$};
% flèches de transformation
\draw[->, thick, popblue, bend left=25] (M) to node[midway, above left, popblue] {$s_{d_1}$} (Mp);
\draw[->, thick, poporange, bend left=25] (Mp) to node[midway, above right, poporange] {$s_{d_2}$} (Mpp);
% vecteur translation
\draw[->, very thick, poppurple] (M) -- (Mpp) node[midway, below, poppurple] {$2\overrightarrow{H_1H_2}$};
\end{tikzpicture}
\end{center}
\legende{Deux demi-tours d'axes parallèles se composent en une translation de vecteur $2\overrightarrow{H_1H_2}$.}
\end{popfigure}

\begin{aretenir}[À retenir]
La composée de deux demi-tours d'axes \textbf{parallèles} est une \cle{translation}. Son vecteur est le double du vecteur $\overrightarrow{H_1H_2}$, où $(H_1H_2)$ est la perpendiculaire commune aux deux axes. Attention à l'ordre : $s_{d_2} \circ s_{d_1}$ donne $2\overrightarrow{H_1H_2}$, tandis que $s_{d_1} \circ s_{d_2}$ donne $2\overrightarrow{H_2H_1} = -2\overrightarrow{H_1H_2}$.
\end{aretenir}

\begin{exemplebox}[Exemple chiffré]
Dans un repère orthonormé $(O\,;\vec{i},\vec{j},\vec{k})$, soit $(d_1)$ la droite passant par $A(0\,;0\,;0)$ de vecteur directeur $\vec{k}$, et $(d_2)$ la droite passant par $B(2\,;0\,;0)$ de même vecteur directeur $\vec{k}$. Les deux axes sont verticaux et parallèles ; la perpendiculaire commune est dirigée par $\vec{i}$, avec $\overrightarrow{H_1H_2} = 2\vec{i}$. Donc $s_{d_2} \circ s_{d_1}$ est la translation de vecteur $4\vec{i}$ : tout point $M(x\,;y\,;z)$ a pour image $M''(x+4\,;y\,;z)$. Vérifions avec $M(1\,;3\,;5)$ : $s_{d_1}(M) = M'(-1\,;-3\,;5)$, puis $s_{d_2}(M') = M''(5\,;3\,;5)$. On a bien $\overrightarrow{MM''} = (4\,;0\,;0) = 4\vec{i}$.
\end{exemplebox}

\courssub{Composée de deux demi-tours d'axes perpendiculaires sécants : c'est un demi-tour}

\textbf{Intuition.} Place-toi dans le coin d'une salle de classe : les trois arêtes issues du coin du sol sont concourantes en $O$ et deux à deux perpendiculaires. Effectue un demi-tour autour de la première arête $(d_1)$, puis un demi-tour autour de la seconde $(d_2)$. Le résultat est exactement le demi-tour autour de la troisième arête $(d_3)$.

\begin{propriete}[Théorème 2 : composée de deux demi-tours d'axes perpendiculaires sécants]
Soit $(d_1)$ et $(d_2)$ deux droites \textbf{perpendiculaires} sécantes en un point $O$. Soit $(d_3)$ la droite passant par $O$ et perpendiculaire au plan $(d_1, d_2)$. Alors :
\[ s_{d_2} \circ s_{d_1} = s_{d_3}. \]
De plus, $s_{d_1} \circ s_{d_2} = s_{d_3}$ également : la composée est ici \cle{commutative}.
\end{propriete}

\begin{experience}[Démonstration guidée du théorème 2]
On note $f = s_{d_2} \circ s_{d_1}$. Choisissons un repère orthonormé $(O\,;\vec{i},\vec{j},\vec{k})$ adapté : $(d_1) = (Ox)$, $(d_2) = (Oy)$, $(d_3) = (Oz)$.

\textbf{Étape 1 : expressions analytiques des demi-tours.}
Le demi-tour d'axe $(Ox)$ conserve l'abscisse et change les signes des deux autres coordonnées :
\[ s_{d_1}(x,y,z) = (x,-y,-z). \]
Le demi-tour d'axe $(Oy)$ conserve l'ordonnée et change les signes des deux autres coordonnées :
\[ s_{d_2}(x,y,z) = (-x,y,-z). \]

\textbf{Étape 2 : calcul de la composée.}
Pour un point quelconque $M(x,y,z)$ :
\[ f(M) = s_{d_2}\big(s_{d_1}(x,y,z)\big) = s_{d_2}(x,-y,-z) = (-x,-y,z). \]

\textbf{Étape 3 : identification.}
Or le demi-tour d'axe $(d_3) = (Oz)$ a précisément pour expression analytique $(x,y,z) \mapsto (-x,-y,z)$. Donc $f = s_{d_3}$.

\textbf{Étape 4 : commutativité.}
On vérifie de même que $s_{d_1} \circ s_{d_2}(x,y,z) = s_{d_1}(-x,y,-z) = (-x,-y,z) = s_{d_3}(x,y,z)$. $\blacksquare$

\textbf{Remarque géométrique :} une isométrie de l'espace qui fixe point par point une droite $(d_3)$ et dont la restriction à un plan perpendiculaire est une symétrie centrale est nécessairement le demi-tour d'axe $(d_3)$ : c'est exactement la définition du demi-tour.
\end{experience}

\begin{popfigure}
\begin{center}
\begin{tikzpicture}[scale=1.15, >=Stealth]
\coordinate (O) at (0,0);
% trois axes deux à deux perpendiculaires (coin d'une pièce)
\draw[very thick, popblue, ->] (O) -- (3.6,0) node[right, popblue] {$(d_1)$};
\draw[very thick, poporange, ->] (O) -- (-1.7,-1.5) node[below left, poporange] {$(d_2)$};
\draw[very thick, popgreen, ->] (O) -- (0,3.4) node[above, popgreen] {$(d_3)$};
% prolongements pointillés
\draw[popblue, dashed] (O) -- (-1.2,0);
\draw[poporange, dashed] (O) -- (0.9,0.8);
\draw[popgreen, dashed] (O) -- (0,-1);
% angles droits
\draw[popmuted] (0.45,0) -- (0.45,0.45) -- (0,0.45);
\draw[popmuted] (0.35,0) -- (0.22,-0.26) -- (-0.13,-0.26);
% point O
\fill[popdark] (O) circle (2pt) node[above right=1pt, popdark] {$O$};
% point M et images
\coordinate (M) at (2.6,2.2);
\coordinate (Mp) at (-0.9,0.3);
\coordinate (Mpp) at (-2.3,-1.1);
\fill[popink] (M) circle (2pt) node[right] {$M$};
\fill[popblue] (Mp) circle (1.8pt) node[above left] {$M'=s_{d_1}(M)$};
\fill[poppurple] (Mpp) circle (2pt) node[below] {$M''=s_{d_3}(M)$};
\draw[->, thick, popblue, bend right=15] (M) to (Mp);
\draw[->, thick, poporange, bend right=15] (Mp) to (Mpp);
\draw[->, very thick, popgreen, bend left=30] (M) to node[midway, above, popgreen] {$s_{d_3}$} (Mpp);
\end{tikzpicture}
\end{center}
\legende{Trois axes concourants en $O$ et deux à deux perpendiculaires : $s_{d_2} \circ s_{d_1} = s_{d_3}$.}
\end{popfigure}

\begin{exoresolu}[Les axes $(Ox)$ et $(Oy)$]
Dans un repère orthonormé $(O\,;\vec{i},\vec{j},\vec{k})$, on considère les demi-tours $s_{(Ox)}$ et $s_{(Oy)}$ d'axes les axes des abscisses et des ordonnées. Montrer que $s_{(Oy)} \circ s_{(Ox)}$ est le demi-tour d'axe $(Oz)$, et vérifier sur le point $A(1\,;2\,;3)$.
\tcblower
\textbf{Solution.} Les droites $(Ox)$ et $(Oy)$ sont sécantes en $O$ et perpendiculaires. La droite passant par $O$ et perpendiculaire au plan $(xOy)$ est l'axe $(Oz)$. D'après le théorème 2, $s_{(Oy)} \circ s_{(Ox)} = s_{(Oz)}$.

Vérification directe avec $A(1\,;2\,;3)$ :
\begin{align*}
s_{(Ox)}(A) &= A'(1\,;-2\,;-3) \quad \text{(on change les signes de } y \text{ et } z\text{)},\\
s_{(Oy)}(A') &= A''(-1\,;-2\,;3) \quad \text{(on change les signes de } x \text{ et } z\text{)}.
\end{align*}
Or le demi-tour d'axe $(Oz)$ envoie $(x,y,z)$ sur $(-x,-y,z)$ : il envoie bien $A(1,2,3)$ sur $(-1,-2,3) = A''$. Le résultat est confirmé. $\blacksquare$
\end{exoresolu}

\begin{attention}[Ne pas appliquer les théorèmes hors de leurs hypothèses]
Deux erreurs sont fréquentes au Baccalauréat :
\begin{itemize}
\item Si $(d_1)$ et $(d_2)$ sont \textbf{orthogonales mais non sécantes} (non coplanaires), alors $s_{d_2} \circ s_{d_1}$ n'est \textbf{ni une translation ni un demi-tour} : c'est un \cle{vissage} (composée d'un demi-tour et d'une translation), hors programme détaillé. N'écris jamais $s_{d_2} \circ s_{d_1} = s_{d_3}$ sans avoir vérifié que les axes sont \emph{sécants}.
\item Si les axes sont parallèles, le vecteur de la translation est $2\overrightarrow{H_1H_2}$ et non $\overrightarrow{H_1H_2}$ : le facteur $2$ est la source d'innombrables points perdus.
\end{itemize}
\end{attention}

\courssub{Lien avec les réflexions : une preuve unificatrice}

Tu as vu dans la section précédente que tout demi-tour d'axe $(d)$ s'écrit comme composée de deux réflexions de plans perpendiculaires se coupant selon $(d)$ : $s_d = s_{Q} \circ s_{P}$ où $(P) \cap (Q) = (d)$ et $(P) \perp (Q)$. Cette écriture permet de \emph{retrouver} les deux théorèmes précédents à partir des résultats sur les réflexions (section 3).

\textbf{Cas des axes parallèles.} Soit $(d_1) \parallel (d_2)$, et $(P)$ un plan perpendiculaire aux deux axes. Écrivons $s_{d_1} = s_P \circ s_{P_1}$ et $s_{d_2} = s_{P_2} \circ s_P$, où $(P_1)$ et $(P_2)$ sont des plans contenant respectivement $(d_1)$ et $(d_2)$, perpendiculaires à $(P)$. Comme $(d_1) \parallel (d_2)$, on peut choisir $(P_1) \parallel (P_2)$. Alors :
\[ s_{d_2} \circ s_{d_1} = s_{P_2} \circ s_P \circ s_P \circ s_{P_1} = s_{P_2} \circ s_{P_1}, \]
car $s_P \circ s_P = \mathrm{Id}$. Or la composée de deux réflexions de plans \textbf{parallèles} est une translation de vecteur $2\overrightarrow{H_1H_2}$ : on retrouve le théorème 1.

\textbf{Cas des axes perpendiculaires sécants.} Écrivons $s_{d_1} = s_{P} \circ s_{Q}$ et $s_{d_2} = s_{R} \circ s_{P}$ où $(P)$ est le plan $(d_1, d_2)$, $(Q)$ le plan contenant $(d_1)$ perpendiculaire à $(P)$, et $(R)$ le plan contenant $(d_2)$ perpendiculaire à $(P)$. Alors :
\[ s_{d_2} \circ s_{d_1} = s_{R} \circ s_{P} \circ s_{P} \circ s_{Q} = s_{R} \circ s_{Q}. \]
Or $(Q)$ et $(R)$ sont deux plans \textbf{perpendiculaires} (car $(d_1) \perp (d_2)$) dont l'intersection est la droite $(d_3)$ passant par $O$ perpendiculaire au plan $(d_1,d_2)$. La composée de deux réflexions de plans perpendiculaires est le demi-tour d'axe leur intersection : on retrouve le théorème 2.

\begin{savaistu}[Pourquoi cette astuce est puissante]
L'astuce consistant à « insérer » $s_P \circ s_P = \mathrm{Id}$ entre deux isométries s'appelle une \emph{conjonction de réflexions}. Elle est au cœur de la classification générale des isométries : toute isométrie de l'espace est composée d'au plus quatre réflexions. C'est ainsi que les mathématiciens démontrent qu'il n'existe que quelques types d'isométries de l'espace : translations, rotations, réflexions, symétries centrales, vissages et antirotations.
\end{savaistu}

\courssub{Méthode de reconnaissance et synthèse générale}

\begin{methode}[Reconnaître la nature de $s_{d_2} \circ s_{d_1}$]
Face à une composée de deux demi-tours, procède toujours ainsi :
\begin{enumerate}
\item \textbf{Déterminer la position relative des axes} $(d_1)$ et $(d_2)$ : comparer leurs vecteurs directeurs (colinéaires ? orthogonaux ?) et chercher un point commun.
\item \textbf{Si les axes sont strictement parallèles} : conclure que la composée est une \cle{translation}, puis déterminer la perpendiculaire commune $(H_1H_2)$ et écrire le vecteur $2\overrightarrow{H_1H_2}$ (attention à l'ordre de composition).
\item \textbf{Si les axes sont perpendiculaires et sécants en $O$} : conclure que la composée est le \cle{demi-tour} d'axe $(d_3)$, droite passant par $O$ perpendiculaire au plan $(d_1,d_2)$.
\item \textbf{Si les axes sont sécants non perpendiculaires} : la composée est une \cle{rotation} d'axe la perpendiculaire commune en $O$ (cas général, hors programme détaillé).
\item \textbf{Si les axes sont non coplanaires} : aucun des théorèmes ne s'applique ; la composée est un vissage (hors programme).
\item \textbf{Vérifier} le résultat sur un point bien choisi, comme dans l'exercice résolu.
\end{enumerate}
\end{methode}

\begin{popfigure}
\begin{center}
\begin{tikzpicture}[
  box/.style={draw=popblue, thick, rounded corners=4pt, fill=popblueL, text width=3.6cm, align=center, minimum height=1.1cm, font=\small},
  res/.style={draw=popgreen, thick, rounded corners=4pt, fill=popgreenL, text width=3.9cm, align=center, minimum height=1.1cm, font=\small},
  >=Stealth]
\node[box, align=center] (a) at (0,2.4) {\textbf{Réflexions}\\ plans $(P_1) \parallel (P_2)$};
\node[res, align=center] (b) at (5.6,2.4) {\textbf{Translation}\\ vecteur $2\overrightarrow{H_1H_2}$};
\node[box, align=center] (c) at (0,0.6) {\textbf{Réflexions}\\ plans $(P_1) \perp (P_2)$};
\node[res, align=center] (d) at (5.6,0.6) {\textbf{Demi-tour}\\ d'axe $(P_1)\cap(P_2)$};
\node[box, align=center] (e) at (0,-1.4) {\textbf{Demi-tours}\\ axes $(d_1) \parallel (d_2)$};
\node[res, align=center] (f) at (5.6,-1.4) {\textbf{Translation}\\ vecteur $2\overrightarrow{H_1H_2}$};
\node[box, align=center] (g) at (0,-3.4) {\textbf{Demi-tours}\\ axes $\perp$ sécants en $O$};
\node[res, align=center] (h) at (5.6,-3.4) {\textbf{Demi-tour} d'axe $(d_3)$\\ $(d_3) \perp$ plan $(d_1,d_2)$ en $O$};
\draw[->, very thick, popdark] (a) -- (b);
\draw[->, very thick, popdark] (c) -- (d);
\draw[->, very thick, popdark] (e) -- (f);
\draw[->, very thick, popdark] (g) -- (h);
\end{tikzpicture}
\end{center}
\legende{Tableau-synthèse des composées de deux réflexions ou de deux demi-tours.}
\end{popfigure}

\begin{aretenir}[L'essentiel de la section]
\begin{itemize}
\item Axes \textbf{parallèles} : $s_{d_2} \circ s_{d_1} = t_{2\overrightarrow{H_1H_2}}$ (translation).
\item Axes \textbf{perpendiculaires sécants en $O$} : $s_{d_2} \circ s_{d_1} = s_{d_3}$, où $(d_3)$ est la perpendiculaire en $O$ au plan $(d_1,d_2)$. La composition est alors commutative.
\item Ces résultats se retrouvent en écrivant chaque demi-tour comme composée de deux réflexions et en utilisant $s_P \circ s_P = \mathrm{Id}$.
\item Hors de ces deux configurations, aucune conclusion simple n'est autorisée.
\end{itemize}
\end{aretenir}

\begin{exercice}[Pour t'entraîner]
Dans un repère orthonormé $(O\,;\vec{i},\vec{j},\vec{k})$, on considère les droites $(d_1)$ d'équations $x = 0,\ z = 1$ (parallèle à $(Oy)$) et $(d_2)$ d'équations $x = 3,\ z = 1$ (parallèle à $(Oy)$).
\begin{enumerate}
\item Justifier que $(d_1)$ et $(d_2)$ sont parallèles et déterminer leur perpendiculaire commune.
\item En déduire la nature et les éléments caractéristiques de $s_{d_2} \circ s_{d_1}$.
\item Donner l'image du point $M(1\,;2\,;4)$ par cette composée, puis vérifier en calculant successivement $s_{d_1}(M)$ puis son image par $s_{d_2}$.
\end{enumerate}
\end{exercice}

\coursec{Reconnaissance d'isométries et résolution de problèmes}

Dans la section précédente, nous avons appris à composer deux demi-tours et à identifier la nature de la composée : une translation lorsque les axes sont parallèles, un demi-tour lorsque les axes sont orthogonaux et concourants. Nous disposons désormais de toute la panoplie des isométries de l'espace étudiées au programme : réflexions, demi-tours, translations, et leurs composées. Reste la question décisive, celle que l'on rencontre systématiquement au Baccalauréat : \emph{on te donne une application par son expression analytique, ou par une construction géométrique — sauras-tu dire qui elle est, le prouver, et t'en servir pour résoudre un problème ?} C'est l'objet de cette dernière section, qui est aussi la plus stratégique pour l'examen.

\courssub{Reconnaître une isométrie donnée par son expression analytique}

\textbf{L'idée directrice.} Chaque isométrie laisse une « signature » dans son expression analytique. Une réflexion change le signe d'\emph{une seule} coordonnée (celle qui est perpendiculaire au plan miroir). Un demi-tour change le signe de \emph{deux} coordonnées (celles qui sont perpendiculaires à l'axe). Une translation \emph{ajoute} un vecteur constant sans toucher aux signes. Mais attention : cette lecture des signes n'est fiable que lorsque les plans ou axes sont parallèles aux plans de coordonnées. L'outil universel, lui, c'est la recherche des \cle{points fixes}.

\begin{methode}[Fiche de reconnaissance d'une isométrie à partir de son expression analytique]
Soit $f$ l'application de l'espace définie par $f(x\,;y\,;z) = (x'\,;y'\,;z')$.
\begin{enumerate}
\item \textbf{Observer la forme de l'expression.} Si $x'$, $y'$, $z'$ s'obtiennent en ajoutant des constantes à $x$, $y$, $z$ (par exemple $x' = x + 2$), c'est une \cle{translation}. Sinon, passer à l'étape 2.
\item \textbf{Résoudre l'équation des points fixes} $f(M) = M$, c'est-à-dire le système $x' = x$, $y' = y$, $z' = z$.
\begin{itemize}
\item L'ensemble des solutions est un \textbf{plan} $\mathcal{P}$ : $f$ est la \cle{réflexion} de plan $\mathcal{P}$.
\item L'ensemble des solutions est une \textbf{droite} $(\Delta)$ : $f$ est le \cle{demi-tour} d'axe $(\Delta)$.
\item L'ensemble des solutions est \textbf{vide} : $f$ n'est ni une réflexion ni un demi-tour ; revenir à la forme analytique (translation, ou composée à étudier).
\end{itemize}
\item \textbf{Contrôler avec la signature des signes} : une coordonnée change de signe $\Rightarrow$ réflexion ; deux coordonnées changent de signe $\Rightarrow$ demi-tour.
\item \textbf{Conclure en rédigeant} : nature de l'isométrie \emph{et} éléments caractéristiques (plan de la réflexion, ou axe du demi-tour donné par un point et un vecteur directeur, ou vecteur de la translation).
\end{enumerate}
\end{methode}

\begin{popfigure}
\begin{tikzpicture}[
  node distance=0.55cm and 1.1cm,
  decision/.style={diamond, draw=poporange, fill=poporangeL, text width=3.1cm, align=center, inner sep=1pt, font=\small},
  block/.style={rectangle, rounded corners, draw=popblue, fill=popblueL, text width=3.4cm, align=center, font=\small},
  start/.style={rectangle, rounded corners, draw=popdark, fill=popdark, text=white, text width=4.6cm, align=center, font=\small},
  arr/.style={-{Stealth[length=2.5mm]}, thick, popdark}]
\node[start] (start) {Expression analytique de $f$ : $M(x;y;z) \mapsto M'(x';y';z')$};
\node[decision, below=of start] (trans) {Forme $x'=x+a$, $y'=y+b$, $z'=z+c$ ?};
\node[block, right=of trans, xshift=3.6cm] (tr) {\textbf{Translation} de vecteur $\vec{u}(a\,;b\,;c)$};
\node[decision, below=of trans] (fixe) {Résoudre $f(M)=M$ : ensemble des points fixes ?};
\node[block, below left=0.6cm and -1.2cm of fixe] (plan) {Un plan $\mathcal{P}$ fixe $\Rightarrow$ \textbf{réflexion} de plan $\mathcal{P}$};
\node[block, below=0.6cm of fixe] (droite) {Une droite $(\Delta)$ fixe $\Rightarrow$ \textbf{demi-tour} d'axe $(\Delta)$};
\node[block, below right=0.6cm and -1.2cm of fixe] (vide) {Aucun point fixe $\Rightarrow$ composée : analyser la forme};
\draw[arr] (start) -- (trans);
\draw[arr] (trans) -- node[above, font=\small]{oui} (tr);
\draw[arr] (trans) -- node[right, font=\small]{non} (fixe);
\draw[arr] (fixe) -| (plan);
\draw[arr] (fixe) -- (droite);
\draw[arr] (fixe) -| (vide);
\end{tikzpicture}
\legende{Arbre de décision : de l'expression analytique à la nature de l'isométrie.}
\end{popfigure}

\begin{exemplebox}[Identifier $f(x\,;y\,;z) = (-x\,;\,y\,;\,-z+4)$]
Cherchons les points fixes : $f(M) = M$ équivaut à
\[
\begin{cases} -x = x \\ y = y \\ -z + 4 = z \end{cases}
\iff
\begin{cases} x = 0 \\ y \in \mathbb{R} \\ z = 2. \end{cases}
\]
L'ensemble des points fixes est la droite $(\Delta)$ passant par $\Omega(0\,;0\,;2)$ et dirigée par $\vec{j}(0\,;1\,;0)$. Deux coordonnées ($x$ et $z$) changent de signe : $f$ est le \textbf{demi-tour d'axe $(\Delta)$}, droite d'équations $x = 0$ et $z = 2$.
\end{exemplebox}

\begin{attention}[Une isométrie sans point fixe n'est pas forcément une translation]
Si l'équation $f(M) = M$ n'a pas de solution, ne conclus pas trop vite « c'est une translation ». Par exemple, la composée d'un demi-tour et d'une translation de vecteur parallèle à l'axe (un \emph{vissage}) ou la composée d'une réflexion et d'une translation de vecteur parallèle au plan (une \emph{symétrie glissée}) n'ont aucun point fixe sans être des translations. Vérifie \emph{toujours} la forme analytique : une translation s'écrit $x' = x + a$, $y' = y + b$, $z' = z + c$ avec les coordonnées \emph{inchangées de signe}. Si des signes changent \emph{et} qu'il n'y a pas de point fixe, tu es en présence d'une composée qu'il faut décomposer.
\end{attention}

\begin{exoresolu}[Reconnaissance complète d'une isométrie]
L'espace est rapporté à un repère orthonormé $(O\,;\vec{i},\vec{j},\vec{k})$. On considère l'application $g$ qui à tout point $M(x\,;y\,;z)$ associe $M'(x'\,;y'\,;z')$ avec
\[
x' = x + 2, \qquad y' = -y, \qquad z' = z - 1.
\]
Déterminer la nature de $g$ et ses éléments caractéristiques.
\tcblower
\textbf{Étape 1 : points fixes.} $g(M) = M$ donne le système $x + 2 = x$, $-y = y$, $z - 1 = z$. La première équation $2 = 0$ est impossible : \textbf{aucun point fixe}. Ce n'est donc ni une réflexion ni un demi-tour.

\textbf{Étape 2 : lecture de la forme.} La coordonnée $y$ change de signe : cela évoque la réflexion $s$ de plan $\mathcal{P}$ d'équation $y = 0$ (le plan $(O\,;\vec{i},\vec{k})$). Les coordonnées $x$ et $z$ reçoivent des constantes : cela évoque la translation $t$ de vecteur $\vec{u}(2\,;0\,;-1)$.

\textbf{Étape 3 : vérification par composition.} Calculons $t \circ s$ : si $M(x;y;z)$, alors $s(M) = (x\,;-y\,;z)$ puis $t(s(M)) = (x+2\,;-y\,;z-1) = g(M)$. Donc $g = t \circ s$. Comme $\vec{u}$ est \emph{parallèle au plan} $\mathcal{P}$ (sa composante selon $\vec{j}$ est nulle), $t$ et $s$ commutent et $g$ est une \textbf{symétrie glissée} : réflexion de plan $\mathcal{P} : y = 0$ suivie de la translation de vecteur $\vec{u}(2\,;0\,;-1)$.
\end{exoresolu}

\courssub{Problème de distance minimale : la technique du point symétrique}

\textbf{La situation.} Un technicien d'Orange Cameroun doit poser un câble fibre optique d'un central $A$ à un relais $B$, en passant obligatoirement par un point $M$ d'un mur (plan $\mathcal{P}$) où se trouve une armoire de connexion. Où placer $M$ pour que la longueur totale $AM + MB$ soit minimale ? C'est le problème classique du \cle{plus court chemin avec réflexion}, et l'outil miracle est la réflexion de plan $\mathcal{P}$.

\begin{methode}[Technique du point symétrique]
Pour minimiser $AM + MB$ lorsque $M$ décrit un plan $\mathcal{P}$, $A$ et $B$ étant du même côté de $\mathcal{P}$ :
\begin{enumerate}
\item Construire $A'$, symétrique de $A$ par rapport à $\mathcal{P}$ (image de $A$ par la réflexion de plan $\mathcal{P}$).
\item Pour tout point $M$ de $\mathcal{P}$, on a $AM = A'M$ (la réflexion est une isométrie et $M$ est fixe), donc $AM + MB = A'M + MB$.
\item Par l'inégalité triangulaire, $A'M + MB \geq A'B$, avec égalité si et seulement si $M$ appartient au segment $[A'B]$.
\item Le point optimal est donc $M_0 = [A'B] \cap \mathcal{P}$, et la distance minimale vaut $A'B$.
\end{enumerate}
\end{methode}

\begin{popfigure}
\begin{tikzpicture}[scale=1.05, >=Stealth]
% Plan P vu de profil (le mur)
\draw[very thick, popblue] (0,-1.6) -- (0,3.4) node[above, popblue]{mur : plan $\mathcal{P}$};
\fill[pattern=north east lines, pattern color=popblue!40] (0,-1.6) rectangle (0.25,3.4);
% Points
\coordinate (A) at (-3,2.6);
\coordinate (Apr) at (3,2.6);
\coordinate (B) at (-2.2,-0.9);
\coordinate (M0) at (0,0.62);
\coordinate (M) at (0,2.1);
% Chemin optimal
\draw[very thick, popgreen] (A) -- (M0) -- (B);
\draw[dashed, popgreen] (M0) -- (Apr);
\draw[dashed, popmuted] (A) -- (Apr);
% Autre chemin
\draw[thick, poporange, dashed] (A) -- (M) -- (B);
% Points et labels
\fill[popink] (A) circle (2pt) node[left]{$A$};
\fill[popink] (B) circle (2pt) node[left]{$B$};
\fill[poppurple] (Apr) circle (2pt) node[right]{$A' = s_{\mathcal{P}}(A)$};
\fill[popgreen] (M0) circle (2.2pt) node[below right=1pt, popgreen]{$M_0$};
\fill[poporange] (M) circle (2pt) node[above right=0pt, poporange]{$M$};
\end{tikzpicture}
\legende{Plus court chemin de $A$ à $B$ via le plan $\mathcal{P}$ : le chemin vert $A \to M_0 \to B$ se « déplie » en segment $[A'B]$. Tout autre point $M$ donne un chemin plus long (orange).}
\end{popfigure}

\begin{exoresolu}[Câble de longueur minimale]
Dans un repère orthonormé, le mur est le plan $\mathcal{P}$ d'équation $x = 0$, le central est $A(4\,;1\,;2)$ et le relais $B(2\,;5\,;6)$. Déterminer le point $M_0$ de $\mathcal{P}$ minimisant $AM + MB$ et calculer cette longueur minimale.
\tcblower
\textbf{Étape 1 : symétrique de $A$.} La réflexion de plan $\mathcal{P} : x = 0$ envoie $(x;y;z)$ sur $(-x;y;z)$, donc $A'(-4\,;1\,;2)$.

\textbf{Étape 2 : intersection de $[A'B]$ avec $\mathcal{P}$.} Paramétrons : un point de $(A'B)$ s'écrit $M(t) = A' + t\overrightarrow{A'B}$ avec $\overrightarrow{A'B}(6\,;4\,;4)$, soit $M(t) = (-4 + 6t\,;\,1 + 4t\,;\,2 + 4t)$. La condition $x = 0$ donne $-4 + 6t = 0$, donc $t = \dfrac{2}{3}$. D'où
\[
M_0\left(0\,;\,1 + \tfrac{8}{3}\,;\,2 + \tfrac{8}{3}\right) = \left(0\,;\,\tfrac{11}{3}\,;\,\tfrac{14}{3}\right).
\]

\textbf{Étape 3 : longueur minimale.} Elle vaut
\[
A'B = \sqrt{6^2 + 4^2 + 4^2} = \sqrt{36 + 16 + 16} = \sqrt{68} = 2\sqrt{17} \approx 8{,}25.
\]
Le câble minimal mesure donc $2\sqrt{17}$ unités, posé via le point $M_0$.
\end{exoresolu}

\courssub{Utiliser les isométries pour démontrer des propriétés géométriques}

Les isométries \emph{conservent les distances, les angles, les alignements, les milieux et l'orthogonalité}. C'est un outil de démonstration redoutable, souvent plus rapide qu'un calcul de coordonnées.

\begin{propriete}[Conservation par les isométries]
Toute isométrie de l'espace (réflexion, demi-tour, translation, et leurs composées) conserve : les distances ($MN = M'N'$), les milieux, l'alignement, l'orthogonalité, et transforme un plan en plan, une droite en droite, une sphère en sphère de même rayon.
\end{propriete}

\textbf{Application 1 : plan médiateur.} Le plan médiateur $\mathcal{P}$ d'un segment $[AB]$ est le plan passant par le milieu $I$ de $[AB]$ et perpendiculaire à $(AB)$. Montrons, à l'aide de la réflexion $s_{\mathcal{P}}$, que c'est exactement l'ensemble des points équidistants de $A$ et $B$.
\begin{itemize}
\item \textbf{Sens direct.} La réflexion $s_{\mathcal{P}}$ échange $A$ et $B$ (car $\mathcal{P}$ est perpendiculaire à $(AB)$ en son milieu). Si $M \in \mathcal{P}$, alors $M$ est fixe par $s_{\mathcal{P}}$, et comme une réflexion conserve les distances : $MA = s_{\mathcal{P}}(M)\,s_{\mathcal{P}}(A) = MB$.
\item \textbf{Réciproque.} Soit $M$ tel que $MA = MB$. Le plus efficace est le produit scalaire : $MA = MB$ équivaut, en développant $MA^2 = MB^2$ avec $I$ milieu de $[AB]$, à $\overrightarrow{IM} \cdot \overrightarrow{AB} = 0$, c'est-à-dire $M \in \mathcal{P}$. La réflexion fournit le sens géométrique, le produit scalaire sécurise la réciproque.
\end{itemize}

\textbf{Application 2 : symétries d'un cube.} Un cube $ABCDEFGH$ admet de nombreux éléments de symétrie. Ses \textbf{plans de symétrie} sont de deux types : les 3 plans parallèles aux faces passant par le centre du cube (par exemple le plan médiateur de $[AD]$, qui échange les faces $ABFE$ et $DCGH$) et les 6 plans diagonaux contenant deux arêtes opposées parallèles (par exemple le plan $(ACGE)$). Le cube possède donc \textbf{9 plans de symétrie}. Il admet aussi des \textbf{axes de demi-tour} : les droites joignant les centres de deux faces opposées, et les droites joignant les milieux de deux arêtes opposées. Chacune de ces réflexions ou demi-tours laisse le cube \emph{globalement invariant}.

\begin{exoresolu}[Plan de symétrie d'un cube]
Soit $ABCDEFGH$ un cube d'arête $a$ ($ABCD$ face inférieure, $EFGH$ face supérieure, $E$ au-dessus de $A$). Soient $I$, $J$, $K$, $L$ les milieux respectifs de $[AB]$, $[DC]$, $[EF]$ et $[HG]$. Montrer que le plan $\mathcal{P} = (IJLK)$ est un plan de symétrie du cube, et en déduire une propriété des sections planes du cube.
\tcblower
\textbf{Nature de $\mathcal{P}$.} Les arêtes $[AB]$, $[DC]$, $[EF]$ et $[HG]$ sont parallèles et de même longueur. La droite $(AB)$ est perpendiculaire à $(AD)$ et à $(AE)$, donc au plan $(ADHE)$ ; les quatre milieux $I$, $J$, $K$, $L$ sont à la distance $\frac{a}{2}$ de ce plan. Le plan $\mathcal{P}$ passe donc par $I$ et est perpendiculaire à $(AB)$ : c'est le \textbf{plan médiateur de $[AB]$}, et aussi de $[DC]$, $[EF]$ et $[HG]$.

\textbf{Vérification en repère.} Prenons le repère orthonormé tel que $A(0;0;0)$, $B(a;0;0)$, $D(0;a;0)$, $E(0;0;a)$. Alors $\mathcal{P}$ a pour équation $x = \dfrac{a}{2}$ et la réflexion $s$ de plan $\mathcal{P}$ s'écrit $s(x;y;z) = (a - x\,;\,y\,;\,z)$. Sommet par sommet :
\[
s(A) = B,\quad s(B) = A,\quad s(C) = D,\quad s(D) = C,\quad s(E) = F,\quad s(F) = E,\quad s(G) = H,\quad s(H) = G.
\]
L'ensemble des huit sommets est globalement invariant ; comme $s$ est une isométrie, elle conserve les segments et les faces, donc l'image du cube est le cube lui-même : $\mathcal{P}$ est bien un \textbf{plan de symétrie du cube}.

Conséquence : toute section du cube par un plan $\mathcal{Q}$ est transformée par $s$ en la section par $s(\mathcal{Q})$. Si $\mathcal{Q}$ est globalement invariant par $s$, sa section est une figure symétrique par rapport à la droite $\mathcal{P} \cap \mathcal{Q}$. C'est ainsi qu'on démontre élégamment que certaines sections planes du cube sont des rectangles ou des hexagones admettant un axe de symétrie.
\end{exoresolu}

\courssub{Décomposer une isométrie en composée de réflexions ou de demi-tours}

Rappel des résultats des sections précédentes : la composée de deux réflexions de plans \emph{parallèles} est une translation de vecteur double du vecteur reliant les plans ; la composée de deux réflexions de plans \emph{perpendiculaires} est le demi-tour d'axe leur droite d'intersection. On peut « remonter » ces résultats pour décomposer une isométrie.

\begin{methode}[Décomposer une translation ou un demi-tour]
\begin{enumerate}
\item \textbf{Translation de vecteur $\vec{u}$} : choisir deux plans parallèles $\mathcal{P}_1$ et $\mathcal{P}_2$, perpendiculaires à $\vec{u}$, distants de $\dfrac{\|\vec{u}\|}{2}$, $\mathcal{P}_2$ étant situé dans le sens de $\vec{u}$ par rapport à $\mathcal{P}_1$. Alors $t_{\vec{u}} = s_{\mathcal{P}_2} \circ s_{\mathcal{P}_1}$. Il y a une infinité de couples possibles.
\item \textbf{Demi-tour d'axe $(\Delta)$} : choisir deux plans perpendiculaires $\mathcal{P}_1$ et $\mathcal{P}_2$ se coupant selon $(\Delta)$. Alors $s_{(\Delta)} = s_{\mathcal{P}_2} \circ s_{\mathcal{P}_1}$. Là encore, infinité de choix.
\item \textbf{Demi-tour comme composée de demi-tours} : la composée de deux demi-tours d'axes orthogonaux et concourants est un demi-tour dont l'axe est perpendiculaire aux deux premiers ; réciproquement, un demi-tour d'axe $(\Delta)$ peut s'écrire comme composée de deux demi-tours d'axes concourants, orthogonaux à $(\Delta)$ et orthogonaux entre eux.
\end{enumerate}
\end{methode}

\begin{exemplebox}[Décomposition chiffrée]
Soit $t$ la translation de vecteur $\vec{u}(0\,;0\,;6)$. Prenons $\mathcal{P}_1 : z = 0$ et $\mathcal{P}_2 : z = 3$ (plans parallèles, perpendiculaires à $\vec{u}$, distants de $3 = \tfrac{6}{2}$). Alors $s_{\mathcal{P}_1}(x;y;z) = (x;y;-z)$ et $s_{\mathcal{P}_2}(x;y;z) = (x;y;6-z)$. Composons :
\[
s_{\mathcal{P}_2}\big(s_{\mathcal{P}_1}(x;y;z)\big) = s_{\mathcal{P}_2}(x;y;-z) = (x\,;\,y\,;\,6+z) = t(M).
\]
On retrouve bien $t = s_{\mathcal{P}_2} \circ s_{\mathcal{P}_1}$.
\end{exemplebox}

\courssub{Entraînement type Bac : rédaction complète}

\begin{exoresolu}[Sujet type Baccalauréat C/E]
L'espace est muni d'un repère orthonormé $(O\,;\vec{i},\vec{j},\vec{k})$. On considère les points $A(1\,;0\,;2)$, $B(3\,;0\,;2)$ et l'application $f$ définie par
\[
f\big(M(x\,;y\,;z)\big) = M'(4 - x\,;\,-y\,;\,z).
\]
\begin{enumerate}
\item Déterminer la nature de $f$ et ses éléments caractéristiques.
\item Écrire $f$ comme composée de deux réflexions que l'on précisera.
\item Soit $\mathcal{S}$ la sphère de centre $A$ et de rayon $2$. Déterminer l'image de $\mathcal{S}$ par $f$.
\end{enumerate}
\tcblower
\textbf{1. Nature de $f$.} Cherchons les points fixes : $f(M) = M \iff \begin{cases} 4 - x = x \\ -y = y \\ z = z \end{cases} \iff \begin{cases} x = 2 \\ y = 0 \\ z \in \mathbb{R}. \end{cases}$

L'ensemble des points fixes est la droite $(\Delta)$ passant par $\Omega(2\,;0\,;0)$ et dirigée par $\vec{k}$. Deux coordonnées changent de signe : $f$ est le \textbf{demi-tour d'axe $(\Delta)$}, droite d'équations $x = 2$ et $y = 0$.

\textbf{2. Décomposition.} L'axe $(\Delta)$ est l'intersection des plans perpendiculaires $\mathcal{P}_1 : x = 2$ et $\mathcal{P}_2 : y = 0$. Vérifions : $s_{\mathcal{P}_1}(x;y;z) = (4 - x\,;\,y\,;\,z)$ puis $s_{\mathcal{P}_2}(4 - x\,;\,y\,;\,z) = (4 - x\,;\,-y\,;\,z) = f(M)$. Donc
\[
f = s_{\mathcal{P}_2} \circ s_{\mathcal{P}_1}, \quad \mathcal{P}_1 : x = 2, \ \mathcal{P}_2 : y = 0.
\]

\textbf{3. Image de la sphère.} $f$ est une isométrie, donc l'image de $\mathcal{S}$ est la sphère de même rayon $2$ et de centre $f(A)$. Or $f(1\,;0\,;2) = (3\,;0\,;2) = B$. L'image de $\mathcal{S}$ est donc la sphère de centre $B(3\,;0\,;2)$ et de rayon $2$. (Vérification de cohérence : le milieu de $[AB]$ est $(2\,;0\,;2)$, qui appartient bien à $(\Delta)$, et $(AB)$, dirigée par $\vec{i}$, est perpendiculaire à $(\Delta)$ : $A$ et $B$ sont bien symétriques par rapport à l'axe du demi-tour.)
\end{exoresolu}

\begin{exercice}[Pour t'entraîner]
\begin{enumerate}
\item Déterminer la nature et les éléments caractéristiques de l'application $h$ définie par $h(x\,;y\,;z) = (x\,;\,4 - y\,;\,z)$.
\item Montrer que l'application $k(x\,;y\,;z) = (-x + 2\,;\,-y\,;\,z + 3)$ n'a aucun point fixe. En remarquant que $k = t \circ s_{(\Delta)}$ où $s_{(\Delta)}$ est un demi-tour à préciser et $t$ une translation de vecteur parallèle à $(\Delta)$, décomposer $k$.
\item Deux villes $A(0\,;0\,;0)$ et $B(6\,;4\,;0)$ doivent être reliées par une canalisation passant par un point $M$ du plan $\mathcal{P} : z = 2$. Déterminer $M$ pour que $AM + MB$ soit minimal.
\end{enumerate}
\end{exercice}

\begin{corrige}[Exercice]
\textbf{1.} Points fixes : $x = x$, $4 - y = y$, $z = z$ donnent $y = 2$, $x$ et $z$ libres : c'est le plan $\mathcal{P} : y = 2$. Une seule coordonnée change de signe : $h$ est la \textbf{réflexion de plan $\mathcal{P} : y = 2$}.

\textbf{2.} $k(M) = M$ exigerait $-x + 2 = x$ soit $x = 1$, $y = 0$, et $z + 3 = z$ : impossible, donc aucun point fixe. Écrivons $k(x;y;z) = (2 - x\,;\,-y\,;\,z) + (0\,;\,0\,;\,3)$ : c'est $t \circ s_{(\Delta)}$ où $s_{(\Delta)}$ est le demi-tour d'axe $(\Delta) : x = 1,\ y = 0$ (droite passant par $(1;0;0)$, dirigée par $\vec{k}$) et $t$ la translation de vecteur $3\vec{k}$, parallèle à $(\Delta)$. C'est un \textbf{vissage} (demi-tour suivi d'une translation de vecteur parallèle à son axe).

\textbf{3.} $A'(0\,;0\,;4)$ est le symétrique de $A$ par rapport à $\mathcal{P}$. La droite $(A'B)$ : $M(t) = (6t\,;\,4t\,;\,4 - 4t)$ ; la condition $z = 2$ donne $t = \tfrac{1}{2}$, d'où $M(3\,;\,2\,;\,2)$ et la longueur minimale $A'B = \sqrt{36 + 16 + 16} = \sqrt{68} = 2\sqrt{17}$.
\end{corrige}

\begin{savaistu}[Les isométries au Baccalauréat]
Au Bac C et E, les isométries de l'espace apparaissent très régulièrement dans les problèmes de géométrie, presque toujours \emph{combinées avec le produit scalaire} : on te demande d'abord de vérifier qu'un plan a telle équation (produit scalaire avec un vecteur normal), puis de reconnaître une application, de déterminer l'image d'une sphère ou d'un plan, ou de résoudre un problème de distance. Les sessions récentes ont aussi aimé les questions de type « démontrer que $f$ est un demi-tour » en passant par les points fixes. Maîtriser la fiche de reconnaissance de cette section, c'est sécuriser 3 à 5 points du problème de géométrie. Hors du lycée, ces mêmes isométries régissent la cristallographie (structure des minerais, comme la bauxite de l'Adamaoua), la robotique et l'imagerie 3D.
\end{savaistu}

\newpage

\coursec{Activité d'intégration}

\coursec{Isométries de l'espace}

\courssub{Objectifs de l'activité}

À l'issue de cette activité d'intégration, tu seras capable de :
\begin{itemize}
\item déterminer les \cle{expressions analytiques} de réflexions de plans parallèles aux plans de coordonnées et les appliquer à des points concrets ;
\item composer deux réflexions de plans \cle{parallèles} et reconnaître la \cle{translation} obtenue, en précisant son vecteur ;
\item composer deux réflexions de plans \cle{perpendiculaires} et reconnaître le \cle{demi-tour} obtenu, en précisant son axe ;
\item utiliser une réflexion pour résoudre un \cle{problème d'optimisation} de trajet (problème du chemin le plus court via un plan) ;
\item produire, en groupe, un raisonnement rédigé, des calculs analytiques et une maquette schématique cotée.
\end{itemize}

\begin{methode}[Déterminer l'expression analytique d'une réflexion par rapport à un plan parallèle à un plan de coordonnées]
\begin{enumerate}
\item Identifier l'équation du plan : $x = a$, $y = b$ ou $z = c$.
\item La coordonnée correspondante change par symétrie par rapport à la cote du plan : $x' = 2a - x$, $y' = 2b - y$ ou $z' = 2c - z$.
\item Les deux autres coordonnées restent inchangées.
\item Vérifier : le milieu du segment $[MM']$ appartient au plan et le vecteur $\overrightarrow{MM'}$ est orthogonal au plan (colinéaire au vecteur normal).
\end{enumerate}
\end{methode}

\begin{attention}[Pièges fréquents]
\begin{itemize}
\item \textbf{Ordre de la composée} : $s_Q \circ s_P$ signifie « d'abord $s_P$, ensuite $s_Q$ ». Ne pas inverser l'ordre, le vecteur de translation changerait de sens.
\item \textbf{Signe dans la formule} : Pour le plan $y = b$, l'image de $y$ est $2b - y$ (et non $y - 2b$ ou $b - y$). Tester avec un point du plan : si $y=b$, alors $y'=b$.
\item \textbf{Plans perpendiculaires} : Vérifier l'orthogonalité des vecteurs \textbf{normaux}, pas seulement des plans. Deux plans sont perpendiculaires $\iff$ leurs vecteurs normaux sont orthogonaux.
\item \textbf{Demi-tour vs Rotation} : Un demi-tour est une rotation d'angle $\pi$ (180\textdegree). L'axe est la droite d'intersection des deux plans miroirs. Ne pas confondre avec une symétrie centrale (point).
\end{itemize}
\end{attention}

\courssub{Situation}

Monsieur Etoundi tient, avenue Kennedy à Yaoundé, un showroom de prêt-à-porter. Pour donner une impression d'espace et permettre aux clients de s'admirer sous tous les angles, il fait installer deux grands miroirs plans verticaux. L'atelier du menuisier métallique travaille dans un repère orthonormé $(O\,;\vec{\imath},\vec{\jmath},\vec{k})$ de l'espace, où le plan $(O\,;\vec{\imath},\vec{\jmath})$ est le sol, l'unité étant le mètre.

\begin{itemize}
\item Le premier miroir est adossé au mur du fond : c'est le plan $P$ d'équation $y = 0$.
\item Le second miroir est une cloison mobile parallèle au mur : c'est le plan $Q$ d'équation $y = 2{,}5$.
\end{itemize}

Un mannequin de vitrine est placé au point $M(1\,;\,0{,}8\,;\,1{,}2)$ : il est à $\SI{0,8}{m}$ du mur $P$ et sa tête culmine à $\SI{1,2}{m}$ du sol.

Monsieur Etoundi observe un phénomène amusant : dans le miroir $Q$, il voit non seulement le reflet du mannequin, mais aussi le reflet \emph{du reflet} dans $P$. Il se demande de combien ce « reflet du reflet » est décalé par rapport au mannequin réel. Plus tard, il envisage de remplacer la cloison par un miroir \emph{perpendiculaire} au mur, le plan $R$ d'équation $x = 3$, et se demande quelle transformation relie alors le mannequin à son double reflet.

Enfin, un client debout à hauteur d'yeux constante $\SI{1,6}{m}$ se déplace du point $A(0\,;\,1\,;\,1{,}6)$ vers le point $B(4\,;\,3\,;\,1{,}6)$ en touchant le miroir $P$ en un point $H$ (pour y ajuster une veste exposée contre le mur). Où doit-il toucher le miroir pour que son trajet $AH + HB$ soit le plus court possible ?

\courssub{Consignes}

Travail en groupes de quatre. Chaque groupe rédige les calculs analytiques complets et produit une maquette schématique cotée (vue de dessus dans le plan $z = 0$ et vue en perspective).

\begin{enumerate}
\item \textbf{Reflets successifs.} Déterminer les coordonnées de $M_1$, image de $M$ par la réflexion $s_P$ de plan $P$, puis celles de $M_2$, image de $M_1$ par la réflexion $s_Q$ de plan $Q$.
\item \textbf{Nature de la composée.} Montrer que $s_Q \circ s_P$ est une translation dont on précisera le vecteur $\vec{u}$. Calculer la distance $MM_2$ et interpréter concrètement le résultat pour Monsieur Etoundi (décalage apparent du « reflet du reflet »).
\item \textbf{Miroirs perpendiculaires.} On remplace $Q$ par le plan $R$ d'équation $x = 3$. Déterminer la nature et les éléments caractéristiques de $s_R \circ s_P$, puis les coordonnées de l'image $M'$ de $M$ par cette composée.
\item \textbf{Le trajet minimal.} En utilisant $A_1 = s_P(A)$, justifier que pour tout point $H$ de $P$, on a $AH + HB = A_1H + HB$, et en déduire le point $H$ du miroir $P$ qui minimise le trajet du client, ainsi que la longueur minimale du trajet.
\end{enumerate}

\courssub{Ressources à mobiliser}

\begin{aretenir}[Boîte à outils]
\begin{itemize}
\item \textbf{Réflexion de plan $P$} : $M' = s_P(M)$ si et seulement si $P$ est le plan médiateur de $[MM']$ (la droite $(MM')$ est perpendiculaire à $P$ et le milieu de $[MM']$ appartient à $P$). Pour $P : y = 0$, on a $s_P(x\,;\,y\,;\,z) = (x\,;\,-y\,;\,z)$.
\item \textbf{Expression analytique générale} : Pour un plan $y = b$, $s_{(y=b)}(x\,;\,y\,;\,z) = (x\,;\,2b-y\,;\,z)$. De même pour $x=a$ et $z=c$.
\item \textbf{Composée de deux réflexions de plans parallèles} distants de $d$ : c'est une \cle{translation} de vecteur $2\,\overrightarrow{K_PK_Q}$ (vecteur normal aux plans, de norme $2d$, dirigé du premier plan vers le second).
\item \textbf{Composée de deux réflexions de plans perpendiculaires} : c'est un \cle{demi-tour} d'axe la droite d'intersection des deux plans.
\item \textbf{Demi-tour d'axe $(\Delta)$} : $M'$ est l'image de $M$ si le milieu de $[MM']$ est le projeté orthogonal de $M$ sur $(\Delta)$ et $MM' \perp (\Delta)$. Pour l'axe $(O\,;\vec{k})$, $(x\,;\,y\,;\,z) \mapsto (-x\,;\,-y\,;\,z)$.
\item \textbf{Conservation des distances} : toute isométrie conserve les distances ; en particulier $AH = A_1H$ si $A_1 = s_P(A)$ et $H \in P$.
\item \textbf{Inégalité triangulaire} : $A_1H + HB \geq A_1B$, avec égalité si et seulement si $H \in [A_1B]$.
\end{itemize}
\end{aretenir}

\courssub{Démarche et corrigé commenté}

\begin{exoresolu}[Le showroom aux deux miroirs : corrigé complet]
\textbf{Énoncé rappelé.} $P : y = 0$, $Q : y = 2{,}5$, $R : x = 3$, $M(1\,;\,0{,}8\,;\,1{,}2)$, $A(0\,;\,1\,;\,1{,}6)$, $B(4\,;\,3\,;\,1{,}6)$.

\tcblower

\textbf{Question 1 : les deux reflets.}

Le plan $P$ a pour équation $y = 0$ : la réflexion $s_P$ transforme $(x\,;\,y\,;\,z)$ en $(x\,;\,-y\,;\,z)$. Donc :
\[
M_1 = s_P(M) = (1\,;\,-0{,}8\,;\,1{,}2).
\]
\emph{Vérification} : le milieu de $[MM_1]$ est $(1\,;\,0\,;\,1{,}2)$, qui appartient bien à $P$, et $\overrightarrow{MM_1} = (0\,;\,-1{,}6\,;\,0)$ est bien normal à $P$ (colinéaire à $\vec{\jmath}$).

Pour $s_Q$, le plan $Q : y = 2{,}5$ est parallèle à $P$. On utilise la formule générale $y' = 2 \times 2{,}5 - y = 5 - y$. Ainsi :
\[
s_Q : (x\,;\,y\,;\,z) \longmapsto (x\,;\,5 - y\,;\,z), \qquad M_2 = s_Q(M_1) = (1\,;\,5 - (-0{,}8)\,;\,1{,}2) = (1\,;\,5{,}8\,;\,1{,}2).
\]
\emph{Vérification} : le milieu de $[M_1M_2]$ est $(1\,;\,2{,}5\,;\,1{,}2) \in Q$ et $\overrightarrow{M_1M_2} = (0\,;\,6{,}6\,;\,0) \parallel \vec{\jmath}$.

\textbf{Question 2 : nature de $s_Q \circ s_P$.}

Pour tout point $N(x\,;\,y\,;\,z)$ :
\[
s_Q \circ s_P (x\,;\,y\,;\,z) = s_Q(x\,;\,-y\,;\,z) = (x\,;\,5 - (-y)\,;\,z) = (x\,;\,y + 5\,;\,z).
\]
C'est l'expression analytique de la \cle{translation} de vecteur $\vec{u} = 5\,\vec{\jmath}$, c'est-à-dire $\vec{u}(0\,;\,5\,;\,0)$.

\emph{Cohérence avec le cours} : les plans $P$ et $Q$ sont parallèles, distants de $d = 2{,}5$ m ; la composée est la translation de vecteur $2\,\overrightarrow{K_PK_Q}$ où $\overrightarrow{K_PK_Q} = 2{,}5\,\vec{\jmath}$ (du plan $P$ vers le plan $Q$), donc $\vec{u} = 5\,\vec{\jmath}$. \checkmark

On a $MM_2 = \|\vec{u}\| = \SI{5}{m}$. \textbf{Interprétation} : le « reflet du reflet » apparaît à Monsieur Etoundi décalé de $5$ mètres derrière le mur, soit le double de l'écart entre les deux miroirs : c'est ce qui donne au showroom son impression de profondeur infinie.

\textbf{Question 3 : miroirs perpendiculaires.}

$P : y = 0$ a pour vecteur normal $\vec{n}_P = \vec{\jmath}$ ; $R : x = 3$ a pour vecteur normal $\vec{n}_R = \vec{\imath}$. Comme $\vec{\imath} \cdot \vec{\jmath} = 0$, les vecteurs normaux sont orthogonaux, donc les plans $P$ et $R$ sont \cle{perpendiculaires}. Leur intersection est la droite $(\Delta)$ d'équations cartésiennes $\begin{cases} y = 0 \\ x = 3 \end{cases}$, c'est-à-dire la droite passant par $\Omega(3\,;\,0\,;\,0)$ et dirigée par $\vec{k}$ (droite verticale, l'arête entre le mur et la cloison).

D'après le cours, $s_R \circ s_P$ est le \cle{demi-tour} d'axe $(\Delta)$.

Pour trouver l'image de $M$, on compose les deux réflexions :
\begin{itemize}
\item $s_P(M) = (1\,;\,-0{,}8\,;\,1{,}2)$.
\item Pour $s_R$ (plan $x = 3$) : la formule donne $x' = 2 \times 3 - x = 6 - x$, donc $s_R : (x\,;\,y\,;\,z) \mapsto (6 - x\,;\,y\,;\,z)$.
\end{itemize}
Alors :
\[
M' = s_R \circ s_P(M) = s_R(1\,;\,-0{,}8\,;\,1{,}2) = (6 - 1\,;\,-0{,}8\,;\,1{,}2) = (5\,;\,-0{,}8\,;\,1{,}2).
\]
\emph{Vérification} : le milieu de $[MM']$ est $I\left(\frac{1+5}{2}\,;\,\frac{0{,}8-0{,}8}{2}\,;\,\frac{1{,}2+1{,}2}{2}\right) = (3\,;\,0\,;\,1{,}2)$, qui appartient bien à $(\Delta)$ ($x=3, y=0$), et $\overrightarrow{MM'} = (4\,;\,-1{,}6\,;\,0)$ est orthogonal à $\vec{k}$ (directeur de $\Delta$). Le double reflet est donc le symétrique du mannequin par rapport à l'arête verticale des deux miroirs.

\textbf{Question 4 : le trajet minimal.}

Soit $A_1 = s_P(A) = (0\,;\,-1\,;\,1{,}6)$. Pour tout point $H$ du plan $P$, la réflexion conserve les distances et $H$ est invariant (car $H \in P$), donc :
\[
AH = A_1H \quad \Longrightarrow \quad AH + HB = A_1H + HB.
\]
D'après l'inégalité triangulaire dans le triangle $A_1HB$ : $A_1H + HB \geq A_1B$, avec égalité si et seulement si $H$ appartient au segment $[A_1B]$. Le trajet est minimal lorsque $H$ est l'intersection de la droite $(A_1B)$ avec le plan $P$ ($y = 0$).

Paramétrons la droite $(A_1B)$ : $\overrightarrow{A_1B} = (4\,;\,4\,;\,0)$.
$H = A_1 + t\,\overrightarrow{A_1B} = (0\,;\,-1\,;\,1{,}6) + t(4\,;\,4\,;\,0) = (4t\,;\,-1 + 4t\,;\,1{,}6)$, \quad $t \in \mathbb{R}$.
La condition $H \in P \iff y_H = 0$ donne $-1 + 4t = 0 \iff t = \dfrac{1}{4}$.
D'où :
\[
H\left(1\,;\,0\,;\,1{,}6\right).
\]
La longueur minimale du trajet est :
\[
A_1B = \sqrt{4^2 + 4^2 + 0^2} = \sqrt{32} = 4\sqrt{2} \approx \SI{5,66}{m}.
\]
Le client doit donc toucher le miroir au point situé à $\SI{1}{m}$ de l'axe $(Oy)$, à hauteur de ses yeux ($\SI{1,6}{m}$) : c'est le point où le rayon « réfléchi » de son trajet respecte la loi de la réflexion (angle d'incidence = angle de réflexion), exactement comme en optique géométrique.
\end{exoresolu}

\begin{popfigure}
\begin{center}
\begin{tikzpicture}[scale=0.9,>=Stealth]
% Axes et grille
\draw[popmuted, very thin] (-1,-1.5) grid (6.5,6.5);
\draw[->,popdark] (-1,0) -- (6.5,0) node[below right]{$x$};
\draw[->,popdark] (0,-1.5) -- (0,6.5) node[above left]{$y$};

% Miroirs
\draw[very thick,popblue] (-1,0) -- (6.5,0) node[pos=0.9,below,popblue]{$P : y=0$};
\draw[very thick,poporange] (-1,2.5) -- (6.5,2.5) node[pos=0.9,above,poporange]{$Q : y=2,5$};
\draw[very thick,popgreen] (3,-1.5) -- (3,6.5) node[pos=0.9,right,popgreen]{$R : x=3$};

% Points principaux
\fill[popink] (1,0.8) circle (2.5pt) node[above right]{$M(1;0,8)$};
\fill[poppurple] (1,-0.8) circle (2.5pt) node[below right]{$M_1$};
\fill[poporange] (1,5.8) circle (2.5pt) node[above right]{$M_2$};
\fill[popgreen] (5,-0.8) circle (2.5pt) node[below right]{$M'$};
\fill[popdark] (3,0) circle (2.5pt) node[below left]{$\Omega(3;0)$};

% Flèches réflexions parallèles
\draw[->,dashed,poppurple,thick] (1,0.8) -- (1,-0.8) node[midway,right]{$s_P$};
\draw[->,dashed,poporange,thick] (1,-0.8) -- (1,5.8) node[midway,right]{$s_Q$};
\draw[->,thick,popink] (1.3,0.9) -- (1.3,5.7) node[midway,right]{$\vec{u}=5\vec{\jmath}$};

% Flèche demi-tour
\draw[->,dashed,popgreen,thick] (1,0.8) -- (5,-0.8) node[midway,above left]{$s_R \circ s_P$};

% Trajet minimal (Question 4)
\fill[popblue] (0,1) circle (2pt) node[left]{$A$};
\fill[popblue] (4,3) circle (2pt) node[above right]{$B$};
\fill[poppurple] (0,-1) circle (2pt) node[left]{$A_1$};
\fill[popblue] (1,0) circle (3pt) node[below right,popblue]{$H$};
\draw[thick,popblue] (0,-1) -- (4,3) node[midway,above left]{$A_1B$};
\draw[thick,popblue,dashed] (0,1) -- (1,0);
\draw[thick,popblue,dashed] (1,0) -- (4,3);

% Cotes
\draw[|-|,popmuted] (3.2, -1.2) -- (5.2, -1.2) node[midway,below]{$2$ m};
\draw[|-|,popmuted] (-1.2, 0.2) -- (-1.2, 2.3) node[midway,left]{$2,5$ m};
\draw[|-|,popmuted] (-1.2, 2.7) -- (-1.2, 5.6) node[midway,left]{$3,3$ m};
\end{tikzpicture}
\end{center}
\legende{Vue de dessus (plan horizontal $z=1,2$ pour le mannequin, $z=1,6$ pour le client). En bleu : miroirs $P, R$ et trajet minimal $A \to H \to B$ via $A_1$. En orange : miroirs $P, Q$ et reflets $M_1, M_2$ (translation). En vert : demi-tour d'axe $\Delta$ (intersection $P \cap R$).}
\end{popfigure}

\courssub{Bilan de l'activité}

Cette activité a permis de mettre en évidence, sur une situation concrète d'aménagement commercial, les résultats centraux de la leçon :

\begin{itemize}
\item une \cle{réflexion} se manipule très simplement en coordonnées dès lors que le plan est parallèle à un plan de coordonnées : une seule coordonnée change, par symétrie autour de la cote du plan ($y' = 2b - y$ pour $y=b$) ;
\item la composée de deux réflexions de plans \cle{parallèles} distants de $d$ est une \cle{translation} de vecteur de norme $2d$, perpendiculaire aux plans, dirigé du premier plan vers le second : c'est le phénomène des « reflets en cascade » des miroirs face à face, qui multiplie visuellement l'espace du showroom ;
\item la composée de deux réflexions de plans \cle{perpendiculaires} est un \cle{demi-tour} dont l'axe est la droite d'intersection des deux plans : l'arête verticale joue le rôle de « charnière » du double reflet ;
\item les isométries sont un outil puissant d'\cle{optimisation} : en « dépliant » le trajet par réflexion, le problème du plus court chemin via un miroir se ramène à une simple ligne droite (inégalité triangulaire), exactement comme la lumière choisit son chemin en optique (principe de Fermat / Héron d'Alexandrie).
\end{itemize}

\begin{savaistu}[Des miroirs aux périscopes]
Le principe des deux miroirs parallèles ou inclinés est à la base de nombreux objets techniques : le périscope (sous-marins, chars), les salons de coiffure de nos quartiers où deux miroirs face à face créent une « infinité » de reflets, et même les réflecteurs routiers « yeux de chat ». Les géomètres grecs (Héron d'Alexandrie, I\textsuperscript{er} siècle) utilisaient déjà la réflexion pour prouver que la lumière suit le plus court chemin entre deux points en se réfléchissant sur un plan : c'est exactement la méthode de la question 4, fondement de l'optique géométrique !
\end{savaistu}

\newpage

\coursec{Méthodes et exercices résolus}

\coursec{Réflexion de l'espace : définition et propriétés}

\begin{definition}[Réflexion (symétrie orthogonale par rapport à un plan)]
Soit $(P)$ un plan de l'espace. La \textbf{réflexion de plan $(P)$}, notée $s_{P}$, est l'application qui associe à tout point $M$ le point $M'$ tel que :
\begin{itemize}
\item le milieu $I$ du segment $[MM']$ appartient au plan $(P)$ ;
\item la droite $(MM')$ est orthogonale au plan $(P)$ (c'est-à-dire $\overrightarrow{MM'}$ est colinéaire à un vecteur normal $\vec{n}$ de $(P)$).
\end{itemize}
Si $M\in(P)$, alors $M'=M$ : les points de $(P)$ sont \textbf{fixes}.
\end{definition}

\begin{popfigure}
\begin{tikzpicture}[scale=1.2, >=Stealth]
% Plan
\fill[popblue!10, opacity=0.5] (-2,-1.5) -- (2,-1.5) -- (2,1.5) -- (-2,1.5) -- cycle;
\draw[popblue, thick] (-2,-1.5) -- (2,-1.5) -- (2,1.5) -- (-2,1.5) -- cycle;
\node[popblue, right] at (2,1.5) {$(P)$};
% Normal vector
\draw[poporange, thick, ->] (0,0) -- (0,1.5) node[midway, right] {$\vec{n}$};
% Points M and M'
\coordinate (M) at (-1.5, -1);
\coordinate (Mp) at (1.5, 1);
\coordinate (I) at (0,0);
\fill[popdark] (M) circle (2pt) node[left] {$M$};
\fill[popdark] (Mp) circle (2pt) node[right] {$M'$};
\fill[popgreen] (I) circle (2pt) node[below left] {$I$};
\draw[popmuted, dashed] (M) -- (Mp);
\draw[popgreen, <->] (M) -- (I) node[midway, left] {$MI$};
\draw[popgreen, <->] (I) -- (Mp) node[midway, right] {$IM'$};
% Right angle mark
\draw[popdark] (-0.15, -0.15) -- (0.15, -0.15) -- (0.15, 0.15);
\end{tikzpicture}
\legende{La réflexion $s_{P}$ : $I$ milieu de $[MM']$ et $(MM')\perp(P)$.}
\end{popfigure}

\begin{propriete}[Propriétés fondamentales de la réflexion]
\begin{enumerate}
\item $s_{P}$ est une \textbf{isométrie} : elle conserve les distances, les angles, les aires, les volumes.
\item $s_{P}$ est \textbf{involutive} : $s_{P}\circ s_{P} = \mathrm{Id}$ (appliquer deux fois la réflexion revient à ne rien faire).
\item L'ensemble des points fixes de $s_{P}$ est exactement le plan $(P)$.
\item $s_{P}$ transforme une droite en une droite, un plan en un plan, une sphère en une sphère.
\item Pour tout point $M$, $(P)$ est le \textbf{plan médiateur} du segment $[M\,s_{P}(M)]$.
\end{enumerate}
\end{propriete}

\coursec{Expression analytique d'une réflexion}

\begin{methode}[Expression analytique de la réflexion $s_{P}$ de plan $(P)$]
Dans un repère orthonormé $(O\,;\vec{\imath},\vec{\jmath},\vec{k})$, soit $(P)$ d'équation cartésienne $ax+by+cz+d=0$ ($\vec{n}(a\,;b\,;c)$ normal à $(P)$). Pour déterminer $s_{P}(M)=M'(x'\,;y'\,;z')$ à partir de $M(x\,;y\,;z)$ :
\begin{enumerate}
\item \textbf{Colinéarité :} $\overrightarrow{MM'}\parallel\vec{n}$ $\Rightarrow$ $\exists\lambda\in\mathbb{R}$,\quad $x'=x+\lambda a$, $y'=y+\lambda b$, $z'=z+\lambda c$.
\item \textbf{Milieu dans $(P)$ :} $I\bigl(\frac{x+x'}{2},\frac{y+y'}{2},\frac{z+z'}{2}\bigr)\in(P)$ $\Rightarrow$ $a\frac{x+x'}{2}+b\frac{y+y'}{2}+c\frac{z+z'}{2}+d=0$.
\item \textbf{Substitution :} remplacer $x',y',z'$ dans l'équation du milieu $\Rightarrow$ $\lambda=-\dfrac{2(ax+by+cz+d)}{a^{2}+b^{2}+c^{2}}$.
\item \textbf{Conclusion :} reporter $\lambda$ dans les expressions de $x',y',z'$.
\end{enumerate}
\textbf{Cas particuliers (plans parallèles aux plans de coordonnées) — à connaître par cœur :}
\begin{itemize}
\item $(P):z=c$ (parallèle à $(xOy)$) : $\begin{cases}x'=x\\ y'=y\\ z'=2c-z\end{cases}$
\item $(P):y=b$ (parallèle à $(xOz)$) : $\begin{cases}x'=x\\ y'=2b-y\\ z'=z\end{cases}$
\item $(P):x=a$ (parallèle à $(yOz)$) : $\begin{cases}x'=2a-x\\ y'=y\\ z'=z\end{cases}$
\end{itemize}
\end{methode}

\begin{exoresolu}[Réflexion de plan oblique — type Bac]
L'espace est muni d'un repère orthonormé $(O\,;\vec{\imath},\vec{\jmath},\vec{k})$. On considère le plan $(P)$ d'équation $x-2y+2z-3=0$ et le point $A(1\,;0\,;-1)$.
\begin{enumerate}
\item Déterminer l'expression analytique de la réflexion $s$ de plan $(P)$.
\item Déterminer les coordonnées de $A'$, image de $A$ par $s$. Vérifier que le milieu de $[AA']$ appartient à $(P)$.
\end{enumerate}
\tcblower
\textbf{1. Expression analytique de $s$.}

Le plan $(P)$ a pour vecteur normal $\vec{n}(1\,;-2\,;2)$ ; on a $a^{2}+b^{2}+c^{2}=1+4+4=9$.
Soit $M(x\,;y\,;z)$ un point de l'espace et $M'(x'\,;y'\,;z')=s(M)$.

\textbf{Condition 1 (colinéarité) :} $\overrightarrow{MM'}$ est colinéaire à $\vec{n}$, donc $\exists\lambda\in\mathbb{R}$ tel que :
\[x'=x+\lambda,\qquad y'=y-2\lambda,\qquad z'=z+2\lambda.\]

\textbf{Condition 2 (milieu dans $(P)$) :} le milieu $I\left(\dfrac{x+x'}{2}\,;\dfrac{y+y'}{2}\,;\dfrac{z+z'}{2}\right)$ de $[MM']$ appartient à $(P)$ :
\[\frac{x+x'}{2}-2\cdot\frac{y+y'}{2}+2\cdot\frac{z+z'}{2}-3=0.\]
En remplaçant $x',y',z'$ par leurs expressions en fonction de $\lambda$ :
\begin{align*}
\frac{2x+\lambda}{2}-2\cdot\frac{2y-2\lambda}{2}+2\cdot\frac{2z+2\lambda}{2}-3&=0\\
x+\tfrac{\lambda}{2}-2y+2\lambda+2z+2\lambda-3&=0\\
(x-2y+2z-3)+\tfrac{9}{2}\lambda&=0.
\end{align*}
On en tire $\lambda=-\dfrac{2(x-2y+2z-3)}{9}$.

En reportant dans les expressions de $x',y',z'$ :
\[\boxed{\begin{cases}
x'=x-\dfrac{2}{9}(x-2y+2z-3)\\[4pt]
y'=y+\dfrac{4}{9}(x-2y+2z-3)\\[4pt]
z'=z-\dfrac{4}{9}(x-2y+2z-3)
\end{cases}}\]
Après développement et mise au même dénominateur :
\[\begin{cases}
x'=\dfrac{7x+4y-4z+6}{9}\\[4pt]
y'=\dfrac{4x+y+8z-12}{9}\\[4pt]
z'=\dfrac{-4x+8y+z+12}{9}.
\end{cases}\]

\textbf{2. Image de $A(1\,;0\,;-1)$.}

On calcule $x_{A}-2y_{A}+2z_{A}-3=1-0-2-3=-4$, d'où $\lambda=-\dfrac{2\times(-4)}{9}=\dfrac{8}{9}$.
Alors :
\[x'=1+\frac{8}{9}=\frac{17}{9},\qquad y'=0-2\times\frac{8}{9}=-\frac{16}{9},\qquad z'=-1+2\times\frac{8}{9}=\frac{7}{9}.\]
Donc $A'\left(\dfrac{17}{9}\,;-\dfrac{16}{9}\,;\dfrac{7}{9}\right)$.

\textbf{Vérification :} le milieu de $[AA']$ est $I\left(\dfrac{13}{9}\,;-\dfrac{8}{9}\,;-\dfrac{1}{9}\right)$. On a :
\[\frac{13}{9}-2\times\left(-\frac{8}{9}\right)+2\times\left(-\frac{1}{9}\right)-3=\frac{13+16-2-27}{9}=0.\]
Donc $I\in(P)$. De plus $\overrightarrow{AA'}\left(\tfrac{8}{9}\,;-\tfrac{16}{9}\,;\tfrac{16}{9}\right)=\tfrac{8}{9}\vec{n}$, donc $(AA')\perp(P)$. La réflexion est bien vérifiée.
\end{exoresolu}

\coursec{Composée de deux réflexions}

\begin{popfigure}
\begin{tikzpicture}[scale=1.1, >=Stealth]
% Two parallel planes
\fill[popblue!10, opacity=0.3] (-2,-1.5) -- (2,-1.5) -- (2,-0.5) -- (-2,-0.5) -- cycle;
\fill[popgreen!10, opacity=0.3] (-2,0.5) -- (2,0.5) -- (2,1.5) -- (-2,1.5) -- cycle;
\draw[popblue, thick] (-2,-1.5) -- (2,-1.5) -- (2,-0.5) -- (-2,-0.5) -- cycle;
\draw[popgreen, thick] (-2,0.5) -- (2,0.5) -- (2,1.5) -- (-2,1.5) -- cycle;
\node[popblue, left] at (-2,-1) {$(P)$};
\node[popgreen, left] at (-2,1) {$(P')$};
% Points
\coordinate (M) at (-1.5, -1.5);
\coordinate (M1) at (-1.5, 0.5);
\coordinate (Mp) at (-1.5, 2.5);
\fill[popdark] (M) circle (2pt) node[left] {$M$};
\fill[poporange] (M1) circle (2pt) node[left] {$M_1=s_P(M)$};
\fill[poppurple] (Mp) circle (2pt) node[left] {$M'=s_{P'}(M_1)$};
% Arrows
\draw[popblue, ->, thick] (M) -- (M1) node[midway, right] {$s_P$};
\draw[popgreen, ->, thick] (M1) -- (Mp) node[midway, right] {$s_{P'}$};
\draw[poppurple, <->, thick, dashed] (M) -- (Mp) node[midway, left] {$\vec{v}=2\overrightarrow{HH'}$};
% Distance marks
\draw[popmuted, <->] (2.2, -1) -- (2.2, 1) node[midway, right] {$2d$};
\draw[popmuted, <->] (2.5, -1) -- (2.5, 0) node[midway, right] {$d$};
\draw[popmuted, <->] (2.5, 0) -- (2.5, 1) node[midway, right] {$d$};
\end{tikzpicture}
\legende{Composée de deux réflexions de plans parallèles : translation de vecteur $2\overrightarrow{HH'}$.}
\end{popfigure}

\begin{methode}[Nature de la composée $f=s_{P'}\circ s_{P}$]
Soit $(P)$ et $(P')$ deux plans, $s_{P}$ et $s_{P'}$ les réflexions associées.
\begin{enumerate}
\item \textbf{Comparer les plans} : parallèles (distincts), sécants, ou confondus.
\item \textbf{Cas $(P)\parallel(P')$, distincts :} $f$ est une \cle{translation} de vecteur $\vec{v}=2\,\overrightarrow{H H'}$, où $H\in(P)$ et $H'$ est son projeté orthogonal sur $(P')$. Le vecteur $\vec{v}$ est normal aux deux plans, dirigé de $(P)$ vers $(P')$.
\item \textbf{Cas $(P)$ et $(P')$ sécants selon une droite $(\Delta)$ :} $f$ est une \textbf{rotation} d'axe $(\Delta)$. \textbf{En particulier, si $(P)\perp(P')$, alors $f$ est le \cle{demi-tour} d'axe $(\Delta)=(P)\cap(P')$.}
\item \textbf{Cas $(P)=(P')$ :} $f=\mathrm{Id}$ (une réflexion est involutive).
\end{enumerate}
\textbf{Réflexes indispensables :}
\begin{itemize}
\item L'ordre compte : $s_{P'}\circ s_{P}$ translate de $(P)$ vers $(P')$ ; $s_{P}\circ s_{P'}$ translate du vecteur opposé.
\item Pour prouver qu'une composée est une translation, on peut calculer les expressions analytiques et montrer que $x'=x+\alpha$, $y'=y+\beta$, $z'=z+\gamma$.
\end{itemize}
\end{methode}

\begin{exoresolu}[Composée de deux réflexions de plans parallèles]
L'espace est muni d'un repère orthonormé $(O\,;\vec{\imath},\vec{\jmath},\vec{k})$. On considère les plans $(P):z=1$ et $(P'):z=4$, et les réflexions $s_{P}$ et $s_{P'}$ associées.
\begin{enumerate}
\item Déterminer les expressions analytiques de $s_{P}$ et $s_{P'}$.
\item En déduire l'expression analytique de $f=s_{P'}\circ s_{P}$ et donner sa nature et ses éléments caractéristiques.
\item Un chantier de construction à Yaoundé utilise deux miroirs plans parallèles horizontaux distants de \SI{3}{m}. Quel déplacement subit le rayon lumineux réfléchi successivement par les deux miroirs (en termes de position de l'image) ?
\end{enumerate}
\tcblower
\textbf{1. Expressions analytiques.}

$(P)$ et $(P')$ sont parallèles au plan $(xOy)$. D'après les formules du cours (réflexion par rapport au plan $z=c$ : $x'=x$, $y'=y$, $z'=2c-z$) :
\[s_{P}:\begin{cases}x'=x\\ y'=y\\ z'=2-z\end{cases}\qquad\qquad s_{P'}:\begin{cases}x'=x\\ y'=y\\ z'=8-z\end{cases}\]

\textbf{2. Expression et nature de $f=s_{P'}\circ s_{P}$.}

Soit $M(x\,;y\,;z)$, $M_{1}(x_{1}\,;y_{1}\,;z_{1})=s_{P}(M)$ et $M'(x'\,;y'\,;z')=s_{P'}(M_{1})$. Alors :
\[x_{1}=x,\quad y_{1}=y,\quad z_{1}=2-z\ ;\qquad x'=x_{1},\quad y'=y_{1},\quad z'=8-z_{1}=8-(2-z)=6+z.\]
Donc :
\[f:\begin{cases}x'=x\\ y'=y\\ z'=z+6\end{cases}\]
On reconnaît la \textbf{translation de vecteur $\vec{v}(0\,;0\,;6)$}, c'est-à-dire $\vec{v}=6\vec{k}$.

\textbf{Vérification géométrique :} les plans sont distants de $d=4-1=3$, et le vecteur de la translation doit être $2d\,\vec{k}=6\vec{k}$, dirigé de $(P)$ vers $(P')$ (sens des $z$ croissants). C'est cohérent.

\textbf{3. Application concrète.}

L'image d'un objet par les deux réflexions successives subit une translation de vecteur normal aux miroirs, de norme $2\times 3=\SI{6}{m}$, dirigée du premier miroir vers le second. C'est le principe du décalage d'image dans les systèmes à double miroir (périscopes, dispositifs de visée).
\end{exoresolu}

\begin{exoresolu}[Composée de deux réflexions de plans perpendiculaires — Demi-tour]
Soit $(P):x=0$ (plan $(yOz)$) et $(P'):y=0$ (plan $(xOz)$). On note $s_{P}$ et $s_{P'}$ les réflexions associées.
\begin{enumerate}
\item Justifier que $f=s_{P'}\circ s_{P}$ est un demi-tour dont on précisera l'axe.
\item Retrouver ce résultat par le calcul analytique.
\end{enumerate}
\tcblower
\textbf{1. Étude géométrique.}

Les plans $(P)$ et $(P')$ ont pour vecteurs normaux respectifs $\vec{n}(1\,;0\,;0)$ et $\vec{n}'(0\,;1\,;0)$. Or $\vec{n}\cdot\vec{n}'=0$, donc $(P)\perp(P')$.

Leur intersection est la droite $(\Delta)$ d'équations $\begin{cases}x=0\\ y=0\end{cases}$, c'est-à-dire l'axe $(O\,;\vec{k})$ (axe des cotes).

D'après le théorème du cours, la composée de deux réflexions de plans perpendiculaires est le \textbf{demi-tour d'axe $(\Delta)=(P)\cap(P')$}. Donc $f$ est le demi-tour d'axe $(Oz)$.

\textbf{2. Vérification analytique.}

Réflexion par rapport à $(P):x=0$ : $s_{P}:(x\,;y\,;z)\mapsto(-x\,;y\,;z)$.

Réflexion par rapport à $(P'):y=0$ : $s_{P'}:(x\,;y\,;z)\mapsto(x\,;-y\,;z)$.

Composée : $M(x\,;y\,;z)\xrightarrow{s_{P}}M_{1}(-x\,;y\,;z)\xrightarrow{s_{P'}}M'(-x\,;-y\,;z)$.

Donc $f:\begin{cases}x'=-x\\ y'=-y\\ z'=z\end{cases}$

On reconnaît l'expression analytique du demi-tour d'axe $(Oz)$ (la cote est invariante, l'abscisse et l'ordonnée sont changées en leurs opposées). Le calcul confirme le résultat géométrique.
\end{exoresolu}

\coursec{Demi-tour : définition et expression analytique}

\begin{definition}[Demi-tour (symétrie orthogonale par rapport à une droite)]
Soit $(\Delta)$ une droite de l'espace. Le \textbf{demi-tour d'axe $(\Delta)$}, noté $s_{\Delta}$, est l'application qui associe à tout point $M$ le point $M'$ tel que :
\begin{itemize}
\item le milieu $I$ du segment $[MM']$ appartient à la droite $(\Delta)$ ;
\item la droite $(MM')$ est orthogonale à $(\Delta)$ (c'est-à-dire $\overrightarrow{MM'}\cdot\vec{u}=0$ où $\vec{u}$ est un vecteur directeur de $(\Delta)$).
\end{itemize}
Si $M\in(\Delta)$, alors $M'=M$ : les points de $(\Delta)$ sont \textbf{fixes}.
Le demi-tour est l'isométrie qui « retourne » l'espace autour de l'axe $(\Delta)$ d'un angle de $\pi$ (180°).
\end{definition}

\begin{popfigure}
\begin{tikzpicture}[scale=1.2, >=Stealth]
% Axis Delta
\draw[poppurple, thick, ->] (-2,0) -- (2,0) node[right] {$(\Delta)$};
% Point M and M'
\coordinate (M) at (-1, 1.5);
\coordinate (Mp) at (1, -1.5);
\coordinate (I) at (0,0);
\fill[popdark] (M) circle (2pt) node[above left] {$M$};
\fill[popdark] (Mp) circle (2pt) node[below right] {$M'$};
\fill[popgreen] (I) circle (2pt) node[below left] {$I$};
\draw[popmuted, dashed] (M) -- (Mp);
% Perpendicular mark
\draw[popdark] (-0.15, -0.15) -- (0.15, -0.15) -- (0.15, 0.15);
% Projection lines
\draw[poporange, dashed] (M) -- (-1,0);
\draw[poporange, dashed] (Mp) -- (1,0);
\node[poporange, above] at (-1,0) {$H$};
\node[poporange, below] at (1,0) {$H'$};
\end{tikzpicture}
\legende{Le demi-tour $s_{\Delta}$ : $I$ milieu de $[MM']$ sur $(\Delta)$ et $(MM')\perp(\Delta)$.}
\end{popfigure}

\begin{methode}[Expression analytique du demi-tour $s_{\Delta}$ d'axe $(\Delta)$]
Dans un repère orthonormé, soit $(\Delta)$ passant par $A(x_{0}\,;y_{0}\,;z_{0})$ et de vecteur directeur $\vec{u}(a\,;b\,;c)$. Pour déterminer $s_{\Delta}(M)=M'(x'\,;y'\,;z')$ à partir de $M(x\,;y\,;z)$ :
\begin{enumerate}
\item \textbf{Projeté orthogonal $I$ de $M$ sur $(\Delta)$ :} $I$ est le point de $(\Delta)$ tel que $\overrightarrow{MI}\perp\vec{u}$.
On paramètre $(\Delta)$ : $I = A + t\vec{u}$. La condition $\overrightarrow{MI}\cdot\vec{u}=0$ donne :
\[t = \frac{\overrightarrow{AM}\cdot\vec{u}}{\|\vec{u}\|^{2}}.\]
\item \textbf{Milieu :} $I$ est le milieu de $[MM']$, donc $M' = 2I - M$, c'est-à-dire
\[x' = 2x_{I}-x,\quad y' = 2y_{I}-y,\quad z' = 2z_{I}-z.\]
\end{enumerate}
\textbf{Cas particuliers (axes parallèles aux axes de coordonnées, passant par $A(x_{0}\,;y_{0}\,;z_{0})$) — à connaître :}
\begin{itemize}
\item Axe $\parallel (Oz)$ : $\begin{cases}x'=2x_{0}-x\\ y'=2y_{0}-y\\ z'=z\end{cases}$
\item Axe $\parallel (Oy)$ : $\begin{cases}x'=2x_{0}-x\\ y'=y\\ z'=2z_{0}-z\end{cases}$
\item Axe $\parallel (Ox)$ : $\begin{cases}x'=x\\ y'=2y_{0}-y\\ z'=2z_{0}-z\end{cases}$
\end{itemize}
\textbf{Réflexe :} deux coordonnées « basculent » (celles perpendiculaires à l'axe), celle le long de l'axe est conservée.
\end{methode}

\begin{exoresolu}[Demi-tour d'axe donné — type Bac]
L'espace est muni d'un repère orthonormé $(O\,;\vec{\imath},\vec{\jmath},\vec{k})$. Soit $(\Delta)$ la droite passant par $A(1\,;2\,;0)$ et de vecteur directeur $\vec{u}(0\,;0\,;1)$.
\begin{enumerate}
\item Déterminer l'expression analytique du demi-tour $s_{\Delta}$ d'axe $(\Delta)$.
\item Soit $B(3\,;-1\,;5)$. Déterminer $B'=s_{\Delta}(B)$ et vérifier que $B$ et $B'$ ont même projeté orthogonal sur $(\Delta)$.
\end{enumerate}
\tcblower
\textbf{1. Expression analytique.}

Soit $M(x\,;y\,;z)$ et $M'(x'\,;y'\,;z')=s_{\Delta}(M)$. Notons $I$ le milieu de $[MM']$. Par définition du demi-tour, $I$ est le projeté orthogonal de $M$ sur $(\Delta)$.

Comme $(\Delta)$ est dirigée par $\vec{u}(0\,;0\,;1)$, elle est parallèle à l'axe $(Oz)$. Le projeté orthogonal de $M(x\,;y\,;z)$ sur $(\Delta)$ s'obtient en conservant la cote et en prenant l'abscisse et l'ordonnée de $A$ : $I(1\,;2\,;z)$.

Justifions : $I\in(\Delta)$ car $\overrightarrow{AI}=(0\,;0\,;z)=z\,\vec{u}$, et $\overrightarrow{MI}=(-x+1\,;\,2-y\,;\,0)$ vérifie $\overrightarrow{MI}\cdot\vec{u}=0$. Donc $I$ est bien le projeté orthogonal de $M$ sur $(\Delta)$.

La relation $I$ milieu de $[MM']$ donne $M'=2I-M$ :
\[\boxed{\begin{cases}
x'=2-x\\
y'=4-y\\
z'=z
\end{cases}}\]

\textbf{2. Image de $B(3\,;-1\,;5)$.}

$x'=2-3=-1$ ; $y'=4-(-1)=5$ ; $z'=5$. Donc $B'(-1\,;5\,;5)$.

\textbf{Vérification :} le milieu de $[BB']$ est $I\left(\dfrac{3-1}{2}\,;\dfrac{-1+5}{2}\,;\dfrac{5+5}{2}\right)=(1\,;2\,;5)$. Or $I\in(\Delta)$ car $\overrightarrow{AI}=(0\,;0\,;5)=5\vec{u}$, et $\overrightarrow{BB'}=(-4\,;6\,;0)$ vérifie $\overrightarrow{BB'}\cdot\vec{u}=0$, donc $(BB')\perp(\Delta)$. Le point $I$ est bien le projeté orthogonal commun de $B$ et $B'$ sur $(\Delta)$ : $s_{\Delta}(B)=B'$.
\end{exoresolu}

\coursec{Composée de deux demi-tours}

\begin{methode}[Nature de la composée $g=s_{\Delta'}\circ s_{\Delta}$]
Soit $(\Delta)$ et $(\Delta')$ deux droites, $s_{\Delta}$ et $s_{\Delta'}$ les demi-tours d'axes respectifs.
\begin{enumerate}
\item \textbf{Déterminer la position relative des axes} : parallèles, sécants, ou non coplanaires.
\item \textbf{Cas $(\Delta)\parallel(\Delta')$, distincts :} $g$ est une \cle{translation} de vecteur $\vec{v}=2\,\overrightarrow{KK'}$, où $K\in(\Delta)$ et $K'$ est son projeté orthogonal sur $(\Delta')$ (dans le plan contenant les deux droites). Le vecteur $\vec{v}$ est perpendiculaire aux deux axes.
\item \textbf{Cas $(\Delta)$ et $(\Delta')$ sécants et perpendiculaires en un point $O$ :} $g$ est le \cle{demi-tour} dont l'axe $(D)$ est la droite passant par $O$ et perpendiculaire au plan $(\Delta,\Delta')$.
\item \textbf{Cas $(\Delta)=(\Delta')$ :} $g=\mathrm{Id}$.
\end{enumerate}
\textbf{Réflexe calculatoire :} si les expressions analytiques sont simples (axes parallèles aux axes du repère), composer directement les formules et identifier le résultat :
\begin{itemize}
\item forme $(x+\alpha\,;\,y+\beta\,;\,z+\gamma)$ $\Rightarrow$ translation de vecteur $(\alpha\,;\beta\,;\gamma)$ ;
\item forme $(2a-x\,;\,2b-y\,;\,z)$ $\Rightarrow$ demi-tour d'axe parallèle à $(Oz)$ passant par $(a\,;b\,;0)$.
\end{itemize}
\end{methode}

\begin{exoresolu}[Composée de deux demi-tours d'axes parallèles]
L'espace est muni d'un repère orthonormé $(O\,;\vec{\imath},\vec{\jmath},\vec{k})$. Soit $(\Delta)$ la droite passant par $A(1\,;0\,;0)$ dirigée par $\vec{k}$, et $(\Delta')$ la droite passant par $A'(4\,;0\,;0)$ dirigée par $\vec{k}$.
\begin{enumerate}
\item Donner les expressions analytiques de $s_{\Delta}$ et $s_{\Delta'}$.
\item Déterminer l'expression analytique de $g=s_{\Delta'}\circ s_{\Delta}$, puis sa nature et ses éléments caractéristiques.
\item Comparer $s_{\Delta'}\circ s_{\Delta}$ et $s_{\Delta}\circ s_{\Delta'}$.
\end{enumerate}
\tcblower
\textbf{1. Expressions analytiques.}

Les deux axes sont parallèles à $(Oz)$. D'après les formules du cours (demi-tour d'axe parallèle à $(Oz)$ passant par $(x_{0}\,;y_{0}\,;z_{0})$ : $x'=2x_{0}-x$, $y'=2y_{0}-y$, $z'=z$) :
\[s_{\Delta}:\begin{cases}x'=2-x\\ y'=-y\\ z'=z\end{cases}\qquad\qquad s_{\Delta'}:\begin{cases}x'=8-x\\ y'=-y\\ z'=z\end{cases}\]

\textbf{2. Expression et nature de $g=s_{\Delta'}\circ s_{\Delta}$.}

Soit $M(x\,;y\,;z)$, $M_{1}=s_{\Delta}(M)$ de coordonnées $(2-x\,;-y\,;z)$, puis $M'=s_{\Delta'}(M_{1})$ :
\[x'=8-(2-x)=x+6,\qquad y'=-(-y)=y,\qquad z'=z.\]
Donc $g:\begin{cases}x'=x+6\\ y'=y\\ z'=z\end{cases}$ : c'est la \textbf{translation de vecteur $\vec{v}=6\vec{\imath}$}.

\textbf{Vérification géométrique :} les axes sont parallèles et distants de $AA'=3$ (avec $\overrightarrow{AA'}=3\vec{\imath}$ perpendiculaire aux axes). La composée est la translation de vecteur $2\overrightarrow{AA'}=6\vec{\imath}$. Cohérent.

\textbf{3. Comparaison.}

Par le même calcul, $s_{\Delta}\circ s_{\Delta'}$ est la translation de vecteur $-6\vec{\imath}$. Les deux composées sont donc \textbf{distinctes} : la composition des isométries n'est pas commutative en général. Ce sont deux translations réciproques l'une de l'autre.
\end{exoresolu}

\coursec{Reconnaissance d'isométries et résolution de problèmes}

\begin{methode}[Identifier la nature d'une isométrie à partir de son expression analytique]
On donne $f:\begin{cases}x'=ax+by+cz+\alpha\\ y'=a'x+b'y+c'z+\beta\\ z'=a''x+b''y+c''z+\gamma\end{cases}$. Démarche rigoureuse :
\begin{enumerate}
\item \textbf{Chercher les points fixes} : résoudre le système $f(M)=M$, c'est-à-dire $x'=x$, $y'=y$, $z'=z$.
\begin{itemize}
\item Un \textbf{plan} de points fixes $\Rightarrow$ réflexion (dont le plan est celui des points fixes).
\item Une \textbf{droite} de points fixes $\Rightarrow$ demi-tour (ou rotation) d'axe cette droite.
\item \textbf{Aucun} point fixe $\Rightarrow$ translation (si l'expression est de la forme $x'=x+\alpha$, etc.) ou composée plus complexe.
\item Un \textbf{unique} point fixe $\Rightarrow$ symétrie centrale ou rotation.
\end{itemize}
\item \textbf{Identifier la forme canonique} : comparer aux expressions types :
\begin{itemize}
\item $(x+\alpha\,;\,y+\beta\,;\,z+\gamma)$ : translation de vecteur $(\alpha\,;\beta\,;\gamma)$ ;
\item $(2a-x\,;\,y\,;\,z)$ et analogues : réflexion de plan $x=a$ ;
\item $(2a-x\,;\,2b-y\,;\,z)$ et analogues : demi-tour d'axe parallèle à $(Oz)$ passant par $(a\,;b\,;0)$ ;
\item $(2a-x\,;\,2b-y\,;\,2c-z)$ : symétrie centrale de centre $(a\,;b\,;c)$.
\end{itemize}
\item \textbf{Conclure} en donnant la nature \textbf{ET} les éléments caractéristiques (plan, axe, vecteur, centre).
\end{enumerate}
\end{methode}

\begin{exoresolu}[Reconnaissance d'isométries — type Bac]
L'espace est muni d'un repère orthonormé $(O\,;\vec{\imath},\vec{\jmath},\vec{k})$. Pour chacune des applications suivantes, déterminer sa nature et ses éléments caractéristiques :
\[f_{1}:\begin{cases}x'=-x+4\\ y'=y\\ z'=z\end{cases}\qquad f_{2}:\begin{cases}x'=-x+2\\ y'=-y+6\\ z'=z\end{cases}\qquad f_{3}:\begin{cases}x'=x-1\\ y'=y+3\\ z'=z-2\end{cases}\]
\tcblower
\textbf{Étude de $f_{1}$.} Cherchons les points fixes : $x=-x+4\Leftrightarrow x=2$ ; $y=y$ et $z=z$ toujours vrais. L'ensemble des points fixes est le plan $(P):x=2$. De plus, $f_{1}$ laisse $y$ et $z$ inchangés et envoie $x$ sur $4-x$ : le segment $[MM']$ est dirigé selon $\vec{\imath}$, normal à $(P)$, et son milieu a pour abscisse $\dfrac{x+4-x}{2}=2$, donc appartient à $(P)$.

\textbf{Conclusion :} $f_{1}$ est la \textbf{réflexion de plan $(P):x=2$}.

\textbf{Étude de $f_{2}$.} Points fixes : $x=-x+2\Leftrightarrow x=1$ ; $y=-y+6\Leftrightarrow y=3$ ; $z=z$ toujours vrai. L'ensemble des points fixes est la droite $(\Delta):\begin{cases}x=1\\ y=3\end{cases}$, qui est la droite passant par $(1\,;3\,;0)$ et dirigée par $\vec{k}$. La forme $(2a-x\,;\,2b-y\,;\,z)$ avec $a=1$, $b=3$ est caractéristique.

\textbf{Conclusion :} $f_{2}$ est le \textbf{demi-tour d'axe $(\Delta)$}, droite passant par $A(1\,;3\,;0)$ et de vecteur directeur $\vec{k}$.

\textbf{Étude de $f_{3}$.} Points fixes : $x=x-1$ est impossible. Il n'y a aucun point fixe, et l'expression est de la forme $(x+\alpha\,;\,y+\beta\,;\,z+\gamma)$.

\textbf{Conclusion :} $f_{3}$ est la \textbf{translation de vecteur $\vec{v}(-1\,;3\,;-2)$}.
\end{exoresolu}

\begin{methode}[Résoudre un problème de géométrie de l'espace par les isométries]
\begin{enumerate}
\item \textbf{Repérer la symétrie cachée} dans la configuration : plan médiateur, milieu, droite perpendiculaire, distances égales. Traduire : « $(P)$ est le plan médiateur de $[AB]$ » signifie $s_{P}(A)=B$.
\item \textbf{Choisir l'isométrie adaptée} et écrire précisément ce qu'elle échange.
\item \textbf{Utiliser les propriétés de conservation} : une isométrie conserve les distances, les angles, l'orthogonalité, les milieux, l'alignement, et transforme droites en droites, plans en plans.
\item \textbf{Transporter le problème} : remplacer un point par son image pour se ramener à une configuration plus simple (points alignés, distance facile à calculer).
\item \textbf{Rédiger la conclusion} en revenant à la figure initiale.
\end{enumerate}
\textbf{Réflexe classique (problème de Fermat / Héron) :} pour minimiser une somme de distances $MA+MB$ avec $M$ contraint d'appartenir à un plan $(P)$, remplacer $A$ par son symétrique $A'=s_{P}(A)$ : alors $MA=MA'$ et le minimum est atteint quand $M$, $A'$, $B$ sont alignés (et $M$ est l'intersection de $(A'B)$ avec $(P)$).
\end{methode}

\begin{exoresolu}[Problème de distance minimale — le câble de CAMTEL]
L'espace est muni d'un repère orthonormé $(O\,;\vec{\imath},\vec{\jmath},\vec{k})$ (unité : \SI{100}{m}). Une entreprise de télécommunications doit poser un câble reliant le central $A(0\,;0\,;3)$ au relais $B(6\,;4\,;5)$, en passant par un point $M$ d'une nappe phréatique modélisée par le plan $(P):z=0$, où le câble doit être inspecté. On cherche la position de $M\in(P)$ qui minimise la longueur totale $MA+MB$.
\begin{enumerate}
\item Soit $A'=s_{P}(A)$ où $s_{P}$ est la réflexion de plan $(P)$. Déterminer les coordonnées de $A'$.
\item Montrer que pour tout point $M$ de $(P)$, $MA=MA'$.
\item En déduire la position de $M$ minimisant $MA+MB$, puis la longueur minimale du câble en mètres.
\end{enumerate}
\tcblower
\textbf{1. Coordonnées de $A'$.}

Le plan $(P):z=0$ est le plan $(xOy)$. La réflexion de plan $(P)$ est $s_{P}:(x\,;y\,;z)\mapsto(x\,;y\,;-z)$. Donc $A'(0\,;0\,;-3)$.

\textbf{2. Pour tout $M\in(P)$, $MA=MA'$.}

Soit $M\in(P)$. Comme $M$ appartient au plan de la réflexion, $M$ est invariant par $s_{P}$ : $s_{P}(M)=M$. Or une réflexion est une isométrie, elle conserve les distances :
\[MA=d(M,A)=d\big(s_{P}(M),s_{P}(A)\big)=d(M,A')=MA'.\]
Autre rédaction possible : $(P)$ est le plan médiateur de $[AA']$ (car $(AA')\perp(P)$ et le milieu $(0\,;0\,;0)$ de $[AA']$ appartient à $(P)$), donc tout point de $(P)$ est équidistant de $A$ et $A'$.

\textbf{3. Position optimale et longueur minimale.}

Pour tout $M\in(P)$ : $MA+MB=MA'+MB$. D'après l'inégalité triangulaire :
\[MA'+MB\geq A'B,\]
avec égalité si et seulement si $M$ appartient au segment $[A'B]$.

Vérifions que la droite $(A'B)$ coupe bien $(P)$ : paramétrons $(A'B)$ : $M(t)=A'+t\overrightarrow{A'B}=(6t\,;\,4t\,;\,-3+8t)$, $t\in\mathbb{R}$. La cote s'annule pour $-3+8t=0$, soit $t=\dfrac{3}{8}\in[0\,;1]$ : le point obtenu appartient bien au segment $[A'B]$. Alors :
\[M\left(6\times\frac{3}{8}\,;\,4\times\frac{3}{8}\,;\,0\right)=\left(\frac{9}{4}\,;\,\frac{3}{2}\,;\,0\right).\]

La longueur minimale est :
\[A'B=\sqrt{(6-0)^{2}+(4-0)^{2}+(5-(-3))^{2}}=\sqrt{36+16+64}=\sqrt{116}=2\sqrt{29}.\]

\textbf{Conclusion :} le point d'inspection optimal est $M\left(\dfrac{9}{4}\,;\dfrac{3}{2}\,;0\right)$, soit à \SI{225}{m} Est et \SI{150}{m} Nord de l'origine, et la longueur minimale du câble est $2\sqrt{29}\times 100\approx \SI{1077}{m}$.
\end{exoresolu}

\begin{savaistu}[Isométries et cristallographie au Cameroun]
Les isométries de l'espace (translations, rotations, réflexions, demi-tours, symétries centrales) sont les briques fondamentales de la \textbf{cristallographie}. Les 230 \emph{groupes d'espace} classent toutes les structures périodiques possibles des cristaux. Au Cameroun, l'exploitation minière (bauxite à Minim-Martap, fer à Mbalam, diamant à Mobilong) repose sur la connaissance précise des mailles cristallines. Un géologue du CAPAM (Cadre d'Appui à la Petite Exploitation Minière) utilise la diffraction des rayons X : les taches de diffraction révèlent les symétries du cristal (plans de réflexion, axes de rotation/demi-tour). Comprendre la composée de deux réflexions (translation ou rotation) permet d'identifier le groupe d'espace et donc la structure atomique du minerai, optimisant ainsi son extraction.
\end{savaistu}

\begin{attention}[Les 5 erreurs qui coûtent des points au Bac]
\begin{enumerate}
\item \textbf{Oublier le facteur 2 dans les composées.} La composée de deux réflexions de plans parallèles distants de $d$ est une translation de vecteur de norme $2d$ (et non $d$). De même pour deux demi-tours d'axes parallèles. C'est l'erreur la plus fréquente.
\item \textbf{Confondre le sens de la composition.} $s_{P'}\circ s_{P}$ signifie qu'on applique \emph{d'abord} $s_{P}$, \emph{ensuite} $s_{P'}$ : la translation obtenue est dirigée de $(P)$ vers $(P')$. Inverser l'ordre donne le vecteur opposé — et une réponse fausse si l'énoncé demande les éléments caractéristiques.
\item \textbf{Donner la nature sans les éléments caractéristiques.} Écrire « c'est une translation » sans donner le vecteur, « c'est un demi-tour » sans donner l'axe (point + vecteur directeur), « c'est une réflexion » sans donner le plan : la réponse est incomplète et perd la moitié des points.
\item \textbf{Mal traiter les conditions dans le calcul analytique.} Pour une réflexion, il faut \emph{deux} conditions : milieu dans le plan \textbf{et} $\overrightarrow{MM'}$ colinéaire au vecteur normal. Pour un demi-tour : milieu sur l'axe \textbf{et} $\overrightarrow{MM'}\cdot\vec{u}=0$. En oublier une conduit à un système sous-déterminé.
\item \textbf{Affirmer qu'une application est une isométrie sans vérifier, ou mal rédiger la reconnaissance.} Pour identifier une isométrie, il faut résoudre proprement le système des points fixes $f(M)=M$ et conclure en comparant aux expressions types. Dire « on reconnaît » sans justification ne vaut rien au Bac : écris le système, résous-le, puis conclus par une phrase complète.
\end{enumerate}
\end{attention}

\newpage

\coursec{Exercices}

\courssub{Je vérifie mes connaissances}

\begin{exercice}[QCM : Nature des isométries]
Pour chaque question, une seule réponse est exacte. Coche-la et justifie brièvement.
\begin{enumerate}
\item La composée de deux réflexions de plans \emph{strictement parallèles} est :
\begin{enumerate}
\item une réflexion ; \item une translation ; \item un demi-tour ; \item l'identité.
\end{enumerate}
\item La composée de deux demi-tours d'axes \emph{orthogonaux et sécants} est :
\begin{enumerate}
\item une translation ; \item une réflexion ; \item un demi-tour ; \item une rotation d'angle $\dfrac{\pi}{3}$.
\end{enumerate}
\item L'ensemble des points invariants par la réflexion de plan $\mathcal{P}$ est :
\begin{enumerate}
\item vide ; \item réduit à un point ; \item le plan $\mathcal{P}$ ; \item une droite perpendiculaire à $\mathcal{P}$.
\end{enumerate}
\item Soit $s$ la réflexion de plan $\mathcal{P}$. Alors $s \circ s$ est :
\begin{enumerate}
\item $s$ ; \item l'identité ; \item une translation ; \item un demi-tour.
\end{enumerate}
\end{enumerate}
\end{exercice}

\begin{exercice}[Vrai ou Faux]
Réponds par Vrai ou Faux en justifiant chaque réponse.
\begin{enumerate}
\item Toute réflexion de l'espace conserve les distances.
\item La composée de deux réflexions de plans perpendiculaires est une translation.
\item Un demi-tour d'axe $(d)$ laisse invariant tout point de $(d)$.
\item La composée de deux demi-tours d'axes parallèles distincts est un demi-tour.
\item Si $M'$ est l'image de $M$ par la réflexion de plan $\mathcal{P}$, alors le milieu de $[MM']$ appartient à $\mathcal{P}$.
\end{enumerate}
\end{exercice}

\begin{exercice}[Définitions et propriétés]
\begin{enumerate}
\item Donne la définition de la réflexion de plan $\mathcal{P}$ dans l'espace.
\item Donne la définition du demi-tour d'axe $(d)$ dans l'espace.
\item Cite deux propriétés de conservation communes à toutes les isométries de l'espace.
\item Soit $s_{(d)}$ le demi-tour d'axe $(d)$. Que vaut $s_{(d)} \circ s_{(d)}$ ? Justifie.
\end{enumerate}
\end{exercice}

\begin{exercice}[Calculs directs]
L'espace est muni d'un repère orthonormé $(O\,;\vec{i},\vec{j},\vec{k})$.
\begin{enumerate}
\item Détermine l'image du point $A(3\,;-1\,;2)$ par la réflexion de plan $(xOy)$ d'équation $z=0$.
\item Détermine l'image du point $B(1\,;4\,;-2)$ par le demi-tour d'axe $(Oz)$.
\item Soit $t$ la translation de vecteur $\vec{u}(0\,;0\,;4)$. Écris $t$ comme composée de deux réflexions de plans parallèles (donne les équations de deux plans convenables).
\end{enumerate}
\end{exercice}

\courssub{Je m'entraîne}

\begin{aretenir}[Formules analytiques -- Réflexion et Demi-tour]
Dans un repère orthonormé $(O;\vec{i},\vec{j},\vec{k})$ :
\begin{itemize}
\item \textbf{Réflexion de plan $\mathcal{P}: ax+by+cz+d=0$ (normal $\vec{n}(a,b,c)$) :}
$M'(x',y',z')$ image de $M(x,y,z)$ :
\[
\begin{cases}
x' = x - 2\dfrac{ax+by+cz+d}{a^2+b^2+c^2}\,a \\[6pt]
y' = y - 2\dfrac{ax+by+cz+d}{a^2+b^2+c^2}\,b \\[6pt]
z' = z - 2\dfrac{ax+by+cz+d}{a^2+b^2+c^2}\,c
\end{cases}
\]
\item \textbf{Demi-tour d'axe $(d)$ passant par $A(x_A,y_A,z_A)$ et de direction $\vec{u}(a,b,c)$ ($\|\vec{u}\|=1$) :}
$M'$ est le symétrique de $M$ par rapport à $(d)$. On utilise le projeté orthogonal $H$ de $M$ sur $(d)$ : $\overrightarrow{HM'} = -\overrightarrow{HM}$, d'où $M' = 2H - M$.
Si l'axe est une droite coordonnée (ex: $(Oz)$ : $x=0,y=0$) : $(x,y,z)\mapsto(-x,-y,z)$.
\end{itemize}
\end{aretenir}

\begin{exercice}[Expression analytique d'une réflexion]
L'espace est muni d'un repère orthonormé $(O\,;\vec{i},\vec{j},\vec{k})$. On considère le plan $\mathcal{P}$ d'équation $2x-y+2z-3=0$ et le point $A(1\,;2\,;-1)$.
\begin{enumerate}
\item Vérifie que le vecteur $\vec{n}(2\,;-1\,;2)$ est normal à $\mathcal{P}$ et détermine une équation paramétrique de la droite $(d)$ passant par $A$ et perpendiculaire à $\mathcal{P}$.
\item Détermine les coordonnées du projeté orthogonal $H$ de $A$ sur $\mathcal{P}$.
\item En déduire les coordonnées de $A'$, image de $A$ par la réflexion $s$ de plan $\mathcal{P}$.
\item Établis l'expression analytique de $s$ : pour $M(x\,;y\,;z)$ d'image $M'(x'\,;y'\,;z')$, exprime $x'$, $y'$, $z'$ en fonction de $x$, $y$, $z$.
\item Vérifie que le milieu de $[AA']$ appartient bien à $\mathcal{P}$.
\end{enumerate}
\end{exercice}

\begin{exercice}[Reconnaissance et écriture comme composée]
Soit $f$ l'application de l'espace définie, dans un repère orthonormé, par
\[ f(x\,;y\,;z)=(x\,;-y+4\,;z). \]
\begin{enumerate}
\item Détermine l'ensemble des points invariants par $f$.
\item Reconnais la nature de $f$ et précise ses éléments caractéristiques.
\item Soit $t$ la translation de vecteur $\vec{u}(0\,;4\,;0)$. Écris $t$ comme composée de deux réflexions de plans parallèles de deux façons différentes (donne les équations des plans dans chaque cas).
\end{enumerate}
\end{exercice}

\begin{exercice}[Composée de deux demi-tours]
Dans un repère orthonormé $(O\,;\vec{i},\vec{j},\vec{k})$, on considère les droites $(d_1)$ et $(d_2)$ définies par
\[ (d_1):\begin{cases} x=0 \\ z=1 \end{cases} \qquad (d_2):\begin{cases} x=2 \\ z=1 \end{cases} \]
On note $s_1$ et $s_2$ les demi-tours d'axes respectifs $(d_1)$ et $(d_2)$.
\begin{enumerate}
\item Justifie que $(d_1)$ et $(d_2)$ sont parallèles et précise leur direction.
\item Détermine les expressions analytiques de $s_1$ et de $s_2$.
\item Détermine l'expression analytique de $s_2 \circ s_1$.
\item En déduire la nature et les éléments caractéristiques de $s_2 \circ s_1$.
\item Calcule l'image du point $B(3\,;5\,;0)$ par $s_2 \circ s_1$.
\end{enumerate}
\end{exercice}

\begin{exercice}[Reconnaissance guidée d'une isométrie : le vis]
On considère l'application $g$ de l'espace définie, dans un repère orthonormé, par
\[ g(x\,;y\,;z)=(-x+2\,;-y\,;z-1). \]
\begin{enumerate}
\item Recherche les éventuels points fixes de $g$.
\item Soit $t$ la translation de vecteur $\vec{u}(0\,;0\,;-1)$ et $h$ l'application définie par $h(x\,;y\,;z)=(-x+2\,;-y\,;z)$. Vérifie que $g = t \circ h$.
\item Reconnais la nature de $h$ et précise ses éléments caractéristiques.
\item Conclus : $g$ est la composée d'un demi-tour et d'une translation. Une telle isométrie sans point fixe, dont le vecteur de translation est \textbf{parallèle} à l'axe du demi-tour, est appelée \emph{vis} (ou déplacement hélicoïdal d'angle $\pi$). Vérifie que l'axe du demi-tour $h$ est dirigé par un vecteur \textbf{parallèle} à $\vec{u}$, puis que $g$ peut aussi s'écrire comme composée d'un demi-tour d'axe $(d')$ à préciser et d'une translation (même axe, même translation).
\end{enumerate}
\end{exercice}

\begin{attention}[Pièges classiques]
\begin{itemize}
\item \textbf{Ordre de la composée :} $s_2\circ s_1$ signifie « on applique $s_1$ \emph{puis} $s_2$ ». Pour les translations issues de deux réflexions, $s_{\mathcal{P}_2}\circ s_{\mathcal{P}_1}$ donne un vecteur orienté de $\mathcal{P}_1$ vers $\mathcal{P}_2$.
\item \textbf{Plans parallèles vs perpendiculaires :} Parallèles $\to$ Translation. Perpendiculaires $\to$ Demi-tour (axe = intersection).
\item \textbf{Demi-tours :} Axes parallèles $\to$ Translation. Axes sécants orthogonaux $\to$ Demi-tour (axe $\perp$ aux deux). Axes sécants non orthogonaux $\to$ Rotation d'angle $2\alpha$ (non au programme explicite mais bon à savoir).
\item \textbf{Points fixes :} Réflexion $\to$ un plan. Demi-tour $\to$ une droite. Translation $\to$ aucun (sauf identité). Vis $\to$ aucun.
\item \textbf{Équation de plan :} Vérifie toujours que les points donnés appartiennent au plan avant d'en déduire l'équation (ex: cube).
\end{itemize}
\end{attention}

\courssub{Je me prépare au Bac}

\begin{exercice}[Type Bac : Le cube et les symétries]
L'espace est rapporté au repère orthonormé $(A\,;\overrightarrow{AB},\overrightarrow{AD},\overrightarrow{AE})$. $ABCDEFGH$ est un cube d'arête $1$, avec les notations habituelles :
$A(0,0,0)$, $B(1,0,0)$, $D(0,1,0)$, $E(0,0,1)$, $C(1,1,0)$, $F(1,0,1)$, $H(0,1,1)$, $G(1,1,1)$.
On note $\mathcal{P}$ le plan $(ACGE)$ et $s$ la réflexion de plan $\mathcal{P}$.

\begin{popfigure}
\begin{tikzpicture}[scale=1.2, >=Stealth]
% Coordinates
\coordinate (A) at (0,0,0);
\coordinate (B) at (2,0,0);
\coordinate (D) at (0,2,0);
\coordinate (E) at (0,0,2);
\coordinate (C) at (2,2,0);
\coordinate (F) at (2,0,2);
\coordinate (H) at (0,2,2);
\coordinate (G) at (2,2,2);
% Edges
\draw[popline, thick] (A) -- (B) -- (C) -- (D) -- cycle;
\draw[popline, thick] (E) -- (F) -- (G) -- (H) -- cycle;
\draw[popline, thick] (A) -- (E) (B) -- (F) (C) -- (G) (D) -- (H);
% Hidden edges dashed
\draw[popline, dashed, thick] (A) -- (D) (A) -- (E) (B) -- (C) (D) -- (C) (E) -- (H) (F) -- (G) (H) -- (G);
% Plane ACGE
\fill[popblue!15, opacity=0.6] (A) -- (C) -- (G) -- (E) -- cycle;
\draw[popblue, thick] (A) -- (C) -- (G) -- (E) -- cycle;
% Labels
\foreach \p/\pos in {A/below left, B/below right, C/above right, D/above left, E/below left, F/below right, G/above right, H/above left}
  \node[popdark, circle, fill=popcream, inner sep=1.5pt, label=\pos:\p] at (\p) {};
% Axis AC
\draw[poporange, very thick, ->] (A) -- (C) node[midway, above, sloped, poporange] {$d=(AC)$};
% Normal vector
\draw[popgreen, very thick, ->] (1,1,1) -- (1.5,0.5,1) node[midway, right, popgreen] {$\vec{n}(1,-1,0)$};
\end{tikzpicture}
\legende{Cube unité et plan $\mathcal{P}=(ACGE)$ d'équation $x-y=0$.}
\end{popfigure}

\textbf{Partie A : Étude de la réflexion $s$}
\begin{enumerate}
\item Détermine les coordonnées des points $A$, $C$, $G$, $E$ et montre qu'une équation cartésienne de $\mathcal{P}$ est $x-y=0$.
\item Montre que les points $B$ et $D$ sont échangés par $s$ (on pourra montrer que $\mathcal{P}$ est le plan médiateur de $[BD]$).
\item En déduire, sans calcul, que pour tout point $M$ de l'espace, $BM = D\,s(M)$. Vérifie cette égalité pour $M=A$, $M=C$, $M=H$.
\end{enumerate}

\textbf{Partie B : Composée avec un demi-tour}
\begin{enumerate}
\item On note $s_{(AC)}$ le demi-tour d'axe $(AC)$. Justifie que la droite $(AC)$ est incluse dans le plan $\mathcal{P}$.
\item On admet que la composée d'une réflexion de plan $\mathcal{P}$ et du demi-tour d'axe $(d)$ inclus dans $\mathcal{P}$ est la réflexion par rapport au plan $\mathcal{Q}$ passant par $(d)$ et perpendiculaire à $\mathcal{P}$. Détermine la nature de $s \circ s_{(AC)}$ et donne une équation du plan $\mathcal{Q}$.
\item Détermine l'image du sommet $E$ par $s \circ s_{(AC)}$.
\end{enumerate}
\end{exercice}

\begin{exercice}[Type Bac : La fourmi et le plus court chemin (SODECOTON Douala)]
Dans un entrepôt de la SODECOTON à Douala, on modélise l'espace par un repère orthonormé $(O\,;\vec{i},\vec{j},\vec{k})$ (unité : le mètre), le plan $(xOy)$ d'équation $z=0$ représentant le sol. Une fourmi part du point $A(0\,;0\,;2)$ situé sur un mur vertical, doit obligatoirement toucher le sol en un point $M$, puis rejoindre la mangeoire placée en $B(5\,;4\,;3)$, elle aussi au-dessus du sol. On cherche à minimiser la longueur totale $AM + MB$.

\begin{popfigure}
\begin{tikzpicture}[scale=0.6, >=Stealth]
% Axes
\draw[->, popmuted] (0,0,0) -- (7,0,0) node[right] {$x$};
\draw[->, popmuted] (0,0,0) -- (0,5,0) node[above] {$y$};
\draw[->, popmuted] (0,0,0) -- (0,0,4) node[above] {$z$};
% Ground plane
\fill[popgreen!10, opacity=0.5] (0,0,0) -- (7,0,0) -- (7,5,0) -- (0,5,0) -- cycle;
\draw[popgreen!50!black] (0,0,0) -- (7,0,0) -- (7,5,0) -- (0,5,0) -- cycle;
% Points
\coordinate (A) at (0,0,2);
\coordinate (B) at (5,4,3);
\coordinate (Bp) at (5,4,-3);
\coordinate (M0) at (2,1.6,0);
% Vertical lines
\draw[popmuted, dashed] (A) -- (0,0,0);
\draw[popmuted, dashed] (B) -- (5,4,0);
\draw[popmuted, dashed] (Bp) -- (5,4,0);
% Paths
\draw[popblue, thick] (A) -- (M0) -- (B);
\draw[poporange, thick, dashed] (A) -- (Bp);
% Labels
\node[popblue, left] at (A) {$A(0,0,2)$};
\node[popblue, right] at (B) {$B(5,4,3)$};
\node[poporange, right] at (Bp) {$B'(5,4,-3)$};
\node[popdark, below] at (M0) {$M_0(2,1.6,0)$};
\node[popdark, below left] at (0,0,0) {$O$};
\end{tikzpicture}
\legende{Trajet optimal de la fourmi : $A \to M_0 \to B$ via réflexion du but $B$ en $B'$.}
\end{popfigure}

\begin{enumerate}
\item Soit $B'$ le symétrique de $B$ par rapport au plan du sol. Détermine les coordonnées de $B'$.
\item Soit $M$ un point du plan d'équation $z=0$. Justifie que $MB = MB'$.
\item En déduire que $AM + MB \geq AB'$, et que l'égalité est réalisée lorsque $M$ est le point d'intersection de la droite $(AB')$ avec le plan du sol.
\item Détermine les coordonnées du point de contact $M_0$ qui minimise la distance parcourue.
\item Calcule la distance minimale parcourue par la fourmi. On donnera la valeur exacte puis une valeur approchée à $10^{-2}$ près.
\item Question bonus : la fourmi se déplace à \SI{2}{cm/s}. Combien de temps (en minutes et secondes) met-elle au minimum pour rejoindre la mangeoire ?
\end{enumerate}
\end{exercice}

\begin{exercice}[Type Bac : Sécurisation d'un chantier par symétries (BTP Yaoundé)]
Une entreprise de BTP à Yaoundé modélise un chantier dans un repère orthonormé $(O\,;\vec{i},\vec{j},\vec{k})$. Deux murs verticaux sont modélisés par les plans $\mathcal{P}_1$ d'équation $x=0$ et $\mathcal{P}_2$ d'équation $x=3$. Un projecteur est placé au point $L(1\,;2\,;5)$. On note $s_1$ et $s_2$ les réflexions de plans respectifs $\mathcal{P}_1$ et $\mathcal{P}_2$.

\begin{popfigure}
\begin{tikzpicture}[scale=0.7, >=Stealth]
% Axes
\draw[->, popmuted] (0,0,0) -- (5,0,0) node[right] {$x$};
\draw[->, popmuted] (0,0,0) -- (0,4,0) node[above] {$y$};
\draw[->, popmuted] (0,0,0) -- (0,0,6) node[above] {$z$};
% Planes
\fill[popblue!15, opacity=0.5] (0,0,0) -- (0,4,0) -- (0,4,6) -- (0,0,6) -- cycle; % P1 x=0
\fill[poporange!15, opacity=0.5] (3,0,0) -- (3,4,0) -- (3,4,6) -- (3,0,6) -- cycle; % P2 x=3
\draw[popblue, thick] (0,0,0) -- (0,4,0) -- (0,4,6) -- (0,0,6) -- cycle;
\draw[poporange, thick] (3,0,0) -- (3,4,0) -- (3,4,6) -- (3,0,6) -- cycle;
% Ground
\fill[popgreen!10, opacity=0.3] (0,0,0) -- (5,0,0) -- (5,4,0) -- (0,4,0) -- cycle;
\draw[popgreen!50!black] (0,0,0) -- (5,0,0) -- (5,4,0) -- (0,4,0) -- cycle;
% Point L
\coordinate (L) at (1,2,5);
\coordinate (L1) at (-1,2,5); % s1(L)
\coordinate (L2) at (5,2,5);  % s2(L)
\coordinate (L12) at (7,2,5); % s2(s1(L))
\draw[popblue, ->, thick] (L) -- (L1) node[midway, left] {$s_1$};
\draw[poporange, ->, thick] (L1) -- (L12) node[midway, above] {$s_2$};
\draw[poppurple, ->, thick, dashed] (L) -- (L12) node[midway, below] {$s_2\circ s_1$};
% Labels
\node[popdark, circle, fill=popcream, inner sep=2pt, label={above right:$L(1,2,5)$}] at (L) {};
\node[popblue, circle, fill=popblue!20, inner sep=2pt, label=left:$s_1(L)$] at (L1) {};
\node[poporange, circle, fill=poporange!20, inner sep=2pt, label=right:$s_2(L)$] at (L2) {};
\node[poppurple, circle, fill=poppurple!20, inner sep=2pt, label=above right:$s_2(s_1(L))$] at (L12) {};
\node[popblue, anchor=east] at (0,2,6) {$\mathcal{P}_1: x=0$};
\node[poporange, anchor=west] at (3,2,6) {$\mathcal{P}_2: x=3$};
\node[popgreen!50!black, anchor=north west] at (5,4,0) {$\mathcal{P}_3: z=0$};
\end{tikzpicture}
\legende{Chantier BTP : murs $\mathcal{P}_1, \mathcal{P}_2$, sol $\mathcal{P}_3$, projecteur $L$ et ses images.}
\end{popfigure}

\textbf{Partie A}
\begin{enumerate}
\item Détermine les expressions analytiques de $s_1$ et de $s_2$.
\item Détermine l'expression analytique de $s_2 \circ s_1$, reconnais sa nature et précise ses éléments caractéristiques.
\item Calcule l'image du projecteur $L$ par $s_2 \circ s_1$.
\item Compare $s_1 \circ s_2$ et $s_2 \circ s_1$. La composition des réflexions est-elle commutative en général ?
\end{enumerate}

\textbf{Partie B}
On considère maintenant le plan horizontal $\mathcal{P}_3$ d'équation $z=0$ (le sol) et la réflexion $s_3$ de plan $\mathcal{P}_3$.
\begin{enumerate}
\item Justifie que $\mathcal{P}_1$ et $\mathcal{P}_3$ sont perpendiculaires et détermine leur droite d'intersection $(\Delta)$.
\item En admettant que la composée de deux réflexions de plans perpendiculaires est le demi-tour dont l'axe est la droite d'intersection des deux plans, déduis-en la nature et les éléments caractéristiques de $s_3 \circ s_1$.
\item Détermine l'expression analytique de $s_3 \circ s_1$ et vérifie le résultat de la question précédente en calculant l'image de $L$.
\end{enumerate}
\end{exercice}

\begin{savaistu}[Le vis (déplacement hélicoïdal) en mécanique et architecture]
Un \textbf{vis} (ou déplacement hélicoïdal) est la composée d'une rotation d'angle $\theta$ autour d'un axe $(d)$ et d'une translation de vecteur $\vec{u}$ \textbf{parallèle} à $(d)$. C'est le mouvement d'une vis dans son écrou, d'un tire-bouchon, ou d'un escalier hélicoïdal (comme celui de la tour de la cathédrale de Yaoundé ou les escaliers de certains silos à grains du Nord-Cameroun).
Si $\theta=\pi$, c'est la composée d'un demi-tour et d'une translation parallèle à l'axe : c'est le cas de l'exercice 8 ($g$). Un vis n'a \textbf{aucun point fixe} (sauf si la translation est nulle, auquel cas c'est une simple rotation). Toute isométrie directe de l'espace (déterminant $+1$) sans point fixe est un vis (théorème de Chasles dans l'espace).
\end{savaistu}

\begin{methode}[Reconnaître une isométrie donnée par son expression analytique]
Soit $f(x,y,z) = (x',y',z')$ une application affine de l'espace.
\begin{enumerate}
\item \textbf{Cherche les points fixes} : résous $f(M)=M$.
\begin{itemize}
\item Plan de points fixes $\to$ \textbf{Réflexion} (plan = ensemble des fixes).
\item Droite de points fixes $\to$ \textbf{Demi-tour} (axe = ensemble des fixes) \textbf{ou} \textbf{Rotation} (si angle $\neq \pi$).
\item Aucun point fixe $\to$ \textbf{Translation} (si $f(M)-M$ constant) \textbf{ou} \textbf{Vis} (si partie linéaire est une rotation $\neq$ Id).
\item Tout l'espace $\to$ \textbf{Identité}.
\end{itemize}
\item \textbf{Analyse la partie linéaire} (matrice $A$ telle que $f(M)=A\cdot M + B$).
\begin{itemize}
\item $\det A = -1$ : isométrie \emph{inverse} (réflexion, symétrie centrale, composition impaire de réflexions).
\item $\det A = +1$ : isométrie \emph{directe} (translation, rotation, demi-tour, vis).
\item Valeurs propres de $A$ : $1$ (axe/direction invariante), $-1$ (inversion), complexes $e^{\pm i\theta}$ (rotation).
\end{itemize}
\item \textbf{Conclusion} :
\begin{itemize}
\item $A=\text{Id}$, $B\neq 0$ $\to$ Translation de vecteur $B$.
\item $A=\text{diag}(1,1,-1)$ (ou similaire) $\to$ Réflexion (plan orthogonal à l'axe $-1$).
\item $A=\text{diag}(-1,-1,1)$ $\to$ Demi-tour (axe = direction de la valeur propre $1$).
\item $A=\text{diag}(-1,-1,-1)$ $\to$ Symétrie centrale (point fixe unique).
\item $A$ rotation dans un plan + $1$ sur l'orthogonal, $B\parallel$ axe $\to$ Vis.
\end{itemize}
\end{enumerate}
\end{methode}

\begin{exercice}[Entraînement Bac : Isométries et pavage]
On considère le pavage de l'espace par des cubes unitaires de sommets les points à coordonnées entières $(k,\ell,m)\in\mathbb{Z}^3$. Soit $s$ la réflexion par rapport au plan $\mathcal{P}: x+y+z=0$ et $t$ la translation de vecteur $\vec{v}(1,1,1)$.
\begin{enumerate}
\item Détermine l'expression analytique de $s$.
\item Montre que $t$ laisse invariant le plan $\mathcal{P}$.
\item Étudie la composée $f = t \circ s$. Est-ce une isométrie ? Laquelle ? A-t-elle des points fixes ?
\item Dessine (schéma) l'effet de $f$ sur le cube $[0,1]^3$.
\end{enumerate}
\end{exercice}

\begin{aretenir}[Résumé des compositions d'isométries (Programme Terminale C/E)]
\begin{center}
\begin{tabularx}{\linewidth}{|l|X|X|}\hline
\textbf{Composition} & \textbf{Résultat} & \textbf{Éléments caractéristiques} \\ \hline
$s_{\mathcal{P}_2} \circ s_{\mathcal{P}_1}$ ($\mathcal{P}_1\parallel\mathcal{P}_2$) & Translation & $\vec{v} = 2\overrightarrow{\mathcal{P}_1\mathcal{P}_2}$ (orienté $\mathcal{P}_1\to\mathcal{P}_2$, $\perp$ aux plans) \\ \hline
$s_{\mathcal{P}_2} \circ s_{\mathcal{P}_1}$ ($\mathcal{P}_1\perp\mathcal{P}_2$) & Demi-tour & Axe = $\mathcal{P}_1 \cap \mathcal{P}_2$ \\ \hline
$s_{(d_2)} \circ s_{(d_1)}$ ($d_1\parallel d_2$) & Translation & $\vec{v} = 2\overrightarrow{A_1A_2}$ ($A_i$ projetés sur un plan $\perp$, $\vec{v}\perp d_i$) \\ \hline
$s_{(d_2)} \circ s_{(d_1)}$ ($d_1\perp d_2$ sécantes) & Demi-tour & Axe = droite $\perp$ aux deux par l'intersection \\ \hline
$s_{\mathcal{P}} \circ s_{(d)}$ ($d\subset\mathcal{P}$) & Réflexion & Plan $\mathcal{Q}\supset d$, $\mathcal{Q}\perp\mathcal{P}$ \\ \hline
$s_{(d)} \circ t_{\vec{u}}$ ($\vec{u}\parallel d$) & Vis (angle $\pi$) & Axe $d$, translation $\vec{u}$ (pas de point fixe) \\ \hline
\end{tabularx}
\end{center}
\textbf{Propriétés générales :} Toute isométrie conserve distances, angles, aires, volumes. Composée de deux isométries = isométrie. Inverse d'une isométrie = isométrie. Déterminant de la partie linéaire : $+1$ (directe : translation, rotation, demi-tour, vis), $-1$ (indirecte : réflexion, symétrie centrale, composition impaire de réflexions).
\end{aretenir}

\newpage

\coursec{Corrigés des exercices}

\courssub{Je vérifie mes connaissances}

\begin{corrige}[Exercice 1 -- QCM : Nature des isométries]
\begin{enumerate}
\item \textbf{Bonne réponse : 2 (translation).} Théorème du cours : la composée de deux réflexions de plans \emph{strictement parallèles} $\mathcal{P}_1$ et $\mathcal{P}_2$ est la translation de vecteur $2\vec{d}$, où $\vec{d}$ est le vecteur perpendiculaire aux deux plans, orienté de $\mathcal{P}_1$ vers $\mathcal{P}_2$, de norme égale à la distance entre les plans. Ce n'est ni une réflexion (qui a un plan de points fixes), ni un demi-tour, ni l'identité (les plans sont distincts).
\item \textbf{Bonne réponse : 3 (demi-tour).} Un demi-tour est une rotation d'angle $\pi$. La composée de deux demi-tours d'axes sécants en $O$ et orthogonaux est une rotation d'angle $2\times\dfrac{\pi}{2}=\pi$, c'est-à-dire un \cle{demi-tour}, dont l'axe est la droite passant par $O$ et orthogonale aux deux axes.
\item \textbf{Bonne réponse : 3 (le plan $\mathcal{P}$).} Par définition, tout point de $\mathcal{P}$ est invariant. Réciproquement, si $M\notin\mathcal{P}$, son image $M'$ est distincte de $M$ (le plan $\mathcal{P}$ est le plan médiateur de $[MM']$), donc $M$ n'est pas fixe. L'ensemble des points invariants est exactement $\mathcal{P}$.
\item \textbf{Bonne réponse : 2 (l'identité).} Une réflexion est une \cle{involution} : appliquer deux fois la symétrie ramène chaque point sur lui-même, donc $s\circ s=\mathrm{Id}$.
\end{enumerate}
\textbf{Points de vigilance :} ne pas confondre « plans parallèles » (translation) et « plans perpendiculaires » (demi-tour) ; retenir que réflexions et demi-tours sont des involutions.
\end{corrige}

\begin{corrige}[Exercice 2 -- Vrai ou Faux]
\begin{enumerate}
\item \textbf{Vrai.} Une réflexion est une isométrie : elle conserve les distances (et par suite les angles, les aires et les volumes).
\item \textbf{Faux.} La composée de deux réflexions de plans \emph{perpendiculaires} est un \cle{demi-tour} dont l'axe est la droite d'intersection des deux plans. C'est la composée de deux réflexions de plans \emph{parallèles} qui est une translation.
\item \textbf{Vrai.} Par définition du demi-tour d'axe $(d)$, tout point de $(d)$ est sa propre image : $(d)$ est exactement l'ensemble des points fixes.
\item \textbf{Faux.} La composée de deux demi-tours d'axes \emph{parallèles distincts} est une \cle{translation} de vecteur $2\overrightarrow{A_1A_2}$, où $A_1$ et $A_2$ sont les projetés orthogonaux d'un même point sur les deux axes. Ce serait un demi-tour si les axes étaient sécants et orthogonaux.
\item \textbf{Vrai.} C'est la propriété caractéristique de la réflexion : $\mathcal{P}$ est le plan médiateur de $[MM']$, donc il contient le milieu de $[MM']$ (et est perpendiculaire à $(MM')$).
\end{enumerate}
\end{corrige}

\begin{corrige}[Exercice 3 -- Définitions et propriétés]
\begin{enumerate}
\item La \cle{réflexion} de plan $\mathcal{P}$ est l'application de l'espace qui à tout point $M$ associe le point $M'$ tel que :
\begin{itemize}
\item si $M\in\mathcal{P}$, alors $M'=M$ ;
\item si $M\notin\mathcal{P}$, alors $\mathcal{P}$ est le plan médiateur du segment $[MM']$ (c'est-à-dire : $(MM')\perp\mathcal{P}$ et le milieu de $[MM']$ appartient à $\mathcal{P}$).
\end{itemize}
\item Le \cle{demi-tour} d'axe $(d)$ est l'application de l'espace qui à tout point $M$ associe le point $M'$ tel que :
\begin{itemize}
\item si $M\in(d)$, alors $M'=M$ ;
\item si $M\notin(d)$, alors, dans le plan passant par $M$ et perpendiculaire à $(d)$, $M'$ est le symétrique de $M$ par rapport au point d'intersection $H$ de ce plan avec $(d)$ (autrement dit $\overrightarrow{HM'}=-\overrightarrow{HM}$).
\end{itemize}
C'est la rotation d'angle $\pi$ autour de $(d)$.
\item Propriétés de conservation communes à toutes les isométries : conservation des \cle{distances}, des \cle{angles}, de l'\cle{alignement}, du \cle{parallélisme}, de l'\cle{orthogonalité}, des aires et des volumes. (Deux citations suffisent.)
\item $s_{(d)}\circ s_{(d)}=\mathrm{Id}$. Justification : un demi-tour est une rotation d'angle $\pi$ ; le composer avec lui-même donne une rotation d'angle $\pi+\pi=2\pi$, c'est-à-dire l'identité. On dit que le demi-tour est une \cle{involution}.
\end{enumerate}
\end{corrige}

\begin{corrige}[Exercice 4 -- Calculs directs]
\begin{enumerate}
\item La réflexion de plan $(xOy)$ d'équation $z=0$ conserve $x$ et $y$ et change $z$ en $-z$ : $(x,y,z)\mapsto(x,y,-z)$. Donc $A'(3\,;-1\,;-2)$.
\item Le demi-tour d'axe $(Oz)$ fixe la cote $z$ et, dans chaque plan horizontal, agit comme la symétrie centrale de centre le point de l'axe : $(x,y,z)\mapsto(-x,-y,z)$. Donc $B'(-1\,;-4\,;-2)$.
\item On cherche $\mathcal{P}_1$ et $\mathcal{P}_2$ parallèles tels que $s_{\mathcal{P}_2}\circ s_{\mathcal{P}_1}=t$ avec $\vec{u}(0,0,4)$. Les plans doivent être perpendiculaires à $\vec{u}$ (donc d'équation $z=\text{constante}$) et distants de $\dfrac{\|\vec{u}\|}{2}=2$.
Choix : $\mathcal{P}_1:z=0$ et $\mathcal{P}_2:z=2$.
Vérification : $s_{\mathcal{P}_1}(x,y,z)=(x,y,-z)$ et $s_{\mathcal{P}_2}(x,y,z)=(x,y,4-z)$, donc
\[ s_{\mathcal{P}_2}\circ s_{\mathcal{P}_1}(x,y,z)=s_{\mathcal{P}_2}(x,y,-z)=(x,y,4-(-z))=(x,y,z+4)=t(x,y,z). \]
D'autres choix conviennent : $\mathcal{P}_1:z=1$, $\mathcal{P}_2:z=3$, etc. L'ordre compte : il faut orienter de $\mathcal{P}_1$ vers $\mathcal{P}_2$ dans le sens de $\vec{u}$.
\end{enumerate}
\end{corrige}

\courssub{Je m'entraîne}

\begin{corrige}[Exercice 5 -- Expression analytique d'une réflexion]
\begin{enumerate}
\item Le plan $\mathcal{P}:2x-y+2z-3=0$ admet pour vecteur normal $\vec{n}(2\,;-1\,;2)$ (coefficients de $x$, $y$, $z$). La droite $(d)$ passant par $A(1\,;2\,;-1)$ et dirigée par $\vec{n}$ admet la représentation paramétrique :
\[ (d):\begin{cases} x=1+2t \\ y=2-t \\ z=-1+2t \end{cases}\quad (t\in\mathbb{R}). \]
\item $H$ est l'intersection de $(d)$ et de $\mathcal{P}$. On injecte le paramétrage dans l'équation de $\mathcal{P}$ :
\begin{align*}
2(1+2t)-(2-t)+2(-1+2t)-3&=0\\
2+4t-2+t-2+4t-3&=0\\
9t-5&=0 \iff t=\frac{5}{9}.
\end{align*}
D'où $H\left(1+\dfrac{10}{9}\,;\,2-\dfrac{5}{9}\,;\,-1+\dfrac{10}{9}\right)$, c'est-à-dire $H\left(\dfrac{19}{9}\,;\,\dfrac{13}{9}\,;\,\dfrac{1}{9}\right)$.
\item $H$ est le milieu de $[AA']$, donc $A'=2H-A$ :
\[ A'=\left(\frac{38}{9}-1\,;\,\frac{26}{9}-2\,;\,\frac{2}{9}+1\right)=\left(\frac{29}{9}\,;\,\frac{8}{9}\,;\,\frac{11}{9}\right). \]
\item On applique la formule du cours avec $a=2$, $b=-1$, $c=2$, $d=-3$ et $a^2+b^2+c^2=9$. On pose $k=\dfrac{2x-y+2z-3}{9}$. Alors :
\[ \begin{cases}
x'=x-4k=\dfrac{1}{9}\bigl(x+4y-8z+12\bigr)\\[4pt]
y'=y+2k=\dfrac{1}{9}\bigl(4x+7y+4z-6\bigr)\\[4pt]
z'=z-4k=\dfrac{1}{9}\bigl(-8x+4y+z+12\bigr)
\end{cases} \]
Contrôle avec $A(1,2,-1)$ : $k=\dfrac{2-2-2-3}{9}=-\dfrac{5}{9}$, donc $x'=1+\dfrac{20}{9}=\dfrac{29}{9}$, $y'=2-\dfrac{10}{9}=\dfrac{8}{9}$, $z'=-1+\dfrac{20}{9}=\dfrac{11}{9}$ : on retrouve $A'$.
\item Le milieu de $[AA']$ est $\left(\dfrac{1+\frac{29}{9}}{2}\,;\,\dfrac{2+\frac{8}{9}}{2}\,;\,\dfrac{-1+\frac{11}{9}}{2}\right)=\left(\dfrac{19}{9}\,;\,\dfrac{13}{9}\,;\,\dfrac{1}{9}\right)=H$, et $2\times\dfrac{19}{9}-\dfrac{13}{9}+2\times\dfrac{1}{9}-3=\dfrac{38-13+2-27}{9}=0$, donc ce milieu appartient bien à $\mathcal{P}$.
\end{enumerate}
\textbf{Points de vigilance :} ne pas oublier le facteur $2$ dans la formule ; diviser par $a^2+b^2+c^2$ (ici $9$) et non par $\|\vec{n}\|$ ; vérifier le résultat sur un point particulier.
\end{corrige}

\begin{corrige}[Exercice 6 -- Reconnaissance et écriture comme composée]
\begin{enumerate}
\item $M(x,y,z)$ est invariant si et seulement si $f(M)=M$, c'est-à-dire :
\[ \begin{cases} x=x \\ y=-y+4 \\ z=z \end{cases}\iff y=2. \]
L'ensemble des points invariants est le \cle{plan} $\mathcal{P}$ d'équation $y=2$.
\item $f$ admet un plan de points fixes et s'écrit $(x,y,z)\mapsto(x,\,4-y,\,z)$, c'est-à-dire $y'=2\times 2-y$ : c'est la \cle{réflexion} de plan $\mathcal{P}:y=2$. (La partie linéaire $\mathrm{diag}(1,-1,1)$ a pour déterminant $-1$, ce qui confirme une isométrie indirecte.)
\item On veut $s_{\mathcal{P}_2}\circ s_{\mathcal{P}_1}=t$ avec $\vec{u}(0,4,0)$. Les plans doivent être perpendiculaires à $\vec{u}$ (équations $y=\text{constante}$) et distants de $\dfrac{\|\vec{u}\|}{2}=2$, le second étant dans le sens de $\vec{u}$ par rapport au premier.
\begin{itemize}
\item \textbf{Choix 1 :} $\mathcal{P}_1:y=0$ et $\mathcal{P}_2:y=2$. Alors $s_{\mathcal{P}_1}(x,y,z)=(x,-y,z)$ et $s_{\mathcal{P}_2}(x,y,z)=(x,4-y,z)$, d'où $s_{\mathcal{P}_2}\circ s_{\mathcal{P}_1}(x,y,z)=(x,4-(-y),z)=(x,y+4,z)=t(x,y,z)$.
\item \textbf{Choix 2 :} $\mathcal{P}_1:y=1$ et $\mathcal{P}_2:y=3$. Alors $s_{\mathcal{P}_1}(x,y,z)=(x,2-y,z)$ et $s_{\mathcal{P}_2}(x,y,z)=(x,6-y,z)$, d'où $s_{\mathcal{P}_2}\circ s_{\mathcal{P}_1}(x,y,z)=(x,6-(2-y),z)=(x,y+4,z)=t(x,y,z)$.
\end{itemize}
Il y a une infinité de couples possibles : tout couple $(y=a\,;\,y=a+2)$ convient.
\end{enumerate}
\end{corrige}

\begin{corrige}[Exercice 7 -- Composée de deux demi-tours]
\begin{enumerate}
\item $(d_1)$ : $x=0$, $z=1$ et $(d_2)$ : $x=2$, $z=1$ sont toutes deux dirigées par $\vec{j}(0,1,0)$ (seule $y$ varie). Elles sont donc \cle{parallèles} (et distinctes). $(d_1)$ passe par $A_1(0,0,1)$ et $(d_2)$ par $A_2(2,0,1)$.
\item Pour $M(x,y,z)$, le projeté orthogonal sur $(d_1)$ est $H_1(0,y,1)$. Le demi-tour vérifie $\overrightarrow{H_1M'}=-\overrightarrow{H_1M}$, soit $M'=2H_1-M$ :
\[ s_1(x,y,z)=(-x,\,y,\,2-z). \]
De même, le projeté sur $(d_2)$ est $H_2(2,y,1)$ :
\[ s_2(x,y,z)=(4-x,\,y,\,2-z). \]
\item On compose ($s_1$ d'abord, puis $s_2$) :
\[ s_2\circ s_1(x,y,z)=s_2(-x,y,2-z)=\bigl(4-(-x),\,y,\,2-(2-z)\bigr)=(x+4,\,y,\,z). \]
\item On reconnaît la \cle{translation} de vecteur $\vec{v}(4,0,0)$. C'est conforme au théorème : $\vec{v}=2\overrightarrow{A_1A_2}=2(2,0,0)=(4,0,0)$, vecteur perpendiculaire aux deux axes, orienté de $(d_1)$ vers $(d_2)$.
\item $s_2\circ s_1(3,5,0)=(3+4,5,0)=(7\,;\,5\,;\,0)$.
\end{enumerate}
\textbf{Points de vigilance :} dans $s_2\circ s_1$, c'est $s_1$ qui agit en premier ; le vecteur de la translation est le \emph{double} du vecteur joignant les deux axes, pas le vecteur lui-même.
\end{corrige}

\begin{corrige}[Exercice 8 -- Reconnaissance guidée : le vis]
\begin{enumerate}
\item $g(M)=M\iff \begin{cases} x=-x+2 \\ y=-y \\ z=z-1 \end{cases}\iff \begin{cases} x=1 \\ y=0 \\ 0=-1 \end{cases}$. La dernière équation est impossible : $g$ \textbf{n'admet aucun point fixe}.
\item $t\circ h(x,y,z)=t(-x+2,\,-y,\,z)=(-x+2,\,-y,\,z-1)=g(x,y,z)$. Donc $g=t\circ h$.
\item Points fixes de $h$ : $x=-x+2\Rightarrow x=1$ ; $y=-y\Rightarrow y=0$ ; $z=z$ (toujours vrai). L'ensemble des points fixes est la droite $(d):\begin{cases} x=1 \\ y=0 \end{cases}$, parallèle à $(Oz)$. De plus $h(x,y,z)=(2-x,-y,z)$ : on reconnaît le \cle{demi-tour} d'axe $(d)$ (dans chaque plan horizontal, symétrie centrale de centre $(1,0,z)$).
\item L'axe $(d)$ est dirigé par $\vec{k}(0,0,1)$ et $\vec{u}(0,0,-1)=-\vec{k}$ : les vecteurs sont \cle{colinéaires}, donc $\vec{u}$ est parallèle à $(d)$. $g$ est la composée d'un demi-tour d'axe $(d)$ et d'une translation de vecteur parallèle à $(d)$ : c'est un \cle{vis} d'angle $\pi$, d'axe $(d'):x=1,\,y=0$ et de vecteur $\vec{u}(0,0,-1)$. Comme $\vec{u}$ est parallèle à l'axe, la translation et le demi-tour commutent : $g=t\circ s_{(d')}=s_{(d')}\circ t$.
\end{enumerate}
\textbf{Point de vigilance :} une isométrie sans point fixe n'est pas forcément une translation : il faut examiner la partie linéaire (ici $A=\mathrm{diag}(-1,-1,1)$, rotation d'angle $\pi$, et non l'identité).
\end{corrige}

\courssub{Je me prépare au Bac}

\begin{corrige}[Exercice 9 -- Type Bac : Le cube et les symétries]
\textbf{Partie A}
\begin{enumerate}
\item $A(0,0,0)$, $C(1,1,0)$, $G(1,1,1)$, $E(0,0,1)$. Chacun vérifie $x-y=0$ : $0-0=0$ pour $A$ et $E$ ; $1-1=0$ pour $C$ et $G$. Les quatre points appartiennent donc au plan d'équation $x-y=0$. Comme $A$, $C$, $E$ ne sont pas alignés ($\overrightarrow{AC}(1,1,0)$ et $\overrightarrow{AE}(0,0,1)$ non colinéaires), ce plan est bien le plan $(ACGE)$ : $\mathcal{P}:x-y=0$.
\item $B(1,0,0)$ et $D(0,1,0)$. Le milieu $I$ de $[BD]$ est $I\left(\dfrac{1}{2},\dfrac{1}{2},0\right)$, et $\dfrac{1}{2}-\dfrac{1}{2}=0$, donc $I\in\mathcal{P}$. De plus $\overrightarrow{BD}(-1,1,0)$ est colinéaire au vecteur normal $\vec{n}(1,-1,0)$ de $\mathcal{P}$, donc $(BD)\perp\mathcal{P}$. Ainsi $\mathcal{P}$ est le plan médiateur de $[BD]$, ce qui prouve que $s(B)=D$ et $s(D)=B$.
\item $s$ est une isométrie : pour tout point $M$, $BM=d(B,M)=d\bigl(s(B),s(M)\bigr)=d\bigl(D,s(M)\bigr)=D\,s(M)$.
Vérifications :
\begin{itemize}
\item $M=A$ : $A\in\mathcal{P}$ donc $s(A)=A$ ; $BA=1$ et $DA=1$. Égalité vérifiée.
\item $M=C$ : $C\in\mathcal{P}$ donc $s(C)=C$ ; $BC=1$ et $DC=1$. Égalité vérifiée.
\item $M=H(0,1,1)$ : la réflexion de plan $x-y=0$ échange $x$ et $y$ et conserve $z$, donc $s(H)=(1,0,1)=F$. $BH=\sqrt{1+1+1}=\sqrt{3}$ et $DF=\sqrt{1+1+1}=\sqrt{3}$. Égalité vérifiée.
\end{itemize}
\end{enumerate}
\textbf{Partie B}
\begin{enumerate}
\item $A\in\mathcal{P}$ et $C\in\mathcal{P}$ (vérifié en A.1), donc la droite $(AC)$ est incluse dans $\mathcal{P}$.
\item D'après la propriété admise, $s\circ s_{(AC)}$ est la \cle{réflexion} par rapport au plan $\mathcal{Q}$ contenant $(AC)$ et perpendiculaire à $\mathcal{P}$. Un vecteur normal à $\mathcal{Q}$ est orthogonal à la fois à $\vec{u}(1,1,0)$ (direction de $(AC)$) et à $\vec{n}_{\mathcal{P}}(1,-1,0)$ : on prend $\vec{n}_{\mathcal{Q}}=\vec{n}_{\mathcal{P}}\times\vec{u}=(0,0,2)$, colinéaire à $\vec{k}$. $\mathcal{Q}$ passe par $A(0,0,0)$ : son équation est $z=0$, c'est-à-dire le plan $(ABC)$. Donc $s\circ s_{(AC)}=s_{(ABC)}$.
\item $E(0,0,1)$ a pour image par la réflexion de plan $z=0$ le point $E'(0,0,-1)$. (Contrôle direct : $s_{(AC)}(x,y,z)=(y,x,-z)$ donne $s_{(AC)}(E)=(0,0,-1)$, point de $\mathcal{P}$ car $x-y=0$, donc invariant par $s$ : résultat identique.)
\end{enumerate}
\textbf{Barème indicatif (sur 5 pts) :} A.1 : 1 pt ; A.2 : 1 pt ; A.3 : 1 pt ; B.1 : 0,5 pt ; B.2 : 1 pt ; B.3 : 0,5 pt.
\textbf{Points de vigilance :} pour A.2, il faut justifier \emph{les deux} conditions du plan médiateur (milieu dans le plan \emph{et} perpendicularité) ; pour B.2, citer explicitement la propriété admise par l'énoncé ; le repère $(A;\overrightarrow{AB},\overrightarrow{AD},\overrightarrow{AE})$ est orthonormé car $ABCDEFGH$ est un cube : le dire.
\end{corrige}

\begin{corrige}[Exercice 10 -- Type Bac : La fourmi et le plus court chemin]
\begin{enumerate}
\item La réflexion par rapport au plan $z=0$ transforme $(x,y,z)$ en $(x,y,-z)$. Donc $B(5,4,3)$ a pour image $B'(5,4,-3)$.
\item Le plan $z=0$ est le plan médiateur de $[BB']$ (il contient le milieu $(5,4,0)$ et est perpendiculaire à $(BB')$). Tout point $M$ du plan médiateur d'un segment est équidistant de ses extrémités : $MB=MB'$. (C'est la conservation des distances par la réflexion, $M$ étant invariant.)
\item Pour tout point $M$ du sol : $AM+MB=AM+MB'\geq AB'$ d'après l'inégalité triangulaire appliquée aux points $A$, $M$, $B'$. L'égalité a lieu si et seulement si $M$ appartient au segment $[AB']$, c'est-à-dire lorsque $M$ est l'intersection de la droite $(AB')$ avec le plan $z=0$ (intersection qui existe car $A$ est au-dessus du sol, $z_A=2>0$, et $B'$ en dessous, $z_{B'}=-3<0$).
\item La droite $(AB')$ est dirigée par $\overrightarrow{AB'}(5,4,-5)$ et passe par $A(0,0,2)$ :
\[ \begin{cases} x=5t \\ y=4t \\ z=2-5t \end{cases}\quad(t\in\mathbb{R}). \]
Intersection avec $z=0$ : $2-5t=0\iff t=\dfrac{2}{5}$. Alors $x=5\times\dfrac{2}{5}=2$ et $y=4\times\dfrac{2}{5}=\dfrac{8}{5}=1,6$. Donc $M_0(2\,;\,1,6\,;\,0)$.
\item La distance minimale est
\[ AB'=\sqrt{(5-0)^2+(4-0)^2+(-3-2)^2}=\sqrt{25+16+25}=\sqrt{66}\ \text{m}\approx 8,12\ \text{m}. \]
\item Vitesse : $v=\SI{2}{cm/s}=\SI{0,02}{m/s}$. Temps minimal :
\[ t=\dfrac{\sqrt{66}}{0,02}=50\sqrt{66}\ \text{s}\approx 406,2\ \text{s}\approx 6\ \text{min}\ 46\ \text{s}. \]
\end{enumerate}
\textbf{Barème indicatif (sur 5 pts) :} Q1 : 0,5 pt ; Q2 : 1 pt ; Q3 : 1 pt ; Q4 : 1 pt ; Q5 : 1 pt ; Q6 : 0,5 pt.
\textbf{Points de vigilance :} c'est le symétrique de $B$ (point d'arrivée) qu'il faut utiliser car $M$ doit rester sur le plan ; l'inégalité triangulaire doit être citée avec son cas d'égalité (alignement) ; penser à convertir les unités dans la question bonus (cm/s en m/s).
\end{corrige}

\begin{corrige}[Exercice 11 -- Type Bac : Sécurisation d'un chantier]
\textbf{Partie A}
\begin{enumerate}
\item Réflexion de plan $\mathcal{P}_1:x=0$ : on change $x$ en $-x$ : $s_1(x,y,z)=(-x,\,y,\,z)$.
Réflexion de plan $\mathcal{P}_2:x=3$ : le milieu de $[MM']$ a pour abscisse $3$, donc $x'=2\times 3-x=6-x$ : $s_2(x,y,z)=(6-x,\,y,\,z)$.
\item $s_2\circ s_1(x,y,z)=s_2(-x,y,z)=\bigl(6-(-x),\,y,\,z\bigr)=(x+6,\,y,\,z)$.
C'est la \cle{translation} de vecteur $\vec{v}(6,0,0)$. C'est conforme au théorème : les plans $x=0$ et $x=3$ sont parallèles, distants de $3$, et $\vec{v}=2\vec{d}$ avec $\vec{d}(3,0,0)$ orienté de $\mathcal{P}_1$ vers $\mathcal{P}_2$.
\item $s_2\circ s_1(L)=s_2\circ s_1(1,2,5)=(1+6,\,2,\,5)=(7\,;\,2\,;\,5)$.
\item $s_1\circ s_2(x,y,z)=s_1(6-x,y,z)=\bigl(-(6-x),\,y,\,z\bigr)=(x-6,\,y,\,z)$ : translation de vecteur $(-6,0,0)$, opposé du précédent. Donc $s_1\circ s_2\neq s_2\circ s_1$ : la composition de deux réflexions de plans parallèles \textbf{n'est pas commutative} (les deux résultats sont des translations de vecteurs opposés, donc réciproques l'une de l'autre).
\end{enumerate}
\textbf{Partie B}
\begin{enumerate}
\item $\mathcal{P}_1:x=0$ a pour vecteur normal $\vec{n}_1(1,0,0)$ ; $\mathcal{P}_3:z=0$ a pour vecteur normal $\vec{n}_3(0,0,1)$. Or $\vec{n}_1\cdot\vec{n}_3=0$, donc les plans sont \cle{perpendiculaires}. Leur intersection est la droite $(\Delta):\begin{cases} x=0 \\ z=0 \end{cases}$, c'est-à-dire l'axe $(Oy)$.
\item D'après la propriété admise, $s_3\circ s_1$ est le \cle{demi-tour} d'axe $(\Delta)=(Oy)$.
\item $s_3\circ s_1(x,y,z)=s_3(-x,y,z)=(-x,\,y,\,-z)$. On reconnaît bien le demi-tour d'axe $(Oy)$ : $y$ est invariant, $x$ et $z$ sont changés en leurs opposés. Vérification sur $L(1,2,5)$ : $s_1(L)=(-1,2,5)$ puis $s_3(-1,2,5)=(-1,2,-5)$, et la formule donne directement $s_3\circ s_1(1,2,5)=(-1,2,-5)$ : les deux calculs concordent.
\end{enumerate}
\textbf{Barème indicatif (sur 6 pts) :} A.1 : 1 pt ; A.2 : 1 pt ; A.3 : 0,5 pt ; A.4 : 1 pt ; B.1 : 1 pt ; B.2 : 0,5 pt ; B.3 : 1 pt.
\textbf{Points de vigilance :} pour la perpendicularité de deux plans, justifier par le produit scalaire des vecteurs normaux ; dans A.4, la réponse « non commutative » doit être appuyée par le calcul explicite des deux composées ; ne pas oublier de citer la propriété admise dans B.2.
\end{corrige}

\begin{corrige}[Exercice 12 -- Entraînement Bac : Isométries et pavage]
\begin{enumerate}
\item $\mathcal{P}:x+y+z=0$, vecteur normal $\vec{n}(1,1,1)$, $a^2+b^2+c^2=3$. Avec $k=\dfrac{x+y+z}{3}$ :
\[ s(x,y,z)=(x,y,z)-2k(1,1,1)=\left(\frac{x-2y-2z}{3}\,;\,\frac{-2x+y-2z}{3}\,;\,\frac{-2x-2y+z}{3}\right). \]
Contrôle : $s(0,0,0)=(0,0,0)$ (point de $\mathcal{P}$ fixe) et $s(1,1,1)=(-1,-1,-1)$ (point de la normale renvoyé en son opposé) : cohérent.
\item Soit $M(x,y,z)\in\mathcal{P}$, donc $x+y+z=0$. Son image $t(M)=(x+1,y+1,z+1)$ vérifie $(x+1)+(y+1)+(z+1)=3\neq 0$ : $t(M)\notin\mathcal{P}$. Donc $t$ \textbf{ne laisse pas $\mathcal{P}$ invariant} : elle envoie $\mathcal{P}$ sur le plan parallèle $x+y+z=3$. (En revanche, $t$ conserve globalement la famille des plans $x+y+z=k$, $k\in\mathbb{Z}$, donc elle respecte le pavage.)
\item $f=t\circ s$ est la composée de deux isométries : c'est une \cle{isométrie}. Sa partie linéaire est celle de $s$ (la translation n'en change pas), de déterminant $-1$ : c'est une isométrie indirecte. Cherchons ses points fixes :
\[ f(x,y,z)=\left(\frac{x-2y-2z}{3}+1\,;\,\frac{-2x+y-2z}{3}+1\,;\,\frac{-2x-2y+z}{3}+1\right). \]
$f(M)=M$ équivaut, pour la première coordonnée, à $\dfrac{x-2y-2z}{3}+1=x\iff x-2y-2z+3=3x\iff 2x+2y+2z=3\iff x+y+z=\dfrac{3}{2}$ ; les deux autres équations donnent la même condition. L'ensemble des points fixes est donc le \textbf{plan $\mathcal{P}':x+y+z=\dfrac{3}{2}$}. Une isométrie indirecte admettant un plan de points fixes est une \cle{réflexion} : $f$ est la réflexion de plan $\mathcal{P}'$.
Interprétation : composer une réflexion avec une translation de vecteur \emph{perpendiculaire} au miroir décale le miroir de la moitié du vecteur : $\mathcal{P}$ est translaté de $\dfrac{1}{2}(1,1,1)$, ce qui donne bien $x+y+z=\dfrac{3}{2}$.
\item Le cube $[0,1]^3$ a pour centre $G\left(\dfrac{1}{2},\dfrac{1}{2},\dfrac{1}{2}\right)$, qui appartient à $\mathcal{P}'$ (somme des coordonnées : $\dfrac{3}{2}$). Le plan $\mathcal{P}'$ coupe le cube selon un hexagone régulier passant par les milieux de six arêtes ; $f$ échange les deux moitiés du cube situées de part et d'autre de cette section. On note que $f$ laisse ce cube globalement invariant : par exemple $f(0,0,0)=(1,1,1)$ et $f(1,1,1)=(0,0,0)$, les sommets diagonalement opposés sont échangés. En revanche, $f$ ne conserve pas le pavage de l'espace par les cubes unitaires translatés de $\mathbb{Z}^3$ (l'image du sommet $(1,0,0)$ n'a pas des coordonnées entières).
\end{enumerate}
\textbf{Points de vigilance :} la translation ne modifie pas la partie linéaire, donc pas le déterminant ; la recherche des points fixes est la clé de la reconnaissance ; attention au piège de la question 2 : « laisser invariant » signifie $t(\mathcal{P})\subset\mathcal{P}$, ce qui est faux ici.
\end{corrige}

\newpage

\coursec{Fiche bilan}

\begin{popfigure}
\begin{tikzpicture}[
  mindmap,
  grow cyclic,
  every node/.style={concept, text=white, font=\small\bfseries},
  root concept/.append style={concept color=popdark, font=\large\bfseries},
  level 1/.append style={level distance=4.2cm, sibling angle=60},
  level 2/.append style={level distance=2.6cm, sibling angle=45, font=\footnotesize},
]
\node[root concept]{Isométries de l'espace}
  child[concept color=popblue]{ node {Réflexion}
    child { node {Plan médiateur de $[MM']$} }
    child { node {Conserve distances, orthogonalité} }
  }
  child[concept color=poporange]{ node {Expression analytique}
    child { node {$x'=\pm x$ selon le plan} }
    child { node {Cas général : milieu + vecteur normal} }
  }
  child[concept color=popgreen]{ node {Composée de 2 réflexions}
    child { node {Plans parallèles : translation $2\vec{u}$} }
    child { node {Plans perpendiculaires : demi-tour} }
  }
  child[concept color=popink]{ node {Demi-tour}
    child { node {Axe $(D)$ : $I$ milieu, $(MM')\perp(D)$} }
    child { node {Expression analytique} }
  }
  child[concept color=poppurple]{ node {Composée de 2 demi-tours}
    child { node {Axes parallèles : translation} }
    child { node {Axes orthogonaux sécants : demi-tour} }
  }
  child[concept color=popgold]{ node {Reconnaissance}
    child { node {Forme de l'expression} }
    child { node {Points fixes} }
  };
\end{tikzpicture}
\legende{Carte mentale de la leçon : les six compétences clés autour des isométries de l'espace.}
\end{popfigure}

\begin{aretenir}[L'essentiel à retenir]
\textbf{1. Définitions clés}
\begin{itemize}
  \item \cle{Réflexion} de plan $(P)$ : $M' = s_{(P)}(M)$ si $(P)$ est le \textbf{plan médiateur} du segment $[MM']$, c'est-à-dire le plan passant par le milieu de $[MM']$ et perpendiculaire à la droite $(MM')$ ; et $M'=M$ si $M\in(P)$.
  \item \cle{Demi-tour} d'axe $(D)$ : $M' = s_{(D)}(M)$ si $(D)$ est la \textbf{médiatrice} de $[MM']$ dans le plan $(MM'D)$, c'est-à-dire : le milieu $I$ de $[MM']$ appartient à $(D)$ et $(MM')\perp(D)$ ; et $M'=M$ si $M\in(D)$.
  \item Réflexions et demi-tours sont des \cle{isométries} : elles conservent les distances, les angles, l'alignement, l'orthogonalité, les milieux, les aires et les volumes.
  \item Réflexions et demi-tours sont des \cle{involutions} : $s\circ s = \mathrm{Id}$, donc chacune est sa propre bijection réciproque.
\end{itemize}

\textbf{2. Formules à connaître par c\oe ur}

\begin{center}
\begin{tabularx}{\linewidth}{|l|X|}
\hline
\rowcolor{popblueL} \textbf{Isométrie} & \textbf{Expression analytique (repère orthonormé adapté)} \\
\hline
Réflexion de plan $z=0$ & $x'=x$,\quad $y'=y$,\quad $z'=-z$ \\
\hline
Réflexion de plan $x=0$ & $x'=-x$,\quad $y'=y$,\quad $z'=z$ \\
\hline
Réflexion de plan $y=0$ & $x'=x$,\quad $y'=-y$,\quad $z'=z$ \\
\hline
Demi-tour d'axe $(Oz)$ & $x'=-x$,\quad $y'=-y$,\quad $z'=z$ \\
\hline
Demi-tour d'axe $(Ox)$ & $x'=x$,\quad $y'=-y$,\quad $z'=-z$ \\
\hline
Demi-tour d'axe $(Oy)$ & $x'=-x$,\quad $y'=y$,\quad $z'=-z$ \\
\hline
Symétrie centrale de centre $O$ & $x'=-x$,\quad $y'=-y$,\quad $z'=-z$ \\
\hline
\end{tabularx}
\end{center}

Cas général : pour une réflexion de plan $(P)$ d'équation $ax+by+cz+d=0$, on écrit les deux conditions suivantes : le milieu $I$ de $[MM']$ appartient à $(P)$, et le vecteur $\overrightarrow{MM'}$ est colinéaire au vecteur normal $\vec{n}(a\,;b\,;c)$. On obtient un système qui permet d'exprimer $x',y',z'$ en fonction de $x,y,z$.

\textbf{3. Composées (théorèmes du cours)}
\begin{itemize}
  \item $s_{(P_2)}\circ s_{(P_1)}$ avec $(P_1)\parallel(P_2)$ : \textbf{translation} de vecteur $2\vec{u}$, où $\vec{u}$ est le vecteur orthogonal aux deux plans qui transforme $(P_1)$ en $(P_2)$.
  \item $s_{(P_2)}\circ s_{(P_1)}$ avec $(P_1)\perp(P_2)$ : \textbf{demi-tour} d'axe $(D)=(P_1)\cap(P_2)$.
  \item $s_{(D_2)}\circ s_{(D_1)}$ avec $(D_1)\parallel(D_2)$ : \textbf{translation} de vecteur $2\overrightarrow{H_1H_2}$, où $H_1$ et $H_2$ sont les intersections de $(D_1)$ et $(D_2)$ avec un plan perpendiculaire aux deux axes.
  \item $s_{(D_2)}\circ s_{(D_1)}$ avec $(D_1)$ et $(D_2)$ \textbf{orthogonales et sécantes} en un point $O$ : \textbf{demi-tour} d'axe la droite passant par $O$ et perpendiculaire au plan $(D_1,D_2)$.
\end{itemize}

\textbf{4. Méthodes express}
\begin{itemize}
  \item \emph{Reconnaître une isométrie} : examiner la forme analytique (un seul signe $-$ : réflexion ; deux signes $-$ : demi-tour ; trois signes $-$ : symétrie centrale), puis confirmer en cherchant les points fixes.
  \item \emph{Trouver les éléments caractéristiques} : l'ensemble des points fixes donne le plan (réflexion) ou l'axe (demi-tour) ; pour une translation, calculer $\overrightarrow{MM'}$ pour un point $M$ quelconque.
  \item \emph{Attention} : la composée n'est pas commutative en général : $s_{(P_2)}\circ s_{(P_1)}\neq s_{(P_1)}\circ s_{(P_2)}$ (on obtient la translation de vecteur opposé lorsque les plans sont parallèles).
\end{itemize}
\end{aretenir}

\begin{attention}[Pièges fréquents au Bac]
\begin{itemize}
  \item Ne pas confondre \textbf{plan médiateur} (réflexion, dans l'espace) et \textbf{médiatrice} (demi-tour, dans un plan) : le plan médiateur de $[MM']$ est un plan, la médiatrice est une droite.
  \item Pour la composée de deux demi-tours donnant un demi-tour, les axes doivent être orthogonaux \textbf{et sécants} : deux droites orthogonales non coplanaires ne donnent pas ce résultat.
  \item Dans $f\circ g$, c'est $g$ qui agit \textbf{en premier}. Inverser l'ordre change le vecteur de la translation obtenue.
  \item Une expression analytique du type $x'=-x$, $y'=-y$, $z'=z$ n'est \textbf{pas} une réflexion : c'est un demi-tour (deux signes $-$).
  \item Les formules du tableau ne sont valables que dans un repère \textbf{orthonormé} adapté ; pour un plan quelconque, il faut revenir à la méthode « milieu + vecteur normal ».
\end{itemize}
\end{attention}

\begin{aretenir}[Checklist avant le Bac]
\begin{itemize}
  \item[$\square$] Je sais définir une réflexion à l'aide du plan médiateur d'un segment.
  \item[$\square$] Je sais définir un demi-tour à l'aide de la médiatrice et de l'orthogonalité.
  \item[$\square$] Je connais par c\oe ur les expressions analytiques des réflexions de plans de coordonnées.
  \item[$\square$] Je connais par c\oe ur les expressions analytiques des demi-tours d'axes de coordonnées.
  \item[$\square$] Je sais déterminer l'expression analytique d'une réflexion de plan quelconque (milieu + vecteur normal).
  \item[$\square$] Je sais que la composée de deux réflexions de plans parallèles est une translation de vecteur $2\vec{u}$.
  \item[$\square$] Je sais que la composée de deux réflexions de plans perpendiculaires est un demi-tour d'axe l'intersection des plans.
  \item[$\square$] Je sais que la composée de deux demi-tours d'axes parallèles est une translation.
  \item[$\square$] Je sais que la composée de deux demi-tours d'axes orthogonaux sécants est un demi-tour.
  \item[$\square$] Je sais reconnaître la nature d'une isométrie à partir de son expression analytique et de ses points fixes.
  \item[$\square$] Je pense à vérifier l'ordre de composition : $f\circ g$ signifie « $g$ d'abord, puis $f$ ».
  \item[$\square$] Je sais utiliser une isométrie pour transformer un problème (alignement, distances, angles droits) en un problème plus simple.
\end{itemize}
\end{aretenir}

\findecours

\end{document}
